% concatenated from name20250929_2_.tex
\documentclass[10pt]{article}
\usepackage{geometry, enumerate, amsmath, amssymb, amsthm, amscd, hyperref, graphicx, color, enumerate, mathrsfs}
\usepackage[all]{xy}

\title{Geometric configuration of integrally closed Noetherian domains}
%\\ {\small (previously: Polynomial Krull Domains)}}

\date{\today}


\author{Gyu Whan Chang\footnote{Department of Mathematics Education, Incheon National University, Incheon 22012, Republic of Korea E-mail: whan@inu.ac.kr} \and Giulio Peruginelli\footnote{Department of Mathematics ``Tullio Levi-Civita'' University of Padova, Via Trieste, 63 35121 Padova, Italy. E-mail: gperugin@math.unipd.it, ORCID: https://orcid.org/0000-0001-7694-8920}}




\numberwithin{equation}{section}

\newtheorem{Prop}[equation]{Proposition}
\newtheorem{Thm}[equation]{Theorem}
\newtheorem{Lem}[equation]{Lemma}
\newtheorem{Cor}[equation]{Corollary}
\newtheorem{Conj}[equation]{Conjecture}
\newcommand{\Proof}{\noindent{\bfseries Proof : }}
\theoremstyle{definition}
\newtheorem{Def}[equation]{Definition}
\newtheorem{Ex}[equation]{Example}
\newtheorem{Rem}[equation]{Remark}
\newtheorem{Rems}[equation]{Remarks}
\theoremstyle{definition}
\newtheorem{Ques}[equation]{Question}

\newcommand{\F}{\mathbb{F}}
\newcommand{\Q}{\mathbb{Q}}
\newcommand{\N}{\mathbb{N}}
\newcommand{\Z}{\mathbb{Z}}
\newcommand{\R}{\mathbb{R}}
\newcommand{\C}{\mathbb{C}}
\renewcommand{\O}{\mathbb{O}}
\newcommand{\PP}{\mathbb{P}}
\newcommand{\Int}{\textnormal{Int}}
\newcommand{\IntQ}{\Int_{\Q}}
\newcommand{\olK}{\overline{K}}
\newcommand{\uE}{\underline{E}}
\newcommand{\oQp}{\overline{\Q_p}}
\newcommand{\Br}{\text{Br}}
\newcommand{\Gal}{\text{Gal}}
\newcommand{\Zar}{\text{Zar}}
\newcommand{\K}{\widehat{K}}
\newcommand{\V}{\widehat{V}}
\newcommand{\W}{\widehat{W}}
\newcommand{\M}{\widehat{M}}
\newcommand{\rat}{\text{rat}}
\newcommand{\raf}{\text{raf}}
\newcommand{\ra}{\text{ra}}
\newcommand{\rt}{\text{rt}}
\renewcommand{\t}{\text{t}}
\newcommand{\rk}{\text{rk}}
\newcommand{\abK}{\overline{\K}}
\newcommand{\abV}{\overline{\V}}
\newcommand{\Spec}{\text{Spec}}
\newcommand{\wmW}{\widetilde{\mathcal{W}}}
\newcommand{\wW}{\widetilde{W}}
\newcommand{\wM}{\widetilde{M}}
\newcommand{\ov}{\overline{v}}


\newcommand{\oZp}{\overline{\Z_p}}

\newcommand{\hZ}{\widehat{\mathbb{Z}}}
\newcommand{\ohZ}{\overline{\hZ}}
\newcommand{\Cl}{\text{Cl}}

\newcommand{\br}{\text{br}}
\newcommand{\unr}{\text{unr}}

\newcommand\rc { \color{red}}
\newcommand\bc { \color{blue}}
\newcommand\ec { \color{black}}
%\usepackage[notcite, notref]{showkeys}

\renewcommand{\baselinestretch}{1.05}

\begin{document}



\maketitle

%that is, subrings of $\Q[X]$ containing $\Z[X]$ that are also  Krull domains

% into residually algebraic or residually transcendental extensions of $\mathbb{Z}_{(p)}$

\begin{abstract}
In this paper, we completely describe the family of integrally closed Noetherian domains between $\Z[X]$ and $\Q[X]$. We accomplish this result by classifying the Krull domains between these two polynomial rings. To this end, we first describe the
DVRs of $\mathbb{Q}(X)$ lying over $\mathbb{Z}_{(p)}$ for some prime $p \in \Z$, by  distinguishing them according to whether the extension of the residue fields is algebraic or transcendental. We unify the known descriptions of such valuations by considering ultrametric balls in $\C_p$, the completion of the algebraic closure of the field $\Q_p$ of $p$-adic numbers.
We then study when the intersection $R$ of such DVRs with $\Q[X]$ is of finite character, so that 
$R$ is a Krull domain, and we finally compute the divisor class group of $R$. It turns out that such a ring is formed by those polynomials which simultaneously  map a finite union of ultrametric balls of $\C_p$ to its valuation domain $\O_p$, as $p\in\Z$ ranges through the set of primes. 
By a result of Heinzer, the Krull domains of this class are precisely the integrally closed Noetherian domains between $\Z[X]$ and $\Q[X]$. This novel approach provides a geometric understanding of this class of integrally closed domains. Furthermore, we also describe the UFDs between $\Z[X]$ and $\Q[X]$.
\end{abstract}
\medskip
% which is the $p$-adic analogue of the field of complex numbers

{\small \noindent Keywords: Noetherian domain, integrally closed domain, DVR, Krull domain, UFD.\\
\medskip
\noindent MSC Primary  13F05, 13B22, 13E05, 13B25, 13F20\\}
%13B22: Integral closure of commutative rings and ideals [See also 13A35]; integrally closed rings, related rings (Japanese, etc.)
%13E05(1973–now)Commutative Noetherian rings and modules
%13F20(1973–now)Polynomial rings and ideals; rings of integer-valued polynomials 
%13B25(1973–now)Polynomials over commutative rings
%13F05(1973–now)Dedekind, Prüfer, Krull and Mori rings and their generalizations
%13F20, 13F05, 12J20,  13A18.\\}

%\newpage
\tableofcontents
%\newpage
\vskip1cm
\section*{Introduction}\label{sec:intro}
\addcontentsline{toc}{section}{\nameref{sec:intro}}


Let $\bf{X}$ be a topological space, $\C$ be the field of complex numbers and $\bf{C}(\bf{X}, \C)$ be the ring of continuous complex-valued functions on $\bf{X}$. Then $\bf{C}(\bf{X}, \C)$ is a commutative ring with identity. It is known that $\bf{C}(\bf{X}, \C)$ is an integral domain if and only if $\bf{X}$ is irreducible, i.e., $\bf{X}$ cannot be expressed as a finite union of proper closed subsets. It is also known that $\bf{C}(\bf{X}, \C)$ is Noetherian if and only if $\bf{X}$ is a finite set, and in this case, if $\bf{X}$ is Hausdorff, then $\bf{C}(\bf{X}, \C) \cong \C^{|\bf{X}|}$ which is a zero-dimensional Noetherian ring. 
 The class of Noetherian domains is of the utmost importance in all areas of mathematics and especially in commutative algebra
 and algebraic geometry.

 
Let $\overline{\Q}$ be the algebraic closure of $\Q$ and $X$ be an indeterminate over $\C$, then we have the following diagram of subrings of $\bf{C}(\C, \C)$:

$$
\xymatrix{
&&\overline{\Q}[X]\ar[dr]&&\\
\Z[X]\ar[r]&\Q[X]\ar[ur]\ar[dr]&&\C[X]\ar[r]&\bf{C}(\C, \C)\\
&&\R[X]\ar[ur]&&}
$$

%$$\Z[X] \subseteq \Q[X] \subseteq \overline{\Q}[X] \subseteq \C[X] %\subseteq \bf{C}(\C, \C)$$ 
%$$\Q[X] \subseteq {\R}[X] \subseteq \C[X].$$
Although $\bf{C}(\C, \C)$ is neither an integral domain nor a Noetherian ring, $\bf{C}(\C, \C)$ contains the Noetherian integral domain ${\C}[X]$. Hence, we can ask whether we can classify the Noetherian subrings of $\bf{C}(\C, \C)$. However, it seems to be a very difficult task, so in this paper, as the first step in a long journey towards that end, we will classify all integrally closed Noetherian subrings of  $\Q[X]$ containing $\Z[X]$. For this purpose, we describe the Krull domains between these polynomial rings. We now recall the definition of this class of integral domains.


\medskip
An integral domain is said to be a discrete valuation ring (DVR) if it is a principal ideal domain with exactly one nonzero prime ideal; equivalently, it is a Noetherian valuation domain that is not a field. Let $D$ be an integral domain with quotient field $K$ and $X^1(D)$ be the set of height one prime ideals of $D$.
Then $D$ is a Krull domain if the following three conditions are satisfied:
\begin{enumerate}
\item[{(a)}] $D = \bigcap\limits_{P \in X^1(D)}D_P$,
\item[{(b)}] $D_P$ is a DVR for all $P \in X^1(D)$, and
\item[{(c)}] each nonzero nonunit of $D$ is a unit of $D_P$ for all but a finite number of $P \in X^1(D)$.
\end{enumerate}
The class of Krull domains includes Dedekind domains, unique factorization domains (UFDs),
the integral closure of a Noetherian domain, and polynomial/power series rings over a Krull domain (see, for example, \cite{Fossum} and \cite{Gilm}). In particular, if $D$ is a field, then $X^1(D) = \emptyset$, so it satisfies the three conditions of a Krull domain. Thus, a field is a Krull domain.

Let $\Z$ be the ring of integers, $\Q$ be the field of rational numbers,
and $X$ be an indeterminate over $\Q$. Following \cite{Chang}, we say that
a subring of $\Q[X]$ containing $\Z[X]$ is a polynomial overring of $\Z$.
 In \cite{PerDedekind}, Peruginelli called a polynomial overring of $\mathbb{Z}$
a polynomial Dedekind domain if it is a Dedekind domain. He then characterized
a polynomial Dedekind domain with the property that its residue fields of prime characteristic are
finite fields, and its Picard group. Later, in \cite{PerStacked}, Peruginelli
completely classified the polynomial Dedekind domains $R$ by removing the condition on the cardinality of the residue fields;
in particular, he showed that $R$ is equal to a ring of integer-valued polynomials over an algebra.
In this paper, we say that a polynomial overring of $\Z$ is a polynomial Krull domain if
it is also a Krull domain and we are going to completely characterize the polynomial Krull domains which allows us to give a complete characterization of integrally closed Noetherian domains that are subring $\Q[X]$ containing $\Z[X]$.


It is known that $D$ is a Krull domain if and only if there exists a family $F$ of distinct DVRs with quotient field $K$ such that 
(i) $D = \bigcap\limits_{V \in F}V$ and (ii) every nonzero element of $D$ is a unit in all 
but a finite number of rings in $F$ \cite[Proposition 3]{n55}, and in this case, $F$ is called a defining family for the Krull domain $D$.

 Let $E$ be the ring of entire functions. Then $E = \bigcap\limits_{Q \in X^1(E)}E_Q$,
$E_Q$ is a DVR for all $Q \in X^1(E)$, but the intersection $E = \bigcap\limits_{Q \in X^1(E)}E_Q$
is not of finite character, i.e., $E$ does not satisfy the condition (c) of Krull domains,
so $E$ is not a Krull domain. Hence, in order to characterize the polynomial Krull domains, we
first classify the DVRs of $\Q(X)$, then we have to prove when a family of such DVRs has the finite character property (i.e. property (c) above).

Let $W$ be a DVR of $\mathbb{Q}(X)$ with maximal ideal $M$. Then either $W \cap \mathbb{Q} = (0)$ or $W \cap \mathbb{Q} = \mathbb{Z}_{(p)}$
for some prime number $p \in \mathbb{Z}$. It is well known that $W \cap \mathbb{Q} = (0)$ if and only if
$W = \mathbb{Q}[X]_{f\mathbb{Q}[X]}$ for some irreducible $f \in \mathbb{Q}[X]$ or
$W = \mathbb{Q}[\frac{1}{X}]_{(\frac{1}{X})}$ (see for example \cite{Chevalley}). In this paper, we are
mostly interested in the case $W \cap \mathbb{Q} = \mathbb{Z}_{(p)}$.
In \cite{PerStacked}, the second-named author of this paper
completely classified the DVRs $(W,M)$ such that $W/M$ is algebraic over $\mathbb{Z}/p\mathbb{Z}$.
In this paper, we will recall the classification of the DVRs $(W,M)$ such that 
$W$ is a residually transcendental extension of $\mathbb{Z}_{(p)}$, i.e.,
$W/M$ is transcendental over $\mathbb{Z}/p\mathbb{Z}$, along the line of the work of Alexandru and Popescu (\cite{AlePop}).  
This paper consists of five sections including the introduction. 
In  Sections 1. and 1.2, we first review the notion of 
 the ideal factorization property of Krull domains,
the divisor class group of Krull domains and rings of integer-valued polynomials, which allow us to have a nice representation for these class of polynomial Krull domains.  In Section \ref{res transcendental},  we recall the notion of monomial valuation, which is a standard way to extend a valuation domain $V$ with quotient field $K$ to the field of rational functions $K(X)$ and we use it in order to describe the set $\mathcal W_{\rt}$ of residually transcendental extensions of $V$ to $K(X)$. In Section \ref{section parametrization} we show how it is possible to parametrize the elements of $\mathcal W_{\rt}$  by means of ultrametric balls in $\olK$ with respect to a fixed extension $U$ of $V$. It is well-known that a residually transcendental extension $W$ of $V$ to $K(X)$ is a torsion extension of $V$, i.e., the value group of $W$ is torsion over that of $V$. We  study in Section \ref{torsion extensions} the  ordering of torsion extensions $W_1, W_2$ of $V$ to $K(X)$, that is we say that  $W_1\leq W_2$ if $W_1\cap K[X]\subseteq W_2\cap K[X]$. If this happens, in particular we have $W_1 \cap W_2 \cap K[X] = W_1 \cap K[X]$. In Section \ref{res transc completion} we study residually transcendental extension of the completion $\V$ of $V$ and show that $\mathcal W_{\rt}$ is in bijection with the residually transcendental extensions of $\V$ to $\K(X)$.

%the residually transcendental extensions 
In Section 3, we use the results of Section 2 to parametrize the residually transcendental extensions of $\mathbb{Z}_{(p)}$ to $\Q(X)$ for some prime number $p \in \Z$ in terms of ultrametric balls in the algebraic closure $\oQp$ of the $p$-adic field $\Q_p$. We then unify the descriptions of DVRs of $\Q(X)$ which are either residually transcendental or residually algebraic  over $\Z_{(p)}$ by means of ultrametric balls with possibly infinite radius in $\C_p$, the completion of $\oQp$. We show that there is a partial order on this set of DVRs which mirrors the containment between this ultrametric balls, modulo the action of the Galois group $\Gal(\oQp/\Q_p)$ (Theorem \ref{ultrametric balls modulo Galois are DVRs}). Finally, in Section 4, we study first the class of Krull domains between $\Z_{(p)}[X]$ and $\Q[X]$ for a fixed prime $p\in\Z$ (\S \ref{section local case}), and as a special case we characterize the UFDs between these two polynomial rings (\S \ref{section UFD local}). We completely describe the prime spectrum of such a   ring, characterizing when a prime ideal is maximal in geometrical terms. We then globalize these results in Section \ref{section global case}. Namely, we consider a nonempty set $\Lambda$  of integer primes, and, for each $p\in\Lambda$, a finite set of DVRs $W_{p, j}$ of $\mathbb{Q}(X)$ with $X \in W_{p,j}$ and $W_{p,j} \cap \mathbb{Q} = \mathbb{Z}_{(p)}$,
and $R= \bigcap\limits_{p \in \Lambda}(\bigcap\limits_{j=1}^{m_p}(W_{p,j} \cap \mathbb{Q}[X]))$. We study which one of $W_{p,j}$'s is superflous 
and when $R$ is of finite character, so that $R$ is a Krull domain (Theorem \ref{Krull domains finite character}). 
We also compute the divisor class group Cl$(R)$ of the Krull domain $R$, showing that Cl$(R)$ is a direct sum of a countable family of finitely generated abelian groups and, conversely, we show that for any such a group $G$ there exists a Krull non-Dedekind domain between $\Z[X]$ and $\Q[X]$ with divisor class group isomorphic to $G$ (\S \ref{section: construction Krull}). Furthermore, we also characterize UFDs contained between $\Z[X]$ and $\Q[X]$ (Corollary \ref{global UFD}).

%the $v$- and $t$-operation, the ($t$-)ideal class group, and

%\section{The $t$-operation and Int$(D)$ for a domain $D$}

%In this section, we review some definitions, say, the $v$- and $t$-operation,the ($t$-)ideal class group, and rings of integer-valued polynomials for easy understanding of the results in this paper.

%\subsection{The $v$- and $t$-operation}
%A Dedekind domain is an integral domain in which each ideal can be written as a finite product of prime ideals, equivalently,
%each nonzero ideal is invertible. Krull domains have a similar characterization for which
%we need the $v$- and $t$-operation and these operations are very useful when we study the factorization properties of Krull domains.

%Let $F(D)$ be the set of nonzero fractional ideals of $D$, i.e., $F(D)$ is the set of nonzero $D$-submodules $I$ of $K$ such tha $dI \subseteq D$ for some $0 \neq d \in D$. It is clear that for $I \in F(D)$, if we set $I^{-1} = \{x \in K \mid xI \subseteq D\}$, then $I^{-1} \in F(D)$, whence $I_v := (I^{-1})^{-1}$ is well defined.
%The $t$-operation on $D$ is defined by $I_t = \bigcup\{J_v \mid J \in F(D), J \subseteq I$, and $J$ is finitely generated$\}$.

%Let $*=v$ or $t$. We say that $I \in F(D)$ is a $*$-ideal if $I_*=I$, a $*$-ideal is a prime $*$-ideal if it is a prime ideal, and a $*$-ideal is a maximal $*$-ideal if it is maximal among proper integral $*$-ideals of $D$.
%Let $*$-Max$(D)$ be the set of maximal $*$-ideals of $D$. Then $t$-Max$(D) \neq \emptyset$ if and only if $D$ is not a field, while $v$-Max$(D) = \emptyset$ even though $D$ is not a field as in the case of a rank one nondiscrete valuation domain.  An integral domain $D$ is said to be a Mori domain if it satisfies the ascending chain condition for integral $v$-ideals of $D$, and, in this case, $v=t$, i.e., $I_v = I_t$ for all $I \in F(D)$.

%An $I$ is said to be $*$-invertible if $(II^{-1})_* = D$.
%Then a $t$-invertible fractional ideal is $v$-invertible but a $v$-invertible
%fractional ideal need not be $t$-invertible. It is known that $D$ is completely integrally closed if and only if every nonzero ideal of $D$ is $v$-invertible \cite[Theorem 34.6]{Gilm}.  It is also known that $D$ is a Krull domain if and only if $D$ is a completely integrally closed Mori domain \cite[Theorem 3.6 and Corollary 3.13]{Fossum}, if and only if every nonzero (prime) ideal of $D$ is $t$-invertible, if and only if every nonzero prime ideal of $D$ contains a $t$-invertible prime ideal \cite[Theorem 3.6]{k89}.
%Finally, if $D$ is a Krull domain, then $X^1(D) = t$-Max$(D)$; so each maximal $t$-ideal of $D$ has height one.

%\subsection{The $(t$-)ideal class group of an integral domain}

%Let Tinv$(D)$ be the set of $t$-invertible fractional $t$-ideals of $D$, Inv$(D)$ be the set of invertible fractional ideals of $D$, and Prin$(D)$ be the set of nonzero principal fractional ideals of $D$. Then Tinv$(D)$ is an abelian group under the multiplication $I*J = (IJ)_t$ for all $I, J \in F(D)$.
%Note that an invertible fractional ideal is a $t$-invertible $t$-ideal, So Inv$(D)$ is a subgroup of Tinv$(D)$. Moreover, Prin$(D)$ is a subgroup of both Tinv$(D)$ and Inv$(D)$, so Prin$(D) \leq$ Inv$(D) \leq$ Tinv$(D)$.
%The $(t$-)ideal class group and the Picard group of $D$ are defined by the factor groups
%$$\text{Cl}(D) = \text{Tinv}(D)/\text{Prin}(D) \text{  and  } \text{Pic}(D) = \text{Inv}(D)/\text{Prin}(D),$$ respectively. It is clear that Pic$(D)$ is a subgroup of Cl$(D)$;
%Cl$(D)$ is the usual divisor class group of $D$ when $D$ is a Krull domain; 
%and Cl$(D) =$ Pic$(D)$ when 
%$D$ is a Dedekind domain. The $(t$-)ideal class group of $D$ reflects on the factorization properties of $D$. For example, $D$ is a UFD (resp., PID) if and only if $D$
%is a Krull domain (resp., Dedekind domain) with Cl$(D) = \{0\}$ \cite[Proposition 6.1]{Fossum} 
%(resp., \cite[ ]{ }).
\section{Preliminaries}
%on Krull domains, rings of integer-valued polynomials and residually transcendental extensions
%Title? 
\subsection{Algebraic properties of Krull domains}

Let $D$ be an integral domain with quotient field $K$ and $F(D)$ be the set of nonzero fractional ideals $I$ of $D$, i.e., $I$ is a $D$-submodule of $K$ such that $dI \subseteq D$ for some $0 \neq d \in D$.
For each $I, J \in F(D)$, let $(I:_KJ) = \{x \in K \mid xJ \subseteq I\}$,
then $(I:_KJ) \in F(D)$. In particular, $I^{-1} = (D:_KI)$ and $I_v = (I^{-1})^{-1}$, so $(I^{-1})_v = (I_v)^{-1} = I^{-1}$. An $I \in F(D)$ is said to be invertible if $II^{-1} = D$. It is knonw tht $I$ is invertible if and only if $I$ is finitely generated and $ID_M$ is principal for all maximal ideals $M$ of $D$, and that if $I$ is invertible, then $(II^{-1})_v = D$ and $I_v = I$.


\subsubsection{Ideal factorization property of Krull domains}

A Dedekind domain $D$ is an integral domain in which each nonzero ideal  can be written as a finite product of prime ideals. Then
$D$ is a Dedekind domain if and only if $D$ is an integrally closed one-dimensional Noetherian domain, if and only if each nonzero ideal of $D$ is invertible. It is known that $D$ is a Dedekind domain if and only if 
$D$ is a Krull domain of (Krull) dimension at most one. Furthermore, 
 the integral closure of a Noetherian domain is a Krull domain 
\cite[Theorem 2]{n55}, so 
an overring of $2$-dimensional Noetherian domain is a Krull domain if and only if it is an integrally closed Noetherian domain \cite[Theorem 9]{Heinzer}. Hence, a Krull domain can be considered as a higher dimensional generalization of Dedekind domains, so it is natural to ask whether a Krull domain can be characterized by an ideal factorization property as in Dedekind domains. 
 
For $I \in F(D)$, let $I_t = \bigcup\{J_v \mid J \in F(D), J \subseteq I$ and $J$ is finitely generated$\}$. Then $I_t \in F(D)$ and the map $t: F(D) \rightarrow F(D)$, given by $I \mapsto I_t$, is well defined. 
The $t$-operation also has the following four properties for all $I, J \in F(D)$ and $0 \neq a \in K$: (i) $D_t = D$, (ii) $(aI)_t = aI_t$, (iii) $I \subseteq I_t$ and if $I \subseteq J$, then $I_t \subseteq J_t$ and (iv) $(I_t)_t = I$. An $I \in F(D)$ 
is said to be $t$-invertible if $(II^{-1})_t = D$. 
 We then use the $t$-operation to characterize Krull domains as follows: $D$ is a Krull domain if and only if, given a nonzero ideal $I$ of $D$, 
$I_t = (P_1^{e_1} \cdots P_n^{e_n})_t$ for some $P_1, \dots , P_n \in X^1(D)$ and positive integers $e_1, \dots, e_n$, 
if and only if, given a nonzero element $a$ of $D$, 
$aD = (P_1^{e_1} \cdots P_n^{e_n})_t$ for some $P_1, \dots , P_n \in X^1(D)$ and positive integers $e_1, \dots, e_n$,
if and only if every nonzero (prime) ideal of $D$ is $t$-invertible,
 if and only if every nonzero prime ideal of $D$ contains a $t$-invertible prime ideal;
in this case, $I_v = I_t$ for all $I \in F(D)$ (see, for example, \cite{k89}).


\subsubsection{Divisor class group of a Krull domain}

Let $\mathcal{D}(D) = \{I_v \mid I \in F(D)\}$. Then 
$\mathcal{C}(D)$ is a commutative semigroup under $I*J = (IJ)_v$,
and $\mathcal{D}(D)$ is a group, i.e., $(II^{-1})_v = D$ for all $I \in \mathcal{D}(D)$, if and only if $D$ is completely integrally closed. Next let Inv$(D)$ be the group of invertible factional ideals of $D$ and Prin$(D) = \{xD \mid x \in K, x \neq 0\}$. 
Then Inv$(D)$ is a subgroup of $\mathcal{D}(D)$ and Prin$(D)$ is a subgroup of Inv$(D)$, so Prin$(D) \subseteq$ Inv$(D) \subseteq \mathcal{D}(D)$. 
The divisor class group of a completely integrally closed domain $D$ is defined by the factor group Cl$(D) = \mathcal{D}(D)/$Prin$(D)$ of $\mathcal{C}(D)$ modulo Prin$(D)$ and the ideal class group or Picard group of an integral domain $D$ is defined by the factor group Inv$(D)/$Prin$(D)$ of Inv$(D)$ modulo Prin$(D)$. Hence, if $D$ is completely integrally closed, then Pic$(D)$ is a subgroup of Cl$(D)$. 

It is known that $x \in K$ is almost integral over $D$ if and only if $xI \subseteq I$ for some nonzero ideal $I$ of $D$, so a valuation domain $V$ is completely integrally closed if and only if the (Krull) dimension of $V$ is at most one; in particular, a DVR is completely integrally closed. A Krull domain $D$ is definited by an intersection of DVRs, so $D$ is completely integrally closed and its divisor class group is also well defined.


A $\pi$-domain $D$ is a Krull domain with Pic$(D) =$ Cl$(D)$;
equivalently, every nonzero principal ideal of $D$ can be written as a finite product of prime ideals \cite[IV. $\pi$-domains]{k89}. Hence, a Dedekind domain $D$ is a $\pi$-domain and Pic$(D) =$ Cl$(D)$.
A UFD is a Krull domain, and Krull domain $D$ is a UFD if and only if Cl$(D) = \{0\}$. A Krull domain $D$ is called an almost factorial domain if Cl$(D)$ is torsion; equivalently, given $a, b \in D$,
there is a positive integer $n$, which depends on $a,b$,
so that $a^nD \cap b^nD$ is principal.
Therefore, the divisor class group of a Krull domain $D$ measures how far $D$ is from a UFD and reflects the ideal factorization property of a Krull domain. 


\subsection{Rings of integer-valued polynomials}


Let $D$ be an integral domain with quotient field $K$,
$X$ be an indeterminate over $D$, $D[X]$ be the polynomial ring over $D$,
and $A$ be a torsion-free $D$-algebra. By embedding $A$ and $K$ in the extended $K$-algebra $B=A\otimes_D K$, 
we may evaluate a polynomial $f\in K[X]$ at an element $a\in A$, so that $f(a)\in B$. If $f(a)\in A$, we say that $f(X)$ is integer-valued at $a$. 
Given a subset $S$ of $A$, we define the ring of integer-valued polynomials on $S$ as
$$\Int_K(S,A)=\{f\in K[X] \mid f(s)\in A,\forall s\in S\}$$
Clearly, $D[X]\subseteq\Int_K(S, A)\subseteq K[X]$. If $S=A$, we set $\Int_K(A)=\Int_K(A,A)$, and if $A=D$, then we omit the subscript $K$, i.e., $\Int(S,D)= \Int_K(S,D)$ and $\Int(D) = \Int_K(D,D)$. 

Given a valuation domain $V$ with quotient field $K$, the \emph{polynomial closure} of a subset $S$ of $K$ is defined as the largest subset $\overline{S}$ of $K$ for which
$$\Int(S,V)=\Int(\overline{S},V).$$
A subset $S$ of $K$ is \emph{polynomially closed} if $S=\overline{S}$. If $V$ has rank one, Chabert proves in \cite{ChabPolCloVal}  
that the polynomially closed subsets of $K$ form a topology of closed subsets of $K$. For valuation domains of higher rank, the result is no more true \cite{PSpolclo}.

It is well known that Int$(\Z)$ is a non-Noetherian two-dimensional Pr\"ufer domain \cite[Remarks VI.1.8]{cc97}
and Pic(Int$(\Z))$ is a free abelian group on a countably infinite basis \cite[Corollary 7]{ghls90}.
Let $\PP$ be the set of prime numbers of $\Z$. For $p \in \PP$, let $\Z_p$ be the ring of $p$-adic integers,
$T_p = \{\alpha \in \Z_p \mid \alpha$ is transcendental over $\Q\}$, $\Z_{(p), \alpha} = \{\varphi \in \Q(X) \mid \varphi(\alpha) \in \Z_p\}$,
and $M_{(p), \alpha} = \{\varphi \in \Q(X) \mid \varphi(\alpha) \in p\Z_p\}$ for $\alpha \in T_p$.
Then $\Z_{(p), \alpha}$ is a DVR of $\Q(X)$ with maximal ideal $M_{(p), \alpha}$, 
$$\text{Int}(\Z) = \bigcap_{p \in \PP}(\bigcap_{\alpha \in T_p}(\Z_{(p), \alpha} \cap \Q[X])),$$
ht$(M_{(p), \alpha} \cap \text{Int}(\Z)) = 1$, and $\Z_{(p), \alpha}/M_{(p), \alpha} = \Z/p\Z$.
It is worthwhile to note that if $R =$ Int$(\Z)$, then $R = \bigcap\limits_{P \in X^1(R)} R_P$, $R_P$ is a DVR for all $P \in X^1(R)$,
but $R$ is not a Krull domain.

\subsection{Residually transcendental extensions}\label{res transcendental}

Let $V$ be a valuation domain with quotient field $K$,
$X$ be an indeterminate over $K$, and $K(X)$ be the field of rational function, i.e.,
the quotient field of $K[X]$. We denote by $v$ a valuation on $K$ with associated valuation domain $V$ and by $\Gamma_v$ the value group of $v$.  
We also denote by $\Gamma_{\ov}$ the divisible hull of $\Gamma_v$, namely, $\Gamma_{\ov}=\Gamma_v\otimes_{\Z}\Q$.

In this section we recall some known facts about residually transcendental extensions of $V$ to $K(X)$, that is, valuation domains $W$ of the field $K(X)$ lying above $V$ (i.e., $W\cap K=V$) such that the extension of residue fields is transcendental.  We denote by $\mathcal W_{\rt}$  the set of residually transcendental extensions of $V$ to $K(X)$. In contrast, we say that an extension $W$ of $V$ to $K(X)$ is residually algebraic if the extension of the residue fields is algebraic.

\subsubsection{Monomial valuations}\label{monomial valuations}

In this subsection we review the notion of monomial valuations, which is a particular class of extensions of $V$ to $K(X)$, in order to describe the class of residually transcendental extensions.

%A residually transcendental extension of $V$ to $K(X)$ is a valuation domain $W$ of $K(X)$ lying above $V$ (i.e., $W\cap K=V$) such that the extension of the residue fields is transcendental.  

\begin{Def}
Let $\alpha\in K$ and $\delta$ be an element of a value group $\Gamma$ which contains $\Gamma_v$. For $f\in K[X]$ such that $f(X)=a_0+a_1(X-\alpha)+\ldots+a_n(X-\alpha)^n$, we set:
$$v_{\alpha,\delta}(f)=\inf\{v(a_i)+i\delta \mid i=0,\ldots,n\}$$
The function $v_{\alpha,\delta}$ naturally extends to a valuation on $K(X)$ (\cite[Chapt. VI, \S. 10, Lemme 1]{Bourb}), 
and it is called a \emph{monomial valuation}. We denote by $V_{\alpha,\delta}$ the valuation domain of $K(X)$ associated to $v_{\alpha,\delta}$. 
Clearly, $V_{\alpha,\delta}$ lies over $V$.
Note that $V_{0,0} = V(X)$, i.e.,
the case $\alpha=0$ and $\delta=0$ corresponds to the trivial Gaussian extension. Note also that $V_{\alpha,\delta}\cap K[X]=V[\frac{X-\alpha}{c}]$, where $c\in K$ has value $\delta$ (see for example \cite[Remark 3.4]{PerPrufer}).
\end{Def}

Let $\overline{K}$ be the algebraic closure of $K$, 
$u$ be an extension of $v$ to $\overline{K}$, and 
$U$ be the valuation domain of $u$. Given 
$(\alpha,\delta)\in\overline{K}\times\Gamma_{\overline{v}}$, we denote by $v_{\alpha,\delta}$ the restriction of $u_{\alpha,\delta}$ to $K(X)$; note that this valuation depends on the chosen extension $u$ of $v$ to $\overline{K}$, see \cite[Remark 3.9]{PerPrufer}. If $V_{\alpha,\delta}$ ($U_{\alpha,\delta}$, resp.) is the valuation domain associated to $v_{\alpha,\delta}$ ($u_{\alpha,\delta}$, resp.) then clearly we have $V_{\alpha,\delta}=U_{\alpha,\delta}\cap K(X)$.
We also say that $(\alpha,\delta)\in\overline{K}\times\Gamma_{\overline{v}}$ is a \emph{minimal pair of definition of $V_{\alpha,\delta}$} (with respect to $K$)  if for each $\beta\in\overline{K}$ such that $u(\beta-\alpha)\geq \delta$, we have $[K(\beta):K]\geq[K(\alpha):K]$. Equivalently, among the set of pairs which define the valuation $u_{\alpha, \delta}$, the degree $[K(\alpha):K]$ is minimal (see \cite[Proposition 3]{AlePop} and \cite{APZTheorem}). We stress that when $\alpha\in\olK$, the coefficients $a_0,\ldots,a_n$ of the $(X-\alpha)$-adic expansion of $f\in K[X]$ lie in $\olK$ and may not be in $K$.


\begin{Rem}\label{ra rt}%{\rc expand....}\\
Suppose that $\delta>v(a)$ for every $a\in K$. We denote such $\delta$ by $\infty$. Given $\alpha\in \overline{K}$, for every $f\in K[X]$ with $f(X)=a_0+a_1(X-\alpha)+\ldots+a_n(X-\alpha)^n$ where $a_0,\ldots,a_n\in\olK$ we have
$$v_{\alpha,\delta}(f)=u(a_0)=u(f(\alpha))$$
so in this case $V_{\alpha,\infty}=V_{\alpha}=\{\phi\in K(X)\mid \phi(\alpha)\in U\}$.
\end{Rem}

\begin{Rem}\label{remark rt}
(1) Let $U$ be a fixed extension of $V$ to $\overline{K}$. 
Given $\alpha\in\overline{K}$ and an element $\delta$ in a totally ordered group containing $\Gamma_v$, 
it is well-known that $V_{\alpha,\delta}$ is a residually transcendental extension of $V$ 
if and only if $\delta$ belongs to $\Gamma_{\ov}$ (see for example \cite{APZTheorem}  or \cite[Lemma 3.5]{PerPrufer}). 
In this case, by \cite[Theorem 2.1]{APZTheorem}, the value group of $V_{\alpha,\delta}$ has finite index over the value group $\Gamma_v$ of $V$. 
%(so, in particular, it is contained in $\Gamma_{\ov}$). 
In particular, if $V$ is a DVR, then a residually transcendental extension of $V$ to $K(X)$ is a DVR, also.

(2) Let $V$ be a DVR and $U$ a fixed extension of $V$ to $\olK$. Given a (minimal) pair $(\alpha,\delta)\in\overline{K}\times\Gamma_{\overline v}$,
the intersection $V_{\alpha,\delta}\cap K[X]$ is clearly a Krull domain which is not Pr\"ufer, since the valuation overring $V_{\alpha,\delta}$ is residually transcendental over $V$, see \cite[Theorem 3.15]{PerPrufer}. 

It is known that if we restrict to the family of Krull domains $\{V_{\alpha,\delta}\cap K[X]\mid(\alpha,\delta)\in K\times\Gamma_{\overline v}\}$, this family has a natural ordering, namely:
\begin{equation}\label{inclusion}
V_{\alpha_1,\delta_1}\cap K[X]\subseteq V_{\alpha_2,\delta_2}\cap K[X]\Leftrightarrow \delta_1\leq\delta_2
\text{ and }v(\alpha_1-\alpha_2)\geq\delta_1
\end{equation}
(see \cite[Proposition 1.1]{APZAll}; see also \cite[Remark 3.10]{PerPrufer}), so 
\begin{equation}\label{equality monomials}
V_{\alpha_1,\delta_1}\cap K[X]= V_{\alpha_2,\delta_2}\cap K[X]\Leftrightarrow \delta_1=\delta_2 \text{ and } v(\alpha_1-\alpha_2)\geq \delta_1
\end{equation}
We will see in Lemma \ref{ordering} that the previous conditions are indeed equivalent to $V_{\alpha_1,\delta_1}= V_{\alpha_2,\delta_2}$.

 If $(\alpha,\delta)\in K\times \Gamma_{\overline{v}}$ and we consider the ball $B(\alpha,\delta)=\{x\in K\mid v(x-\alpha)\geq \delta\}$, then condition \eqref{inclusion} holds if and only if $B(\alpha_2,\delta_2)\subseteq B(\alpha_1,\delta_1)$.
In the case $V=\Z_{(p)}$ for some prime $p$, we generalize this ordering to the case of pairs $(\alpha,\delta)\in\C_p\times\Q$ in Section \ref{Polordering monomial}.
\end{Rem}

The main result of Alexandru, Popescu and Zaharescu is the following characterization of the  residually transcendental extensions of $V$ to $K(X)$ (see \cite{AlePop,APZTheorem,APZAll}).

\begin{Thm}\label{description rt Popescu}
Let $V$ be a valuation domain with quotient field $K$ and let $U$ be a fixed extension of $V$ to $\overline{K}$. 
%{\bc Does $U$ depend on $W$ ?} {\rc No, it does not; another choice %of the extension of $V$ to $\overline{K}$ leads to another %description of these extensions}.
Let $W$ be a residually transcendental extension of $V$ to $K(X)$. Then there exists a minimal pair
$(\alpha,\delta)\in\overline K\times\Gamma_{\overline{v}}$ such that $W$ is the restriction to $K(X)$
of the monomial valuation $U_{\alpha,\delta}$ of the field $\overline{K}(X)$.
\end{Thm}


Note that the valuation domain $U_{\alpha,\delta}$ of $\overline K(X)$ in Theorem \ref{description rt Popescu} is a residually transcendental extension of $U$ to $\overline{K}(X)$.
%\vskip0.4cm
%{\rc add pseudo-stationary sequences here and how to represent rt extensions by means of them}
%\vskip0.4cm
The next lemma is a generalization of \cite[Lemma 3.13]{PerPrufer} to valuation domains of arbitrary rank.

\begin{Lem}\label{rt overring of V[X]}
Let $U$ be an extension of $V$ to $\overline{K}$ and let $(\alpha,\delta)\in\overline{K}\times\Gamma_v$. Then
\begin{align*}
V[X]\subset U_{\alpha,\delta}\cap K(X) 
\Leftrightarrow u(\alpha)\geq0\;\text{ and }\delta\geq0.
\end{align*}
\end{Lem}
\begin{proof}
Note that the containment is equivalent to $u_{\alpha,\delta}(X)\geq0$ (and, in particular, equivalent to $U[X]\subset U_{\alpha,\delta}$). We have $u_{\alpha,\delta}(X)=\min\{\delta,u(\alpha)\}$ so that the claim is true.
\end{proof}


\subsubsection{Parametrization of residually transcendental extensions}\label{section parametrization}

In this subsection, we show how to parametrize $\mathcal W_{\rt}$ by means of ultrametric balls contained in $\overline{K}$ 
with respect to some fixed extension $U$ of $V$ to $\overline{K}$. 


Let $G_K$ be the absolute Galois group $\text{Gal}(\overline K/K)$. Note that $G_K$ is also equal to the Galois group of $\overline K(X)/K(X)$.  
Given an extension $U$ of $V$ to $\overline{K}$, we denote by $\mathcal{D}(U)$ the decomposition group of $U$, 
namely: $\mathcal{D}(U)=\{\sigma\in G_K \mid \sigma(U)=U\}$.


\begin{Lem}\label{image rt}
Let $U$ be a fixed extension of $V$ to $\overline{K}$.
Let $(\alpha,\delta)\in\overline{K}\times\Gamma_{\overline{v}}$, 
with $\delta=u(c)$ for some $c\in \overline{K}$, and $\sigma\in G_K$.  Then we have
$$\sigma(U_{\alpha,\delta})=\sigma(U)_{\sigma(\alpha),\delta}.$$
In particular, if $\sigma\in \mathcal{D}(U)$, then $\sigma(U_{\alpha,\delta})=U_{\sigma(\alpha),\delta}$.
\end{Lem}
\begin{proof}
Since $\frac{X-\alpha}{c}$ is a unit in $U_{\alpha,\delta}$ and  $U_{\alpha,\delta}$ is residually transcendental over the residue field $k_U$ of $U$, 
it follows that $\sigma\left(\frac{X-\alpha}{c}\right)=\frac{X-\sigma(\alpha)}{\sigma(c)}$ is a unit in $\sigma(U_{\alpha,\delta})$ 
and $\sigma(U_{\alpha,\delta})$ is residually transcendental over $k_{\sigma(U)}$ (note that $\sigma(U_{\alpha,\delta})$ is an extension of $\sigma(U)$ to $\overline K(X)$).
Recall that the valuation associated to $\sigma(U)=U'$ is equal to $u'=u\circ\sigma^{-1}$ (\cite[Proposition 3.2.16, (1)]{EngPre}), so $u'(\sigma(c))=\delta$.
Then the claim follows by \cite[Chapt. VI, \S 10, Proposition 2]{Bourb}. 

If $\sigma\in \mathcal{D}(U)$, then the last claim follows at once. 
\end{proof}
%Then, by , $u(\sigma(c))=u(c)=\delta$. In particular,  the previous equality becomes $\sigma(U_{\alpha,\delta})=U_{\sigma(\alpha),\delta}$.


We give next a criterion for when two residually transcendental extensions of $U$ to $\overline{K}(X)$ contracts down to the same valuation domain of $K(X)$ 
(which is a residually transcendental extension of $V$ to $K(X)$).
This has been done in a similar way in \cite{NN}.
In particular, given a residually transcendental extension $V_{\alpha,\delta}$ of $V$ to $K(X)$,  where $(\alpha,\delta)\in \overline{K}\times\Gamma_{\overline v}$, 
we describe the set of extensions of $V_{\alpha,\delta}$  to $\overline{K}(X)$ lying above $U$  (see also \cite[Theorem 2.2]{APZMinimal}). 

\begin{Thm}\label{Conjugate Rt ext}
Let $U$ be a fixed extension of $V$ to $\overline{K}$ and
let $(\alpha,\delta),(\alpha',\delta')\in\overline{K}\times\Gamma_{\overline{v}}$.
Then the following conditions are equivalent.
\begin{itemize}
\item[{\em (1)}] $U_{\alpha,\delta}\cap K(X)=U_{\alpha',\delta'}\cap K(X)$.
\item[{\em (2)}] $\sigma(U_{\alpha,\delta})=U_{\alpha',\delta'}$ for some $\sigma\in\mathcal{D}(U)$.
\item[{\em (3)}] $\delta=\delta'$ and $ u(\sigma(\alpha)-\alpha')\geq\delta$, for some $\sigma\in \mathcal{D}(U)$.
\end{itemize}
In particular, the set of extensions of a residually transcendental extension $W=U_{\alpha,\delta}\cap K(X)$ to $\overline{K}(X)$ lying above $U$ is equal to $\{U_{\sigma(\alpha),\delta}\mid \sigma\in \mathcal{D}(U)\}$.
%(((Furthermore, the number of such extensions is finite and  equal to the cardinality of $\{\alpha_{i_1},\ldots,\alpha_{i_k}\}\subseteq\Omega_{\alpha}=\{\alpha_1=\alpha,\ldots,\alpha_n\}$, where $u(\alpha_{i_j}-\alpha_{i_s})<\delta$ for each $j\not=s$ and for each $\alpha_k$ there exists (a unique) $j$ such that $u(\alpha_k-\alpha_{i_j})\geq\delta$\footnote{number of distinct balls of radius $\delta$ which form a disjoint covering of $\Omega_{\alpha}$}.)))
\end{Thm}

\begin{proof}
(1) $\Rightarrow$ (2) By  \cite[Conjugation Theorem 3.2.15]{EngPre},   there exists $\sigma\in G_K$ such that $\sigma(U_{\alpha,\delta})=U_{\alpha',\delta'}$. We claim that $\sigma\in \mathcal D(U)$. In fact, let $x\in U\subset U_{\alpha,\delta}$. Then $\sigma(x)\in\sigma(U_{\alpha,\delta})\cap \overline K=U_{\alpha',\delta'}\cap \overline K=U$. Thus, (2) holds.

(2) $\Rightarrow$ (3) By Lemma \ref{image rt}, $U_{\sigma(\alpha),\delta}= \sigma(U_{\alpha,\delta})= U_{\alpha',\delta'}$. 
By \cite[Proposition 1.1]{APZAll},  $\delta'=\delta$ and $u(\sigma(\alpha)-\alpha')\geq\delta$, which is condition (3).

(3) $\Rightarrow$ (1) The inequality $u(\sigma(\alpha)-\alpha')\geq\delta$ is equivalent to $\alpha'\in B_u(\sigma(\alpha),\delta)=\{x\in\olK \mid u(\sigma(\alpha)-x)\geq \delta\}$, so that $(\alpha',\delta)$ is another pair of definition for $U_{\sigma(\alpha),\delta}$ by \cite[Proposition 3]{AlePop} and thus $U_{\alpha',\delta}=U_{\sigma(\alpha),\delta}$. By  Lemma \ref{image rt}, we also have $U_{\sigma(\alpha),\delta}=\sigma(U_{\alpha,\delta})$, so $U_{\alpha,\delta}$  
and $U_{\alpha',\delta}$ contracts down to $K(X)$ to the same valuation domain. Thus, (1) holds.

The last claim follows immediately from the previous part and Lemma \ref{image rt}.
%In particular, the set of extensions of $W=U_{\alpha,\delta}\cap K(X)$ to $\overline{K}(X)$ is equal to $\{\sigma(U_{\alpha,\delta})=U_{\sigma(\alpha),\delta}\mid\sigma\in \mathcal D(U)\}$.
\end{proof}

For the next theorem, we fix the following notation: given an extension $U$ of  $V$ to $\overline{K}$, let $\mathcal{B}_u$ be the set of ultrametric balls $B_u(\alpha,\delta)=\{x\in\overline{K}\mid u(x-\alpha)\geq\delta\}$ in $\overline K$ with respect to $U$, where $(\alpha,\delta)\in \overline{K}\times\Gamma_{\overline{v}}$.
 It is worthwhile to note that condition (3) of Theorem \ref{Conjugate Rt ext} is equivalent to the following: 
$\sigma(B_u(\alpha,\delta))=B_u(\sigma(\alpha),\delta)=B_u(\alpha',\delta)$ for some $\sigma\in \mathcal{D}(U)$.

By Lemma \ref{image rt}, $\mathcal{D}(U)$ acts on $\mathcal{B}_u$, since for $\sigma\in \mathcal{D}(U)$ we have $\sigma(B_u(\alpha,\delta))=B_u(\sigma(\alpha),\delta)$. 
In particular, $\mathcal{D}(U)$ does not change the radius of an ultrametric ball.
We denote by $\mathcal{B}_u/\mathcal{D}(U)$ the set of equivalence classes of $\mathcal{B}_u$ under the action of the decomposition group $\mathcal{D}(U)$. 
By Lemma \ref{image rt}, an element of $\mathcal{B}_u/\mathcal{D}(U)$ is equal to $[B_u(\alpha,\delta)]=\{B_u(\sigma(\alpha),\delta)\mid\sigma\in\mathcal{D}(U)\}$; 
in particular, note that such an equivalence class is formed by finitely many closed balls (i.e., at most as many as the number of conjugates of $\alpha$ over $K$).

 

 
 

\begin{Thm}\label{parametrization rt}
Let $U$ be a fixed extension of $V$ to $\overline{K}$. Then the map
$\Theta:\mathcal{B}_u/\mathcal{D}(U) \to \mathcal W_{\rt}$ given by 
$[B_u(\alpha,\delta)] \mapsto U_{\alpha,\delta}\cap K(X)$ is bijective. 
\end{Thm}
\begin{proof}
In order to show that the map $\Theta$ is well-defined, we observe first that given $(\alpha,\delta)\in \overline{K}\times\Gamma_{\overline{v}}$, we have  $U_{\alpha,\delta}\cap K(X)\in \mathcal W_{\rt}$ by Remark \ref{remark rt}. Moreover, if $\sigma(B_u(\alpha,\delta))=B_u(\sigma(\alpha),\delta)\in [B_u(\alpha,\delta)]$ for some $\sigma\in \mathcal{D}(U)$, then $U_{\alpha,\delta}\cap K(X)=U_{\sigma(\alpha),\delta}\cap K(X)$ by Theorem \ref{Conjugate Rt ext}.

The surjectivity of $\Theta$ follows by Theorem \ref{description rt Popescu}.
Finally, the injectivity follows again by Theorem \ref{Conjugate Rt ext}. In fact, suppose that $\Theta([B_u(\alpha,\delta)])=\Theta([B_u(\alpha',\delta')])$, that is, $U_{\alpha,\delta}\cap K(X)=U_{\alpha',\delta'}\cap K(X)$. Then, by Theorem \ref{Conjugate Rt ext}, $\delta=\delta'$ and $u(\sigma(\alpha)-\alpha')\geq\delta$, so that $B_u(\alpha',\delta')=B_u(\sigma(\alpha),\delta)$. Hence, $[B_u(\alpha,\delta)]=[B_u(\alpha',\delta')]$.
\end{proof}
%{\rc generalizzare su $\Q_p$ considerando anche i DVR res. alg. ??}
%{\rc add equivalence with set of pseudo-stationary sequences??}

\section{Polynomial ordering of torsion extensions}\label{torsion extensions}

We say that an extension $W$ of $V$ to $K(X)$ is torsion if $\Gamma_w\subseteq\Gamma_{\overline{v}}$; we denote by $\mathcal W_{\t}$  the set of all torsion extensions of $V$ to $K(X)$. Note that a residually transcendental extension of $V$ to $K(X)$ is torsion by Remark \ref{remark rt},  
so $\mathcal W_{\rt} \subseteq \mathcal W_{\t}$. We stress that not all the residually algebraic extensions of $V$ to $K(X)$ are torsion. For example, $\{\phi\in K(X)\mid \phi(0)\in V\}$  is a residually algebraic extension of $V$ (its residue field is the same as the residue field of $V$) which is not torsion (its one-dimensional valuation overring is $K[X]_{(X)}$; see Proposition \ref{subvaluation of DVR} and also Remark \ref{why alpha transcendental}). 

\subsection{A partial order on the set of torsion extensions of $V$ to $K(X)$}


Let $W_1,W_2$ be in $\mathcal W_{\t}$. In \cite[page 282]{APZAll}, the following definition is given only for residually transcendental extensions (which are torsion):
$$w_1\leq w_2\Leftrightarrow w_1(f)\leq w_2(f),\forall f\in K[X]$$
where the last inequality is considered in
$\Gamma_{\overline{v}}$. More generally, we give the same definition for torsion extensions. 
We remark that since any $f\in K[X]$ can be written as $\frac{g(X)}{c}$ for some $g\in V[X]$ and $c\in V$, $c\not=0$, 
and the valuations $w_1,w_2$ coincide on $K$, the previous condition is equivalent to:  $w_1(g)\leq w_2(g)$ for each $g\in V[X]$.
Note that clearly we have
$$w_1=w_2\Leftrightarrow w_1(f)=w_2(f),\;\forall f\in K[X]\Leftrightarrow W_1=W_2.$$
In particular, we also say that $w_1<w_2$ if and only if $w_1\leq w_2$ and $w_1\not=w_2$.
%, that is:
%$$w_1<w_2\Leftrightarrow w_1\leq w_2,\;\text{ and }\exists f\in K[X]\;\text{ such that }w_1(f)<w_2(f).$$

 

\begin{Lem}\label{ordering}
Let $W_i \in \mathcal W_{\t}$ and $M_i$ be the maximal ideal of $W_i$ for $i=1,2$.
\begin{enumerate}
\item[{\em (1)}] The following statements are equivalent.
\begin{enumerate}
\item[{\em (a)}] $w_1\leq w_2$.
\item[{\em (b)}] $W_1\cap K[X]\subseteq W_2\cap K[X].$ 
\item[{\em (c)}] $M_1\cap K[X]\subseteq M_2\cap K[X]$.
\end{enumerate}
\item[{\em (2)}] $W_1=W_2$ if and only if $W_1\cap K[X]=W_2\cap K[X]$, if and only if $M_1\cap K[X] = M_2\cap K[X]$. 
\item[{\em (3)}] $(\mathcal W_{\t}, \leq)$ is a partially ordered set.
\end{enumerate}
\end{Lem} 

\begin{proof}
(1) (a) $\Rightarrow$ (c). Let $f\in M_1\cap K[X]$ be nonzero. Then $0 < w_1(f)\leq w_2(f)$, so $f\in M_2\cap K[X]$.
Thus $M_1\cap K[X]\subseteq M_2\cap K[X].$ 

(c) $\Rightarrow$ (b). Let $g\in W_1\cap K[X]$ be nonzero. If $g \in M_1$, then $g \in M_2 \cap K[X] \subseteq W_2 \cap K[X]$.
Next assume that $g \not\in M_1$ and $g \not\in W_2 \cap K[X]$. Since $W_2$ is a torsion extension of $V$,
there are $n \in \mathbb{N}$ and $b \in V$ such that $w_2(g^n) = nw_2(g) = -v(b) <0$ or $w_2(g^nb) =0$.
Note that $w_1(g^nb) = nw_1(g)+w_1(b) = v(b) >0$, so $g^nb \in M_1 \cap K[X] \subseteq M_2 \cap K[X]$
by assumption. Hence $w_2(g^nb) > 0$, a contradiction, which implies $g \in W_2 \cap K[X]$.
Thus $W_1\cap K[X]\subseteq W_2\cap K[X].$

(b) $\Rightarrow$ (a). Assume that $f\in K[X]$ is nonzero. Then either $f \in W_1$ or $\frac{1}{f} \in W_1$, and since
%{\rc [[simplified this proof by correcting a minor mistake: in general it is not true that  $V[X]\subseteq W_1\cap K[X]$]]}
%, $f(X)=\sum_i a_i X^i$. We may write $a_i=\frac{b_i}{c}$, for some $b_i,c\in V$; in particular, $f(X)=\frac{\sum_i b_i X^i}{c}$. 
%Therefore, we may restrict to polynomials $f\in V[X]$ in order to prove the inequality $w_1(f)\leq w_2(f)$. 
%Note that $V[X]\subseteq W_1\cap K[X]\subseteq W_2\cap K[X]$. In particular, for such a $f$, we have $w_1(f)\geq0$.
$W_1$ is torsion over $V$, there exist a positive integer  $n\in\N$ and $0 \neq c\in K$ such that $nw_1(f)=v(c)$. 
Hence, $w_1(\frac{f^n}{c}) = 0$. 
In particular, $\frac{f^n}{c}\in W_1\cap K[X]\subseteq W_2\cap K[X]$, so that we also have
$$0\leq w_2\left(\frac{f^n}{c}\right)\Leftrightarrow w_2(f)\geq\frac{v(c)}{n}=w_1(f),$$
so the inequality holds true. 

(2) It is clear that $W_1=W_2$ implies both of $W_1\cap K[X]=W_2\cap K[X]$ 
and $M_1\cap K[X]= M_2\cap K[X]$. Conversely, if $W_1\cap K[X]=W_2\cap K[X]$ or $M_1\cap K[X]= M_2\cap K[X]$,  
then $w_1(f)=w_2(f)$ for every $f\in K[X]$ by (1), which implies that $w_1=w_2$ on $K(X)$. Thus, $W_1=W_2$.

(3) This follows directly from (1) and (2) above.
\end{proof}

Let the notation be as in Lemma \ref{ordering}.
Then $w_1<w_2$ if and only if $W_1\cap K[X]\subset W_2\cap K[X]$,
if and only if $M_1\cap K[X]\subset M_2\cap K[X]$. The next corollary gives an alternative characterization of this strict inequality.
%, which are equivalent to say that $M_1\cap K[X]\subset M_2\cap(W_1\cap K[X])$ by the next corollary.

\begin{Cor}\label{w1<w2}
Let $W_i \in \mathcal W_{\t}$ and $M_i$ be the maximal ideal of $W_i$ for $i=1,2$. Then
the following statements are satisfied. 
\begin{enumerate}
\item[{\em (1)}] $w_1<w_2$ if and only if $M_1\cap K[X]\subset M_2\cap(W_1\cap K[X])$, and in this case, 
$M_1\cap K[X]$ is a prime non-maximal ideal of $W_1\cap K[X]$. 
\item[{\em (2)}] $W\in \mathcal{W}_{\t}$ is a maximal element of $(\mathcal W_{\t}, \leq)$ if and only if $W$ is residually algebraic.
%\item[{\em (3)}] {\rc CHECK  not true in general, consider an algebraically closed and spherically complete field} $W_i$ is contained in a maximal element, say $W_i'$, of $(\mathcal W_{\t}, \leq)$. Moreover, if $W_1$ and $W_2$ are not comparable, then $W_1' \neq W_2'$. 
\end{enumerate}
\end{Cor}
\begin{proof}
(1) Assume that $w_1<w_2$. Then, by Lemma \ref{ordering}, there exists $f\in (W_2\cap K[X])\setminus (W_1\cap K[X])$, so 
$w_1(f)<0\leq w_2(f)$. Since $W_1$ is a torsion extension of $V$, 
there exist $n\in\N$ and $0 \neq a\in K$ such that $w_1(f^n)=v(a)$. Then, we have
$$0=w_1\left(\frac{f^n}{a}\right)<w_2\left(\frac{f^n}{a}\right),$$
so $\frac{f^n}{a}\in M_2\cap (W_1\cap K[X])$ but $\frac{f^n}{a}\notin M_1\cap K[X]$. Conversely, suppose that 
$M_1\cap K[X]\subset M_2\cap(W_1\cap K[X])$. Then $M_2\cap(W_1\cap K[X]) \subseteq M_2 \cap K[X]$
implies $M_1\cap K[X]\subset M_2 \cap K[X]$. Thus $w_1<w_2$ by Lemma \ref{ordering}.

(2) $(\Rightarrow)$ Suppose that $W$ is residually transcendental. By Theorem \ref{description rt Popescu}, 
there exists $(\alpha,\delta)\in \overline{K}\times\Gamma_{\overline v}$ such that $W=U_{\alpha,\delta}\cap K(X)$ 
(where $U$ is a fixed extension of $V$ to $\overline K$). If $\delta'\in\Gamma_v$ 
is strictly greater than $\delta$, we have 
 $W\cap K[X] \subsetneq W' \cap K[X]$, where $W'=U_{\alpha,\delta'}\cap K(X)$. In fact, $U_{\alpha,\delta}\cap \overline K[X]\subsetneq U_{\alpha,\delta'}\cap\overline{K}[X]$ because $\frac{X-\alpha}{c'}$ is in the latter ring but not in the former (where $c'\in K$ has value $\delta'$) and by Theorem \ref{Conjugate Rt ext}, $U_{\alpha,\delta}\cap K[X]\subsetneq U_{\alpha,\delta'}\cap K[X]$.
 
$(\Leftarrow)$ Suppose that $W\in \mathcal{W}_{\t}$ is residually algebraic. We have the following containments, 
 where $M$ is the maximal ideal of $W$:
$$\frac{V}{M \cap K} \subseteq\frac{W\cap K[X]}{M\cap K[X]}\subseteq\frac{W}{M}.$$
It follows that $(W\cap K[X])/(M\cap K[X])$ is a field and thus $M\cap K[X]$ is a maximal ideal of $W\cap K[X]$. 
By the first part of the Corollary, $W$ is a maximal element of $(\mathcal W_{\t}, \leq)$.
%{\rc proof of (3)}
\end{proof}

It is worthwhile to note that if $K$ is algebraically closed and maximally complete (i.e., it does not admit any proper immediate extension), then it is straightforward to see that there are no residually algebraic torsion extensions of $V$ to $K(X)$.
Hence, in this case, $\mathcal W_{\t}$ has no maximal element by Corollary \ref{w1<w2}.

%Note that in Corollary \ref{w1<w2} we only use that $W_1$ and $W_2$ are torsion extensions of $V$ to $K(X)$.

\subsection{Residually transcendental extensions over the completion}\label{res transc completion}

%old title \subsection{Relationship between $\widetilde{\mathcal W}_{\rt}$ and ${\mathcal W}_{\rt}$} 


\bigskip

Let $\V$ (resp. $\K$) denote the $v$-adic completion of $V$ (resp. $K$), 
$\widetilde{\mathcal W}_{\t}$ be the set of all torsion extensions of $\V$  to $\K(X)$,
and $\widetilde{\mathcal W}_{\rt}$ be the set of all residually transcendental 
extensions of $\V$ to $\K(X)$. In this subsection,
we show that the map of $\widetilde{\mathcal W}_{\t}$ (resp. $\widetilde{\mathcal W}_{\rt}$)
into $\mathcal W_{\t}$ (resp. $\mathcal W_{\rt}$) given by $W \mapsto W \cap K(X)$ is bijective.

%\mathcal W_{\rt}$ has already been defined above

\begin{Lem}\label{contraction valuation}
Let $\wW$ be an extension of $\V$ to $\K(X)$ such that  $\tilde w(X)>v(a)$, for some $a\in K$. If $W=\wW\cap K(X)$, then we have
$$\frac{\wW\cap\K[X]}{\wM\cap\K[X]}=\frac{W\cap K[X]}{M\cap K[X]},$$
where $\wM,M$ are the maximal ideals of $\wW,W$, respectively.
\end{Lem}


\begin{proof}
Clearly, we have the containment $(\supseteq)$. For the other containment, 
we need to show that for every $f\in \wW\cap\K[X]$ there exists $g\in W\cap K[X]$ such that $\tilde w(f-g)>0$. 
Suppose that $f(X)=\sum\limits_{i= 0}^d \alpha_i X^i$. We choose $g\in K[X]$, $g(X)=\sum\limits_{i\geq0}^d a_i X^i$,
such that $v(\alpha_i-a_i)+iv(a)>0$, for each $i = 0,\ldots,d$. In particular,
\begin{equation}\label{f-g}
\tilde w(f-g)\geq \min\{v(\alpha_i-a_i)+i\tilde w(X)\}>\min\{v(\alpha_i-a_i)+iv(a)\}>0
\end{equation}
so that $f-g\in\wM$. Note that $g\in W$, so we get the claim.
\end{proof}

Let $\V_{\infty}=\{\psi\in\K(X)\mid \psi(\infty)\in\V\}$, where $\psi(\infty)$ is defined as $\varphi(0)$ with $\varphi(X)=\psi(1/X)$ (see \cite{PerTransc}). Then
the assumption on $\wW$ of Lemma \ref{contraction valuation} is equivalent to say $\wW\not=\V_{\infty}$. 
Note that every $W\in\mathcal{W}_{\t}$ satisfies the assumption of Lemma \ref{contraction valuation}.


\begin{Lem}\label{overring polynomial domain}
Let $\widetilde R$ be an integral domain such that $\V[X/a]\subseteq \widetilde R\subseteq\K[X]$ 
for some nonzero $a\in K$ and let $R=\widetilde R\cap K[X]$. Let $\wW$ be a valuation domain 
of $\K(X)$ extending $\V$  and let $W=\wW\cap K(X)$. Then
$R\subset W$ if and only if $\widetilde R\subset \wW$.
\end{Lem}
%{\rc NB: questo risultato puo' essere usato per mostrare che $\IntQ(S,\overline{\Z_{(p)}})$ e' Prufer (ossia non ha sovranelli di valutazione res. transc.) se e solo se $\Int_{\Q_p}(S_p,\oZp)$ e' Prufer (ossia non ha sovranelli di valutazione res. transc.).}
\begin{proof}
The implication $(\Leftarrow)$ is clear. Conversely, suppose that $R\subset W$ and let $f\in \widetilde R$. In particular, $f\in\K[X]=\K[X/a]$, so $f(X)=\sum\limits_{i=0}^d \alpha_i (X/a)^i$, for some $\alpha_i\in\K$. Let $g(X)=\sum\limits_{i=0}^d a_i (X/a)^i\in K[X]$ be such that $v(\alpha_i-a_i)\geq0$ for each $i=0,\ldots,d$. Then $h:=f-g\in \V[X/a]\subseteq\widetilde{R}$ so that $g=f-h\in \widetilde R\cap K[X]=R$. By assumption, $g\in W$ and note that $h\in \wW$, because $V[X/a]\subset W$ implies that $\V[X/a]\subset \wW$. Hence,  $f=h+g\in \wW$, as claimed.
\end{proof}

The following corollary follows at once.

\begin{Cor}\label{contraction W polynomial}
Let $\wW_i$ be an extension of $\V$ to $\K(X)$ such that $\tilde w_i(X)>v(a)$ for some $a\in K$, for $i=1,2$. Let $W_i=\wW_i\cap K(X)$, for $i=1,2$. Then we have
\begin{equation}\label{contaiments}
W_1\cap K[X]\subseteq W_2\cap K[X]\Leftrightarrow \wW_1\cap \K[X]\subseteq \wW_2\cap \K[X].
\end{equation}
\end{Cor}
\begin{proof}
It is sufficient to notice that the assumption of $
\tilde w_1$ implies that $\V[X/a]\subseteq\wW_1\cap\K[X]$ and then apply Lemma \ref{overring polynomial domain}.
\end{proof}

\begin{Lem}\label{extension rt ext}
Let $W$ be a residually transcendental extension of $V$ to $K(X)$
and $\widehat{K(X)}$ $($resp. $\W)$ be the completion of $K(X)$ $($resp. $W)$ with respect to $W$.
Then the following statements hold true:
\begin{itemize}
\item[{\em (1)}] $\K(X)$ is a subfield of $\widehat{K(X)}$.
\item[{\em (2)}] $\W$ is immediate over $\W\cap\K(X)$.
\item[{\em (3)}] $\W\cap\K(X)$ is a residually transcendental extension of $\V$.
\end{itemize} 
\end{Lem}
\begin{proof}
Since $\Gamma_w/\Gamma_v$ is torsion, it follows that $\K$ is equal to the topological closure of $K$ in $\widehat{K(X)}$. 
Hence, we may embed $\K(X)$ into $\widehat{K(X)}$. 

Let $\widetilde W=\W\cap \K(X)$.
Since $V\subseteq\V$ and $W\subseteq\widehat{W}$ are immediate extensions, it follows that $\widehat{W}$ is immediate over $\widetilde W$ and $\W$ is residually transcendental over $\V$. 
In particular, since $K(X)\subseteq\K(X)\subseteq\widehat{K(X)}$, it follows that also $W\subseteq\widetilde{W}$ is 
an immediate extension and so $\widetilde W$ is also a residually transcendental extension of $\V$. If $X$ were algebraic over $\K$, 
then  $\K\subseteq\K(X)$ would be a finite extension and so $\widetilde W$ could not be residually transcendental over $\V$. 
Then $X$ is transcendental over $\K$. 
\end{proof}



\begin{Lem}\label{contraction rt ext}
Let $\wW$ be a residually transcendental extension of $\V$ to $\K(X)$. Then $\wW\cap K(X)$ is a residually transcendental extension of $V$.
\end{Lem}
\begin{proof}
By \cite[Proposition 1]{AlePop}, there exists $f\in\K[X]\cap \wW$ such that its residue modulo $\widetilde M$ is transcendental over $\V/\M=V/M$. By Lemma \ref{contraction valuation}, 
if $g\in K[X]$ is a polynomial of the same degree as $f$ sufficiently $v$-adically close to $f$,  
then $g\in \wW\cap K(X)$ and $f$ and $g$ have the same residue modulo $\widetilde M$ so $g$ is residually transcendental over $V/M$. 
Thus, $\wW\cap K(X)$ is a residually transcendental extension of $V$.
\end{proof}

%Let $\widetilde{\mathcal W}_{\t}$  be the set of all torsion extensions of $\V$  to $\K(X)$.  already defined above
\begin{Thm}\label{contraction torsion exts}
If $\wW\in \widetilde{\mathcal W}_{\t}$ then $\wW\cap K(X)\in \mathcal W_{\t}$. Moreover,  if $\wW_i\in \widetilde{\mathcal W}_{\t}$ and $W_i=\wW_i\cap K(X)$ for $i=1,2$, then $w_1\leq w_2$ if and only if $\tilde w_1\leq \tilde w_2$. In particular, $W_1=W_2$ if and only if $\wW_1=\wW_2$.
\end{Thm}
\begin{proof}
It is clear that if $\wW\in\widetilde{\mathcal W}_{\t}$, then $\wW\cap K(X)\in \mathcal W_{\t}$, 
because $\Gamma_v=\Gamma_{\hat v}$.
Hence, by Lemma \ref{ordering}, the claim $w_1\leq w_2\Leftrightarrow\tilde w_1\leq \tilde w_2$ 
is equivalent to \eqref{contaiments}, which holds by Corollary \ref{contraction W polynomial}.

For the final claim, $(\Leftarrow)$ is straightforward. Conversely, suppose $W_1=W_2$. 
Then $W_1\cap K[X]= W_2\cap K[X]$, so by \eqref{contaiments}, $\wW_1\cap \K[X]=\wW_2\cap \K[X]$. 
Thus, by Lemma \ref{ordering}, $\wW_1=\wW_2$.
\end{proof}

%already defined above!
%Let $\widetilde{\mathcal W}_{\rt}$ ($\mathcal W_{\rt}$, resp.) be the set of all residually transcendental extensions of $\V$ ($V$, resp.) to $\K(X)$ ($K(X)$, resp.). 

\begin{Cor}\label{bijection rt}
The map $\Phi:\widetilde{\mathcal W}_{\rt} \to \mathcal W_{\rt}$ given by 
$\wW \mapsto \wW\cap K(X)$ is bijective. 
%\begin{align*}
%\Phi:\widetilde{\mathcal W}_{\rt}&\to \mathcal W_{\rt}\\
%\wW&\mapsto \wW\cap K(X).
%\end{align*}
\end{Cor}
\begin{proof}
Note that if $\wW\in\widetilde{\mathcal W}_{\rt}$, then $W=\wW\cap K(X)\in \mathcal W_{\rt}$ by Lemma \ref{contraction rt ext}.
Hence, $\Phi$ is well-defined. Then $\Phi$ is  surjective by Lemma \ref{extension rt ext}
and $\Phi$ is injective by Theorem \ref{contraction torsion exts}.
\end{proof}

\begin{Rem}
Let $\Psi$ be the map from $\widetilde{\mathcal W}_{\t}$ to $\mathcal W_{\t}$ which sends $\wW$ to $\wW\cap K(X)$. Clearly, the restriction of $\Psi$ to $\widetilde{\mathcal W}_{\rt}$ is equal to $\Phi$.  As in the proof of Corollary \ref{bijection rt}, $\Psi$ is well-defined and injective by Theorem \ref{contraction torsion exts}. However, $\Psi$ is not surjective: if $\alpha\in\K$ is transcendental over $K$, then $V_{\alpha}=\{\phi\in K(X) \mid \phi(\alpha)\in\V\}$ is a torsion extension of $V$, while $\V_{\alpha}=\{\psi\in \K(X) \mid \psi(\alpha)\in\V\}$ is not, since $\V_{\alpha}$ is contained in the DVR $\K[X]_{(X-\alpha)}$.
\end{Rem}

%DELETE THIS COROLLARY?
%\begin{Cor}
%Let $U$ be an extension of $V$ to $\overline{K}$ and let $\alpha,\beta\in\overline{K}$ be such that $V_{\alpha}\cap K[X]\subseteq V_{\beta}\cap K[X]$. Then $V_{\alpha}=V_{\beta}$.
%\end{Cor}
%\begin{proof}
%Same argument as in the proof of \cite[Theorem 3.2, iii)$\Rightarrow$ iv)]{PerTransc}.

%Let $p\in K[X]$ be the minimal polynomial of $\alpha$ over $K$. If $p(\beta)\not=0$, let $c\in V$ be such that $\frac{p(\beta)}{c}\notin U$. However, $\frac{p(X)}{c}\in V_{\alpha}\cap K[X]$ , so $\frac{p(X)}{c}\in V_{\beta}$ by assumption, a contradiction. Hence, $p(\beta)=0$, so $\alpha$ and $\beta$ are conjugate over $K$, so by \cite[Theorem 3.2]{PerTransc}, $V_{\alpha}=V_{\beta}$.
%\end{proof}



\section{Complete list of DVRs of $\Q(X)$}\label{list DVRs}

In this section we give a complete description of the family of DVRs with quotient field $\Q(X)$.

Let $\PP$ be the set of prime numbers in $\Z$, so $\{p\Z \mid p \in \PP\}$ is the set of nonzero (principal) 
prime ideals of the ring $\Z$ of integers.
Let $W\subset\Q(X)$ be a DVR with maximal ideal $M$. We have one of the following two cases:
\begin{itemize}
\item $W\cap \Q=\Q$.
\item $W\cap \Q=\Z_{(p)}$, for some $p\in\PP$.
\end{itemize}
In the first case we say that $W$ is a \emph{non-unitary} valuation domain, and in the second case, that $W$ is a \emph{unitary} valuation domain. %\footnote{For a given $p\in\PP$, we could use the expression "$p$-unitary valuation domain" to denote those unitary valuation domains lying above $\Z_{(p)}$}.
For each $p\in\PP$, let $\mathcal{W}_p$ be the set of DVRs $W$ of $\Q(X)$ such that $W \cap \Q = \Z_{(p)}$.

%\subsection{Non-unitary valuation domains}

It is a classical result (see for example \cite[p. 3]{Chevalley}) that if $W$ is a non-unitary valuation domain,
then $W$ is equal to one of the following DVRs, according to whether $X\in W$ or $X\not\in W$:
\begin{itemize}
\item[-] $W=\Q[X]_{(g)}$, for some irreducible $g\in\Q[X]$.
\item[-] $W=\Q[1/X]_{(1/X)}.$
\end{itemize}
%It works over any field $K$ in place of $\Q$.

%\subsection{Unitary valuation domains}
\noindent
If $W\in \mathcal{W}_p$ for some $p\in\PP$, 
%is a unitary valuation domain, i.e., $W\cap\Q=\Z_{(p)}$, for some  $p\in\PP$, 
we distinguish between these two subcases:
\begin{enumerate}
\item $W$ is residually algebraic over $\Z_{(p)}$.
%, that is, the extension of the residue fields is algebraic
\item $W$ is residually transcendental over $\Z_{(p)}$.
%, that is, the extension of the residue fields is transcendental.
\end{enumerate}

\noindent
Throughout this paper, for each $p\in\PP$ we set:
\begin{itemize}
    \item[-] $\mathcal{W}_{p,\ra} = \{W \in \mathcal{W}_p \mid W$ is residually algebraic over $\Z_{(p)}\}$;
    \item[-] $\mathcal{W}_{p,\rt} = \{W \in \mathcal{W}_p \mid W$ is residually transcendental over $\Z_{(p)}\}$.
\end{itemize}
In particular, $\mathcal{W}_{p} = \mathcal{W}_{p,\ra} \cup \mathcal{W}_{p,\rt}$ and  $\mathcal{W}_{p,\ra} \cap \mathcal{W}_{p,\rt}= \emptyset$.
Moreover, let 
$$\mathcal{W} = \bigcup\limits_{p \in \PP}\mathcal{W}_p,\;\; \mathcal{W}_{\ra} = \bigcup\limits_{p \in \PP}\mathcal{W}_{p,\ra},\;\;\mathcal{W}_{\rt} = \bigcup\limits_{p \in \PP}\mathcal{W}_{p,\rt},$$
so
$\mathcal{W}$ is the set of DVRs of $\Q(X)$,
$\mathcal{W}_{\ra}$ is the set of DVRs of $\Q(X)$ that are residually algebraic extensions of $\Z_{(p)}$ for some $p\in\PP$, and
 $\mathcal{W}_{\rt}$ is the set of DVRs of $\Q(X)$ that are residually transcendental extensions of $\Z_{(p)}$ for some $p\in\PP$.


\subsection{Extensions of $\Z_{(p)}$ to $\Q(X)$}


%\subsection{Residually algebraic extensions of $\Z_{(p)}$}\label{res alg Z_p}

Let $\Q_p,\Z_p$ be the field of $p$-adic numbers and the ring of $p$-adic integers, respectively, which are the completion of $\Q,\Z_{(p)}$ with respect to the topology induced by the $p$-adic valuation. Let $\oQp$ be a fixed algebraic closure of $\Q_p$ and $\oZp$ be the integral closure of $\Z_p$ in $\oQp$; since $\Q_p$ is Henselian, $\oZp$ is a one-dimensional non-discrete valuation domain. We denote by $v_p$ the unique extension of the $p$-adic valuation to $\oQp$. The ring $\oZp$ is the valuation domain of $v_p$ (i.e., $\oZp=\{x\in\oQp\mid v_p(x)\geq0\}$). Let $\C_p$ be the completion of $\oQp$ with respect to the topology induced by the valuation $v_p$ and $\O_p$ the completion of $\oZp$. It is well-known that $\C_p$ is algebraically closed (\cite[Proposition 5.7.8]{g97}).
Clearly, each element $\alpha\in\C_p\setminus\oQp$ is transcendental over $\Q_p$, and so, in particular, such an $\alpha$ is transcendental over $\Q$; on the contrary, an element  $\alpha\in\oQp$ can be either algebraic or transcendental over $\Q$. We denote by $\C_p^{\br}$ the subset of those $\alpha\in\C_p$ such that the ramification index of the extension $\Q_p(\alpha)/\Q_p$ is finite. We say that an element $\alpha\in\C_p^{\br}$ is unramified over $\Q_p$ if the ramification index of $\Q_p(\alpha)/\Q_p$ is $1$ (we stress that such an element may not be algebraic over $\Q_p$). Note that $\oQp\subset\C_p^{\br}\subset\C_p$. We also set $\O_p^{\br}=\O_p\cap\C_p^{\br}$.

%This subset is a field (CIT.)

\subsubsection{Residually algebraic extensions}


The following theorem characterizes the DVRs of $\Q(X)$ which are residually algebraic extensions of $\Z_{(p)}$, for some $p\in\PP$ (see \cite[Corollary 2.28]{PerStacked}).

\begin{Thm}\label{DVR ra}
Let $W$ be a DVR of $\Q(X)$ which is a residually algebraic extension of $\Z_{(p)}$ for some $p\in\PP$. 
Then there exists $\alpha\in\C_p^{\br}$, transcendental over $\Q$, such that
\begin{equation}\label{Zpalpha}
W=\Z_{(p),\alpha}=\{\phi\in\Q(X)\mid \phi(\alpha)\in \mathbb O_p\}.
\end{equation}
In particular, $\mathcal{W}_{p,\ra} = \{\Z_{(p), \alpha} \mid \alpha\in\C_p^{\br}$, $\alpha$ transcendental over $\Q\}$.
Furthermore,  $\alpha \in \oQp$ if and only if the extension of the residue fields $W/M\supseteq\Z/p\Z$ is finite.
\end{Thm}

%NEEDED?
\begin{Rem}\label{why alpha transcendental}
We remark that if $\alpha\in\oQp$, we may define in general the following valuation domain as above:
$$\Z_{(p),\alpha}=\{\phi\in\Q(X)\mid \phi(\alpha)\in \oZp\}$$
By \cite[Proposition 2.2]{PerTransc}, $\Z_{(p),\alpha}$ is a DVR precisely if $\alpha$ is transcendental over $\Q$, it has rank $2$ otherwise. In this last case, if $q\in\Q[X]$ is the minimal polynomial of $\alpha$, we have $\Z_{(p),\alpha}\subset\Q[X]_{(q)}$.
\end{Rem}

%{\bc Is it true that if $\mathcal{W}_{ra}$ is the set of DVRs of
%$\Q(X)$ which are residually algebraic extensions of $\Z_{(p)}$, then
%$\mathcal{W}_{ra} = \{\Z_{(p),\alpha} \mid \alpha\in\C_p^{\br}$, transcendental over $\Q\}$ ?}


\subsubsection{Residually transcendental extensions}


%\subsection{Residually transcendental extensions of %$\Z_{(p)}$}\label{res tr}


Let $W$ be a DVR of $\Q(X)$ which is residually transcendental over $\Z_{(p)}$ 
for some $p\in\PP$. 
%As we have already discussed in \S \ref{list DVRs}, $W$ can be either  residually transcendental or residually algebraic over $\Z_{(p)}$. 
The next important theorem gives an alternative characterization of these valuations than Theorem \ref{description rt Popescu}, 
namely, we show that we can parametrize them by means of pairs of $\oQp\times\Q$ rather than pairs of $\overline{\Q}\times\Q$.   

\begin{Thm}\label{rt over completion}
Let $p\in\PP$ and let $W$ be a residually transcendental extension  of $\Z_{(p)}$ to $\Q(X)$. 
Then $W$ is a DVR such that $W=\Z_{p,\alpha,\delta}\cap\Q(X)$ for some $(\alpha,\delta)\in\oQp\times\Q$.
\end{Thm} 
\begin{proof}
We have already said that $W$ is a DVR in Remark \ref{remark rt}. By Corollary \ref{bijection rt}, $W$ is equal to the restriction to $\Q(X)$ of a (unique) residually transcendental extension $\wW$ of $\Z_p$ to $\Q_p(X)$. By Theorem \ref{description rt Popescu}, $\wW=\Z_{p,\alpha,\delta}$ for some $(\alpha,\delta)\in\oQp\times\Q$. Hence, $W=\Z_{p,\alpha,\delta}\cap\Q(X)$. 
\end{proof}

The advantage of working over the completion of $\Q$ with respect to the $p$-adic valuation $v_p$ is that there exists a unique extension of $v_p$ to $\oQp$ since $\Q_p$ is Henselian. Therefore we don't have to worry about fixing an extension of $v_p$ to $\oQp$ in advance (as in Theorem \ref{description rt Popescu}).

%\vskip1.5cm
%\color{red}[NEW SUBSECTION??]\normalcolor\\
%If instead $W$ is a DVR of $\Q(X)$ which is residually algebraic over $\Z_{(p)}$, then by Theorem \ref{DVR ra} there exists $\alpha\in\C_p^{\br}$ transcendental over $\Q$,  such that $W=\Z_{(p),\alpha}$


\subsubsection{Unitary DVRs of $\Q(X)$}\label{res tr}

%extensions of $\Z_{(p)}$
%\subsection{Extensions of $\Z_{(p)}$ to $\Q(X)$ {\bc and 
%ultrametric balls in $\C_p$}}

In order to unify the aforementioned two cases, namely, residually algebraic 
and transcendental extensions of $\Z_{(p)}$ to $\Q(X)$, we consider  pairs $(\alpha,\delta)\in\C_p\times(\Q\cup\{\infty\})$. In the following we will always assume the following conditions for such a pair:
\begin{itemize}
    \item[-] if $\delta=\infty$, then $\alpha$ is transcendental over $\Q$. 
    \item[-] if $\delta\in\Q$, then $\alpha\in \oQp$.
\end{itemize}
%\vskip0.5cm
%we will always make the following assumption:
%\begin{center}
%{\bf for} $(\alpha,\delta)\in \C_p\times\Q$, $\alpha$ {\bf is in }$\oQp$.
%\end{center}
The first assumption is due to Theorem \ref{DVR ra}. The second assumption is due to the following fact : in the case of a monomial valuation $v_{p,\alpha,\delta}$ as introduced in section \ref{monomial valuations}, this valuation can also be defined when $\alpha\in\C_p\setminus\oQp$, but by a density argument we may take $\alpha'\in \oQp$ such that $v_{p,\alpha,\delta}=v_{p,\alpha',\delta}$ (see Remark \ref{ball containment} (2) below). 
Consequently, given a pair $(\alpha,\delta)\in \C_p\times(\Q\cup\{\infty\})$ with $\alpha\in\C_p\setminus\oQp$, automatically $\delta=\infty$. The converse is not true: we can consider pairs $(\alpha,\infty)$, where $\alpha$ is either in $\oQp$ or in $\C_p\setminus\oQp$. 
%We remark that the results of this section hold more generally over $\C_p$ rather than in the proper subset $\C_p^{\br}$.


\begin{Rem} \label{remark35}
Let $(\alpha,\delta)\in\C_p\times(\Q\cup\{\infty\})$ with the above assumptions. We recall that the value group of $\Z_{(p),\alpha,\delta}$ is of finite index over the value group of $\Z_{(p)}$ (by Remark \ref{remark rt} (1) if $\delta\not=\infty$ and by \cite[Proposition 2.2]{PerTransc}  if $\delta=\infty$). In particular, if $\delta\in\Q$, then $\Z_{(p),\alpha,\delta}$ is a DVR. Furthermore, if $\delta=\infty$, by \cite[Proposition 2.14, Remark 2.18, Corollary 2.28]{PerStacked}, $\Z_{(p),\alpha,\delta}$ is a DVR precisely when $\alpha\in\C_p^{\br}$ (we are assuming that $\alpha$ is transcendental over $\Q$). If $\alpha\in\C_p\setminus\C_p^{\br}$, then $\Z_{(p),\alpha,\infty}$ is a one-dimensional non-discrete valuation domain (by the same references).
\end{Rem}


Given $(\alpha,\delta)\in\C_p\times(\Q\cup\{\infty\})$, we denote by $\Z_{(p),\alpha,\delta}$ the valuation domain $\Z_{p,\alpha,\delta}\cap\Q(X)$. In particular,  $\Z_{(p),\alpha}=\Z_{(p),\alpha,\infty}$ by Remark \ref{ra rt}.
The next theorem summarizes the results discussed so far and give a complete description of the unitary DVRs of $\Q(X)$.

\begin{Thm}\label{unitary DVRs over Z(p)}
Let $W$ be a DVR of $\Q(X)$ lying above $\Z_{(p)}$, for some $p\in\PP$. Then
the following statements are satisfied.
\begin{enumerate}
\item[{\em (1)}] $W=\Z_{(p),\alpha,\delta}$ for some $(\alpha,\delta)\in\C_p^{\br}\times(\Q\cup\{\infty\})$. In this case, $X\in W$ if and only if $(\alpha,\delta)\in\O_p^{\br}\times(\Q_{\geq0}\cup\{\infty\})$.
\item[{\em (2)}] $W$ is residually transcendental (algebraic, respectively) over $\Z_{(p)}$ 
if and only if $\delta<\infty$ ( $\alpha$ is transcendental over $\Q$ and $\delta=\infty$, respectively).
\end{enumerate}
In particular, we have:
\begin{eqnarray*}
    \mathcal W_{p,\ra} &= &\{\Z_{(p),\alpha} \mid\alpha \in \C_{p}^{\br} \text{ and }\alpha \text{ is transcendental over }\Q\}\\
    \mathcal W_{p,\rt}  &= &\{\Z_{(p),\alpha,\delta}\mid (\alpha, \delta)\in\oQp\times\Q\}
\end{eqnarray*}
\end{Thm}

\begin{proof}
If $W$ is residually transcendental over $\Z_{(p)}$, then there exists $(\alpha,\delta)\in\oQp\times\Q$ 
such that $W=\Z_{(p),\alpha,\delta}$ (Theorem \ref{rt over completion}).
Next if $W$ is residually algebraic over $\Z_{(p)}$, then there exists $\alpha\in\C_p^{\br}$ 
such that  $\alpha$ is transcendental over $\Q$ and  $W=\Z_{(p),\alpha}$ (Theorem \ref{DVR ra}). 

The second part of item (1) follows by Lemma \ref{rt overring of V[X]}. The final claim follows immediately.
\end{proof}




%\begin{Cor} \label{coro37}
%For a fixed $p\in\PP$, the following statements hold.
%\begin{enumerate} 
%\item[{\em (1)}] $\mathcal W_{p,\ra} = \{\Z_{(p),\alpha} \mid\alpha \in \C_{p}^{\br} \text{ and }\alpha \text{ is transcendental over }\Q\}$.   
%\item[{\em (2)}] $\mathcal W_{p,\rt}  = \{\Z_{(p),\alpha,\delta}\mid (\alpha, \delta)\in\oQp\times\Q\}$.
%\item[\em (3)] If $\Z_{(p), \alpha, \delta}$ is a DVR containing $\Z_{(p)}[X]$,
%then $(\alpha,\delta)\in\O_p^{br} \times(\Q_{\geq 0}\cup\{\infty\})$.
%\end{enumerate}
% from (3):= \{\Z_{(p),\alpha,\delta}\mid (\alpha, \delta)\in\C_p^{\br}\times\Q\}
%\end{Cor}

%\begin{proof}
%(1) This follows from Theorem \ref{DVR ra}.
%(2) This follows from Theorem \ref{rt over completion}.
%\end{proof}
%{\rc the following sentence is already written in lines 7-9 of section %3.1}\\
%{\bc Recall that every $\alpha\in \C_p^{\br}\setminus\oQp$ is %transcendental over $\Q_p$, hence also over $\Q$; on the contrary, an %element $\alpha\in\oQp\subset\C_p^{\br}$ can be either algebraic or %transcendental over $\Q$.}

%Actually,  by our convention after Theorem \ref{rt over completion}, we have that $\mathcal W_{p,\rt} = \{\Z_{(p),\alpha,\delta}\mid (\alpha, \delta)\in\oQp\times\Q\}.$
%\end{proof}

\subsection{Polynomial ordering of unitary  valuation domains and ultrametric balls in $\C_p$}\label{Polordering monomial}

%in the set of extensions of $\Z_{(p)}$ to $\Q(X)$
%monomial valuations}

%PREVIOUS TITLE: \subsection{Polynomial ordering in $\{\Z_{(p),\alpha,\delta} \mid (\alpha,\delta)\in \C_p\times(\Q\cup\{\infty\})\}$}

For $p \in \PP$, we set 
$$\mathcal{S}_p = \{\Z_{(p),\alpha,\delta} \mid (\alpha,\delta)\in \C_p\times(\Q\cup\{\infty\})\}.$$ 
In this subsection, we study under which conditions we have $\Z_{(p),\alpha',\delta'}\cap\Q[X]\subseteq \Z_{(p),\alpha,\delta}$ for 
$(\alpha,\delta),(\alpha',\delta')\in\C_p\times(\Q\cup\{\infty\})$, and we show that $\mathcal{S}_p$ is a 
partially ordered set. 

%\bc 
%We first recall the following corollary which completely classifies all valuation rings on $\Q(X)$ of the form $\Z_{(p),\alpha,\delta}$ for $(\alpha,\delta)\in \C_p\times (\Q \cup \{\infty\})$.

%\begin{Cor} \label{coro38}
%Let $p\in\PP$ and $(\alpha,\delta)\in\C_p \times(\Q\cup\{\infty\})$. Then 
%the following statements hold: 
%\begin{enumerate}
%\item[\em (1)] $(\alpha,\delta)\in \C_p\times \Q$ if and only if $\Z_{(p),\alpha,\delta}$ is residually transendental over $\Z_{(p)}$. In this case, $\Z_{(p),\alpha,\delta}$ is a DVR and there is an $\alpha' \in \oQp$, so that 
%$\Z_{(p), \alpha, \delta} = \Z_{(p), \alpha', \delta}$.
%\item[\em (2)] $(\alpha,\delta)\in \C_p\times \{\infty\}$, i.e., $\alpha \in \C_p$ and $\delta = \infty$, if and only if $\Z_{(p),\alpha,\delta}$ is residually algebraic over $\Z_{(p)}$. In this case, we have the following three cases:
%\begin{enumerate}
%\item[\em (a)] $\alpha \in \C_p^{br}$ and $\alpha$ is transcendental over $\Q$ if and only if $\Z_{(p),\alpha,\delta}$ is a DVR.
%\item[\em (b)] $\alpha \in \C_p^{br}$ and $\alpha$ is algebraic over $\Q$ if and only if $\Z_{(p),\alpha,\delta}$ is two-dimensional. 
%In this case, if $q$ is the minimal polynomial of $\alpha$ over $\Q$, then $\Z_{(p),\alpha,\delta} \subset \Q[X]_{(q)}$.
%\item[\em (c)] $\alpha \in \C_p \setminus \C_p^{br}$ if and only if $\Z_{(p),\alpha,\delta}$ is a non-discrete rank-one valuation domain.
%\end{enumerate}
%\item[\em (3)] $\Z_{(p)}[X] \subseteq \Z_{(p), \alpha, \delta}$ if and only if $(\alpha,\delta)\in\O_p \times(\Q_{\geq 0}\cup\{\infty\})$.
%\item[\em (4)] If $\Z_{(p), \alpha, \delta}$ is a DVR containing $\Z_{(p)}[X]$,
%then $(\alpha,\delta)\in\O_p^{br} \times(\Q_{\geq 0}\cup\{\infty\})$. 
%\end{enumerate}
%\end{Cor}

%\begin{proof}
%(1) was noted at the beginning of Section \ref{res tr}. 
%(2) By (1) above, the first statement holds. Then (a) follows from Theorem \ref{unitary DVRs over Z(p)}. (b) If $\alpha \in \C_p^{br}$ and $\alpha$ is algebraic over $\Q$, then $\alpha \in \oQp$,
%and hence $\Z_{(p), \alpha, \delta}$ is
%two-dimensional. Conversely, if $\Z_{(p), \alpha, \delta}$ is 
%two-dimensional, then $\alpha \in \C_p^{br}$ by Remark \ref{remark35} 
%and $\alpha$ is algebraic over $\Q$ by (1) above. 
%Finally, the ``In this case" was noted in Remark \ref{why alpha transcendental}. 
%(c) follows from (a), (b) and Remark \ref{remark35}. 
%(3) is from Lemma \ref{rt overring of V[X]} and 
%(4) is an immediate consequence of (3) and 
%Theorem \ref{unitary DVRs over Z(p)}.  
%\end{proof}
%\ec

Let 
$$\mathcal{S} = (\bigcup\limits_{p \in \PP} \mathcal{S}_p) \cup \{\Q[X]_{(q)} \mid q \in \Q[X] \text{ is irreducible }\}.$$
It is noteworthy that if $W$ is a two-dimensional valuation overring of $\Z[X]$ with $W \subset \Z_{(p),0,0}$, then $W \neq \Z_{(p), \alpha, \delta}$ for all $(\alpha,\delta)\in\C_p \times(\Q\cup\{\infty\})$. Therefore $\mathcal{S}$ is not the set of all valuation rings of $\Q(X)$, while $\mathcal{S}$ contains the set of all DVRs of $\Q(X)$ by Theorem \ref{unitary DVRs over Z(p)}.

Now, in order to show the correspondence between the valuation domains of $\mathcal{S}_p$ and ultrametric balls  in $\C_p$, we need to consider ultrametric balls possibly having infinite radius (thus, a singleton). Let $(\alpha,\delta)\in\C_p\times(\Q\cup\{\infty\})$. The ball in $\C_p$ of center $\alpha$ and radius $\delta$ is $\bar B_p(\alpha,\delta)=\{x\in\C_p\mid v_p(x-\alpha)\geq\delta\}$. Note that $\bar B_p(\alpha,\delta)=\{\alpha\}$ precisely when $\delta=\infty$. If $\delta\in\Q$, 
then $\bar B_p(\alpha,\delta)$ has a non-empty intersection with $\oQp$, because $\oQp$ lies dense in $\C_p$, and so we may find an element in  $\bar B_p(\alpha,\delta)\cap\oQp$ which can be considered as the center of the ball.

We will show that the ordering on $\mathcal{S}_p$ reflects the containment among ultrametric balls in $\C_p$ under the action of the Galois group $G_p$ (Proposition \ref{overring intersection DVRs} and Remarks \ref{ball containment}).
We first need a couple of lemmas. The first lemma uses the notion of pseudo-stationary sequence introduced in \cite{ChabPolCloVal}, see also \cite{PerPrufer}.

\begin{Lem}\label{lem:ball containment}
Let $(\alpha,\delta)\in \C_p\times(\Q\cup\{\infty\})$ and suppose that $\bar B_p(\alpha,\delta)\subseteq \bigcup\limits_{j=1,\ldots,n}B_j$, 
where the $B_j$'s are polynomially closed subsets of $\C_p$. Then $\bar B_p(\alpha,\delta)\subseteq B_i$ for some $i\in\{1,\ldots,n\}$.
\end{Lem}
\begin{proof}
If $\delta=\infty$, then $\bar B_p(\alpha,\delta)=\{\alpha\}$ and the claim clearly is true.

Suppose that $\delta\in\Q$. Let $E=\{s_n\}_{n\in\N}\subset \bar B_p(\alpha,\delta)$ be a pseudo-stationary sequence with set of pseudo-limits $\mathcal{L}_E=\bar B_p(\alpha,\delta)$.  
By the pigeonhole principle, a subsequence of $E$, say $E'=\{s_{n_k}\}_{k\in\N}$, is contained in $B_i$, for some $i\in\{1,\ldots,n\}$; 
note that $E'$ is still pseudo-stationary with $\mathcal{L}_{E'}=\bar B_p(\alpha,\delta)$. By \cite[Proposition 4.8]{ChabPolCloVal}, 
$\bar B_p(\alpha,\delta)\subseteq B_i$, as claimed.
\end{proof}

\begin{Rem} More generally, Lemma \ref{lem:ball containment} holds for any valuation domain $V$ having infinite residue field; 
we don't need that the polynomial closure is a topological closure, 
which in general is not true for valuation domains of arbitrary rank (\cite{PSpolclo}).
\end{Rem}

We denote by $G_p$ the Galois group of $\oQp/\Q_p$. We recall that $G_p$ is identified with 
the group of all continuous $\Q_p$-automorphism of $\C_p$ (see \cite[\S 3, 1.]{APZclosed}) and so $G_p$ acts on $\C_p$. 
Moreover, for each $\sigma\in G_p$, we have $v_p\circ\sigma=v_p$, because $\Q_p$ is Henselian.

The following remark contains important informations about the set of conjugates of an element $\alpha\in\C_p\setminus\Q_p$ under the action of the group $G_p$.

\begin{Rem}\label{conjugate transcendental element}
Let $\alpha\in\C_p$. By \cite[Remark 3.2]{APZclosed}, $G_p(\{\alpha\})=\{\sigma(\alpha)\mid \sigma\in G_p\}$ is 
an infinite subset of $\C_p$ precisely when $\alpha\notin\oQp$. In this case, by \cite[Theorem 3.5]{APZclosed}, $G_p(\{\alpha\})$ 
is totally disconnected, closed and compact (note that this is true also when $G_p(\{\alpha\})$ is a finite set). 
In particular, by \cite[Proposition 2.4]{ChabPolCloVal}, $G_p(\{\alpha\})$ is polynomially closed.
\end{Rem}

\begin{Lem}\label{union balls pol closed}
Let $(\alpha_i,\delta_i)\in\C_p\times(\Q\cup\{\infty\})$, for $i=1,\ldots,n$. Then, $G_p(\bigcup\limits_{i=1}^n \bar B_p(\alpha_i,\delta_i))$ is polynomially closed.
\end{Lem}
\begin{proof}
We remark that
\begin{equation}\label{conjugates union balls}
G_p(\bigcup_{i=1}^n \bar B_p(\alpha_i,\delta_i))=\bigcup_{i=1}^n G_p(\bar B_p(\alpha_i,\delta_i))=\bigcup_{i=1}^n \bigcup_{\sigma\in G_p}\bar B_p(\sigma(\alpha_i),\delta_i).
\end{equation}

If $\alpha_i\in\oQp$, then $G_p(\bar B_p(\alpha_i,\delta_i))$ is equal to a finite union of balls $\bar B_p(\sigma(\alpha_i),\delta_i)$, 
for $\sigma\in G_p$ (which may not be necessarily disjoint); 
therefore, by \cite[Proposition 3.4]{ChabPolCloVal}, $G_p(\bar B_p(\alpha_i,\delta_i))$ is polynomially closed.
If $\alpha_i\in\C_p\setminus\oQp$ (so $\delta_i$ is necessarily equal to $\infty$ by our convention), 
then $G_p(\bar B_p(\alpha_i,\delta_i))=G_p(\{\alpha_i\})$ is polynomially closed by Remark \ref{conjugate transcendental element}.
Hence, $G_p(\bigcup\limits_{i=1}^n \bar B_p(\alpha_i,\delta_i))$ is polynomially closed, being a finite union of polynomially closed sets.
\end{proof}

We are now ready to prove the main result of this subsection by which we
can show that $\mathcal{S}_p$ is a partially ordered set. We need first the following remark, which gives the connection with rings of integer-valued polynomials.

\begin{Rem}\label{balls and intval rings}
  Let $(\alpha,\delta)\in\C_p\times(\Q\cup\{\infty\})$. If $\delta\not=\infty$, given $c\in\oQp$ such that $v_p(c)=\delta$ we have the following equalities:
$$
\O_{p,\alpha,\delta}\cap\C_p[X]=\O_p\left[\frac{X-\alpha}{c}\right]=\Int(\bar B_p(\alpha,\delta),\O_p).$$
The first equality follows from \cite[Remark 3.4]{PerPrufer} and the second one  
from \cite[Lemma 4.1]{PerPrufer}, since $\O_p$ is a rank one non-discrete valuation domain with infinite residue field. More generally, it is clear that the equality
\begin{equation}\label{restriction valuation domain on polynomial ring}
\O_{p,\alpha,\delta}\cap\C_p[X]=\Int(\bar B_p(\alpha,\delta),\O_p)
\end{equation}
is true also when $\delta=\infty$.
In particular, contracting down the previous equality to $\Q_p[X]$ (and similarly to $\Q[X]$), we have
\begin{equation}\label{IntQpB}
\Z_{p,\alpha,\delta}\cap\Q_p[X]=\Int_{\Q_p}(\bar B_p(\alpha,\delta),\O_p)
\end{equation}
%Note that when $\alpha\in \oQp$, then the right-hand side of \eqref{IntQpB} is equal to $\Int_{\Q_p}(\bar B_p(\alpha,\delta),\oZp)$. THIS is not true because $\bar B_p(\alpha,\delta)$ is a subset of $\C_p$ not of $\oQp$
Clearly, we have
\begin{equation}\label{IntQpB2}
\Int(\bar B_p(\alpha,\delta),\O_p)=\bigcap_{\alpha'\in \bar B_p(\alpha,\delta)}\O_{p,\alpha'}\cap\C_p[X]
\end{equation}
  
\end{Rem}

\begin{Prop}\label{overring intersection DVRs}
Let $p\in\PP$ and $(\alpha,\delta),(\alpha_i,\delta_i)\in\C_p\times(\Q\cup\{\infty\})$, for $i=1,\ldots,n$.
%and $R = \bigcap_{i=1}^n\Z_{(p),\alpha_i,\delta_i}\cap\Q[X]$.
Then the following conditions are equivalent:
\begin{itemize}
\item[{\em (1)}] $\Z_{(p),\alpha_i,\delta_i}\cap\Q[X]\subset\Z_{(p),\alpha,\delta}$ for some $i\in\{1,\ldots,n\}$.
\item[{\em (2)}] $\bigcap\limits_{i=1}^n\Z_{(p),\alpha_i,\delta_i}\cap\Q[X] \subset\Z_{(p),\alpha,\delta}$.
\item[{\em (3)}] $\delta\geq\delta_i$ and $v_p(\alpha-\sigma(\alpha_i))\geq\delta_i$ for some $\sigma\in G_p$ and $i\in\{1,\ldots,n\}$.
%{\bc \item[(iv)] Either ht$(M_{(p),\alpha,\delta}\cap R) = 2$ or $\Z_{(p),\alpha,\delta} = \Z_{(p),\alpha_i,\delta_i}$ for some $i\in\{1,\ldots,n\}$.}
\end{itemize}
\end{Prop}


%{\rc [[note that in the case $(\alpha_i,\delta_i)\in\oQp\times\{\infty\}$, then $\Z_{(p),\alpha_i,\delta_i}$ is a DVR if and only if
%$\alpha_i$ is transcendental over $\Q$ (see \cite[Proposition 2.2]{PerTransc}).
%Shall we always assume that or not? we can do that along with the other assumption at the very end of page 13.
%If you want to apply Corollary \ref{cor10}(1) for (ii) $\Leftrightarrow$ (iv), you need $R$ to be Krull.
%Moreover, are you implicitly assuming that $R = \bigcap_{i=1}^n\Z_{(p),\alpha_i,\delta_i}\cap\Q[X]$ is an irredundant representation?
%because if that is not the case then some $\Z_{(p),\alpha_i,\delta_i}$ could still be superfluous...]] }

\begin{proof}
 (1) $\Rightarrow$ (2) Clear. 
 
 (2) $\Rightarrow$ (3) By Lemma \ref{overring polynomial domain}, (2) is equivalent to the analogous containment in $\Q_p(X)$:
\begin{equation}\label{overring of intersection}
\bigcap_{i=1}^n(\Z_{p,\alpha_i,\delta_i}\cap\Q_p[X])\subset\Z_{p,\alpha,\delta}.
\end{equation}
By \eqref{IntQpB}, the left-hand side of \eqref{overring of intersection} is equal to
$$\bigcap_{i=1}^n\Int_{\Q_p}(\bar B_p(\alpha_i,\delta_i),\O_p)=\Int_{\Q_p}(B,\O_p),$$
where $B=\bigcup\limits_{i=1}^n \bar B_p(\alpha_i,\delta_i)$.

A straightforward adaptation of \cite[Lemma 2.20]{PerDedekind} shows that 
the integral closure of $\Int_{\Q_p}(B,\O_p)$ in $\oQp(X)$ is equal to $\Int_{\oQp}(G_p(B),\O_p)$. Hence, \eqref{overring of intersection} is actually equivalent to
$$\Int_{\oQp}(G_p(B),\O_p)\subset (\oZp)_{\alpha,\delta}.$$
By applying Lemma \ref{overring polynomial domain} again, we can transfer the previous containment over $\C_p$:
$$\Int(G_p(B),\O_p)\subset\O_{p,\alpha,\delta}\cap\C_p[X]=\Int(\bar B_p(\alpha,\delta),\O_p)$$
which amounts to say that $\bar B_p(\alpha,\delta)$ is contained in the polynomial closure of $G_p(B)$. By Lemma \ref{union balls pol closed}, $G_p(B)$ is polynomially closed, so $\bar B_p(\alpha,\delta)\subseteq G_p(B)$. By Lemma \ref{union balls pol closed} again, 
for each $i\in\{1,\ldots,n\}$, $G_p(\bar B_p(\alpha_i,\delta_i))$ is polynomially closed, 
so by Lemma \ref{lem:ball containment}, $\bar B_p(\alpha,\delta)\subseteq G_p(\bar B_p(\alpha_i,\delta_i))$ for some $i\in\{1,\ldots,n\}$. 
If $\delta_i=\infty$, then necessarily $\delta=\infty$ and $\alpha=\sigma(\alpha_i)$ for some $\sigma\in G_p$, 
because otherwise a similar argument to the proof of Lemma \ref{lem:ball containment} shows 
that $G_p(\{\alpha_i\})$ would contain a pseudo-stationary sequence which is not possible by compactness.  
If $\delta_i<\infty$, again Lemma \ref{lem:ball containment} implies that $\bar B_p(\alpha,\delta)$ is contained 
in $\sigma(\bar B_p(\alpha_i,\delta_i))=\bar B_p(\sigma(\alpha_i),\delta_i)$ for some $\sigma\in G_p$, 
which amounts to say that $\delta\geq\delta_i$ and $v_p(\alpha-\sigma(\alpha_i))\geq\delta_i$, which is precisely condition (3).

(3) $\Rightarrow$ (1) Note that $\Z_{(p),\alpha_i,\delta_i}=\Z_{(p),\sigma(\alpha_i),\delta_i}$ for each $\sigma\in G_p$  
by Theorem \ref{Conjugate Rt ext}. Note also that the inequality $v_p(\alpha-\sigma(\alpha_i))\geq\delta_i$ 
by \eqref{equality monomials} is equivalent to $\Z_{(p),\sigma(\alpha_i),\delta_i}=\Z_{(p),\alpha,\delta_i}$. Therefore,
$$\Z_{(p),\alpha_i,\delta_i}\cap\Q[X]=\Z_{(p),\alpha,\delta_i}\cap\Q[X]\subset\Z_{(p),\alpha,\delta}$$
where the last containment follows from the inequality $\delta\geq\delta_i$.
\end{proof}

%{\bc (ii) $\Leftrightarrow$ (iv) Each overring of $\Z[X]$ is of (Krull) dimension at most $2$,
%so the result follows directly from Corollary \ref{cor10}(1). }


\begin{Rems}\label{ball containment}
%(1) {\rc Note that condition (3) of Proposition \ref{overring intersection DVRs}
%amounts to say that $\bar B_p(\alpha,\delta)\subseteq \bar B_p(\sigma(\alpha_i),\delta_i)$, for some $\sigma\in G_p$.} (2) 
(1) Let $(\alpha_i,\delta_i)\in\C_p\times(\Q\cup\{\infty\})$, for $i=1,2$. By  Proposition \ref{overring intersection DVRs}, we have the following equivalence: 
$$\Z_{(p),\alpha_1,\delta_1}\cap\Q[X]\subset\Z_{(p),\alpha_2,\delta_2}\Leftrightarrow \delta_2\geq\delta_1 \;\text{ and } v_p(\alpha_2-\sigma(\alpha_1)) \geq \delta_1 \;\text{ for some }\sigma\in G_p.$$ 
This last condition is equivalent to saying that 
$$\bar B_p(\alpha_2,\delta_2)\subseteq\sigma(\bar B_p(\alpha_1,\delta_1))=\bar B_p(\sigma(\alpha_1),\delta_1).$$
In particular, the residually algebraic DVRs $\Z_{(p),\alpha}=\Z_{(p),\alpha,\infty}$ of $\mathcal{S}_p$ are precisely the maximal elements of $\mathcal{S}_p$ under this ordering (Corollary \ref{w1<w2} (2)) and each residually transcendental $\Z_{(p),\alpha,\delta}\in\mathcal{S}_p$, with $\delta\in\Q$, is strictly less than a maximal element of $\mathcal{S}_p$, namely, $\Z_{(p),\alpha',\infty}$ for $\alpha'\in\bar B_p(\alpha,\delta)$.

(2) Let $(\alpha_i,\delta_i)\in\C_p\times(\Q\cup\{\infty\})$ for $i=1,2$ be such that $\Z_{(p),\alpha_1,\delta_1}\cap\Q[X] = \Z_{(p),\alpha_2,\delta_2} \cap \Q[X]$. By Proposition \ref{overring intersection DVRs}, we have  
$\delta_1 = \delta_2$, $v_p(\alpha_2-\sigma_1(\alpha_1)) \geq \delta_1$, and 
$v_p(\alpha_1-\sigma_2(\alpha_2)) \geq \delta_2$ for some $\sigma_1, \sigma_2 \in G_p$.
If $\delta_1 = \infty$, then $\alpha_1 = \sigma(\alpha_2)$ for some $\sigma \in G_p$ (that is, $\alpha_1,\alpha_2$ are conjugate over $\Q_p$),
so, in particular, $\Z_{(p),\alpha_1,\delta_1} = \Z_{(p),\alpha_2,\delta_2}$. If instead we have $\delta_i < \infty$
for $i=1, 2$, then $\Z_{(p),\alpha_1,\delta_1} = \Z_{(p), \sigma(\alpha_1),\delta_1}$ by Theorem \ref{Conjugate Rt ext}
and $\Z_{(p), \sigma_1(\alpha_1),\delta_1} = \Z_{(p), \alpha_2,\delta_1}$ by (\ref{equality monomials}),
and hence again we get $\Z_{(p),\alpha_1,\delta_1} = \Z_{(p),\alpha_2,\delta_2}$.
Thus, $\mathcal{S}_p$ is a partially ordered set under $\Z_{(p),\alpha_1,\delta_1}\cap\Q[X] \subseteq \Z_{(p),\alpha_2,\delta_2} \cap \Q[X]$
for $\Z_{(p),\alpha_i,\delta_i} \in \mathcal{S}_p$.
\end{Rems}

%{\bc (3) Let $(\alpha_i,\delta_i)\in\C_p\times(\Q\cup\{\infty\})$ for $i=1,2$. Then $\Z_{(p),\alpha_1,\delta_1} = \Z_{(p),\alpha_2,\delta_2}$ if and only if
%$\delta_1 = \delta_2$ and $\alpha_2 \in \bar{B}_p(\sigma(\alpha_1), \delta_1)$ or $\alpha_1 \in \bar{B}_p(\sigma(\alpha_2), \delta_2)$ for some $\sigma \in G_p$ by (2) above and the proof of ``(3) $\Rightarrow$ (1)" in Proposition \ref{overring intersection DVRs}.}
 

 Let $\mathcal{B}_p$ be the set of ultrametric balls in $\C_p$  with center in $\C_p^{\br}$ and possibly infinite radius, namely:
 $$\mathcal{B}_p=\{\bar B_p(\alpha,\delta)\mid (\alpha,\delta)\in\C_p^{\br}\times(\Q\cup\{\infty\})\}$$
with the conventions that 
\begin{itemize}
    \item[-] if $\delta=\infty$, then $\alpha$ is transcendental over $\Q$.
    \item[-] if $\delta<\infty$, then $\alpha\in\oQp$.
\end{itemize}
Note that the Galois group $G_p=\Gal(\oQp/\Q_p)$ acts on $\mathcal{B}_p$ (Lemma \ref{image rt} and Theorem \ref{Conjugate Rt ext}), we denote by $\mathcal{B}_p/G_p$ the set of equivalence classes of $\mathcal{B}_p$ under this action.  Note that, by Remark \ref{conjugate transcendental element}, the cardinality of an equivalence class of a ball $\bar B_p(\alpha,\delta)$ under the action of $G_p$ is infinite precisely when $\delta=\infty$ and $\alpha\not\in\oQp$.  
\begin{Thm}\label{ultrametric balls modulo Galois are DVRs} 
The map
$$\begin{array}{lcc}
   \Theta_p:&\mathcal{B}_p/G_p\to & \mathcal{W}_p\\
   &{[} \bar{B}_p(\alpha,\delta){]}\mapsto &\Z_{(p),\alpha,\delta}
\end{array}$$
is a bijection.
\end{Thm}
\begin{proof}
   We show first that $\Theta_p$ is well-defined. If $[\bar{B}_p(\alpha,\delta)]=[\bar{B}_p(\alpha',\delta')]$ then there exists $\sigma\in G_p$ such that $\bar{B}_p(\alpha,\delta)=\sigma(\bar{B}_p(\alpha',\delta'))=\bar{B}_p(\sigma(\alpha'),\delta')$, so in particular we get $\delta=\delta'$ and $v_p(\alpha-\sigma(\alpha'))\geq\delta$. Hence, by Proposition \ref{overring intersection DVRs}, we have $\Z_{(p),\alpha,\delta}\cap\Q[X]=\Z_{(p),\alpha',\delta'}\cap\Q[X]$ and thus by Remark \ref{ball containment} (2) we get $\Z_{(p),\alpha,\delta}=\Z_{(p),\alpha',\delta'}$.   

  For the injectivity of $\Theta_p$, if the equality $\Z_{(p),\alpha,\delta}=\Z_{(p),\alpha',\delta'}$ holds, then in particular $\Z_{(p),\alpha,\delta}\cap\Q[X]=\Z_{(p),\alpha',\delta'}\cap\Q[X]$ and the same arguments as above show that $[\bar{B}_p(\alpha,\delta)]=[\bar{B}_p(\alpha',\delta')]$.

   The surjectivity of $\Theta_p$ follows from Theorem \ref{unitary DVRs over Z(p)}.
\end{proof}



\vskip1cm

\section{Polynomial Krull domains }
%as rings of  integer-valued polynomials

We are interested in characterizing polynomial Krull Domains, i.e., Krull domains $R$ such that $\Z[X]\subseteq R\subseteq\Q[X]$.
Note that since $\Z[X]$ is a $2$-dimensional Noetherian domain, by \cite[Theorem 9]{Heinzer} such a Krull domain $R$ is an integrally closed Noetherian domain of Krull dimension at most $2$. Conversely,
it is well-known that every integrally closed Noetherian domain is a Krull domain (for example, see \cite[Chapter VI, \S 13, (b), p. 82]{ZS2}). Hence, we are describing the family of integrally closed Noetherian domains between $\Z[X]$ and $\Q[X]$. Given a prime ideal $Q$ of an integral domain $R$ as above, we say that $Q$ is \emph{unitary} if $Q\cap\Z=p\Z$ for some $p\in\PP$ and that $Q$ is \emph{non-unitary} if $Q\cap\Z=(0)$.

\subsection{Representations of polynomial Krull domains} 


We recall that we have the following representation of the PID $\Q[X]$ as an intersection of DVRs:
$$\Q[X]=\bigcap_{q\in\mathcal{P}^{\text{irr}}}\Q[X]_{(q)}$$
where $\mathcal{P}^{\text{irr}}$ is the set of irreducible polynomials of $\Q[X]$  and $\Q[X]_{(q)}$ is the localization of $\Q[X]$ at the prime ideal generated by $q(X)$.

%between $\Z[X]$ and $\Q[X]$
In general, let $R$ be a polynomial Krull domain,
$\Lambda = \{p \in \PP \mid pR \subset R\}$, $X^1(R)$ be the set
of height one prime ideals of $R$, and $X^1_p(R) = \{Q \in X^1(R) \mid p \in Q\}$ for each $p \in \Lambda$.
Then $X^1_p(R)$ is finite, $R_Q$ is a DVR for all $Q \in X^1_p(R)$, and
$$R = \bigcap_{p \in \Lambda}(\bigcap_{Q \in X^1_p(R)}R_Q) \cap \Q[X].$$
Hence, for each $p \in \Lambda$, there exists a finite set $\{W_{p,1}, \dots , W_{p,m_p}\}$
%\footnote{{\rc I hope you don't mind if I invert the order of the indexes in this notation...}}
of DVRs with quotient field $\Q(X)$, so that $W_{p,j} \cap \Q = \Z_{(p)}$ for $j =1, \dots , m_p$, and
\begin{equation}\label{representation Krull domain}
R = \bigcap_{p \in \Lambda}(\bigcap_{j=1}^{m_p}W_{p,j}) \cap \Q[X] = \bigcap_{p \in \Lambda}(\bigcap_{j=1}^{m_p}(W_{p,j} \cap \Q[X])).
\end{equation}
For each $p \in \Lambda$, let 
\begin{equation}\label{localization R at p}
R_p = \bigcap_{j=1}^{m_p}(W_{p,j} \cap \Q[X])    
\end{equation}
It is clear that
$R_p$ is a Krull domain between $\Z_{(p)}[X]$ and $\Q[X]$ and $R = \bigcap\limits_{p \in \Lambda}R_p$.
\vskip0.5cm
%{\color{red} Question: do we have proved somewhere that the localization of $R$ as in \eqref{representation Krull domain} 
%at a fixed prime $p$ is equal to $R_p$, namely, $(\Z\setminus p\Z)^{-1}R=R_p$?? {\color{blue} No.}  I give a short proof just in case, 
%I think this fact needs to be emphasized in some way.}

%\begin{Lem}
%Let $p\in\PP$ and $R$ as in \eqref{representation Krull domain}. Let $R_{(p)}$ be the localization of $R$ at $p$, namely, $R_{(p)}=(\Z\setminus p\Z)^{-1}R$. Then
%$$R_{(p)}=R_p$$
%where $R_p$ is as in \eqref{localization R at p}.
%\end{Lem}


\begin{Lem} \label{lemma43}
Let $p\in\PP$ and $R$ be a ring between $\Z[X]$ and $\Q[X]$ such that
$R = \bigcap\limits_{Q \in X^1(R)}R_Q$ and $R_Q$ is a DVR for all $Q \in X^1(R)$, $X^1_p(R) = \{Q \in X^1(R) \mid p \in Q\}$ 
and $R_p = (\bigcap\limits_{Q \in X^1_p(R)}R_Q) \cap \Q[X]$. 
Then the following statements hold true. 
\begin{enumerate}
\item[{\em (1)}] $R = \bigcap\limits_{p \in \PP}R_p$.
\item[{\em (2)}] If $R_{(p)}=(\Z\setminus p\Z)^{-1}R$, then $R_{(p)}=R_p$.
\item[{\em (3)}] $R_{(p)} \cap \Q = \Z_{(p)}$.
\end{enumerate} 
\end{Lem}


%[Note that the proof below does not use that $R$ is a Krull domain...]{\color{blue} Your proof is correct, 
%and the result is true because each nonzero integer is a unit of all but a finite number of $R_p$'s}

\begin{proof}
(1) If $Q$ is a prime ideal of $R$, either $Q \cap \Z = (0)$ or $Q \cap \Z = p\Z$ for some $p \in \PP$.
Moreover, if $Q \cap \Z = (0)$, then $\Q[X] \subseteq R_Q$. 
Hence, $R = R \cap \Q[X] = \bigcap\limits_{p \in \PP}((\bigcap\limits_{Q \in X^1_p(R)}R_Q) \cap \Q[X])
= \bigcap\limits_{p \in \PP}R_p$.

(2) $(\subseteq)$ Let $\frac{f}{n}\in R_{(p)}$, where $f\in R$ and $n\in\Z$ is not divisible by $p$. Then clearly $\frac{f}{n} \in R_Q$, 
for every $Q \in X^1_p(R)$ since the value of each of the corresponding valuation of $n$ is zero and $R\subseteq R_p$. 
Thus, $R_{(p)} \subseteq R_p$. 
$(\supseteq)$ Conversely, let $h\in R_p$. Then we may write $h=\frac{g}{mp^k}$, 
for some $g\in\Z[X]$, $m\in\Z$ not divisible by $p$ and $k\geq 0$. 
Then $mh=\frac{g}{p^k}$ is clearly in $R$ (because for each $q\in\PP$, $q\not=p$, $p^k$ is a unit of $R_Q$ for all $Q \in X^1_q(R)$). 
Thus, we also have $R_{(p)} \supseteq R_p$.

(3) The result follows from the following two observations: (i) $R_Q \cap \Q = \Z_{(p)}$ for all $Q \in X^1_p(R)$
and (ii) $R_{(p)} = R_p$ by (2).
\end{proof}


Let $R$ be as in Lemma \ref{lemma43} and $R_p' = \bigcap\limits_{Q \in X^1_p(R)}R_Q$ for all $p \in \PP$. Then $R$ is a Krull domain
if and only if $R_p$ is a Krull domain for all $p \in \PP$ and, for any $0 \neq f \in \Z[X]$, 
$fR_p' = R_p'$ for all but a finitely many $p$'s in $\PP$. 
Hence, in the remaining part of this paper, we first study when $R_p$ is a Krull domain and we then study the Krull domain $R$ by investigating the property that if $f \in \Z[X]$ is nonzero, then $fR_p' = R_p'$ for all but a finitely many $p$'s in $\PP$. The first case will be called ``Local case" and the second one ``Global case".   




\normalcolor

\subsection{Local case}\label{section local case}
Let $p\in\PP$ be fixed. In this section we consider a Krull domain $R$ between $\Z_{(p)}[X]$ and $\Q[X]$; we show that $R$ can be represented as a ring of integer-valued polynomials over a finite union of ultrametric balls  in $\C_p$. Clearly, for every such $R$ we have
\begin{equation}\label{representation R}
R=\bigcap_{j=1}^m W_j\cap\Q[X]
\end{equation}
where $\{W_i\mid j=1,\ldots,m\}$ is a finite family (possibly empty) of DVRs of $\Q(X)$  lying above $\Z_{(p)}$.



\subsubsection{Irredundant representation}

We first investigate when a representation of a Krull domain between $\Z_{(p)}[X]$ and $\Q[X]$ as in \eqref{representation R} is irredundant. Let $D$ be an integral domain with quotient field $K$ and suppose that $D=\bigcap\limits_{i\in I} V_i$, where $\{V_i\}_{i\in I}$ is a family of valuation domains of $K$ (by Krull's Theorem this is equivalent to say that $D$ is integrally closed in $K$). As in \cite{gh68},
we say that an element $V=V_{i_0}$ of this family is superfluous if $D=\bigcap\limits_{i\in I,i\not=i_0}V_i$ and that the representation $D=\bigcap\limits_{i\in I} V_i$ is irredundant if no element $V_i$ of the family $\{V_i\}_{i\in I}$ is superfluous.

%{\rc I removed from the following statement "and $X^1(D)$ be the set of height one prime ideals" since it has already been defined at the beginning}
The following result is classical, see for example \cite[Remark on page 296]{n55} or \cite[Corollary 43.9]{Gilm}. 

\begin{Thm} \label{thm9}
Let $D$ be a Krull domain. Then $\{D_P \mid P \in X^1(D)\}$ is
the unique defining family for $D$, that is,
if $\{V_{\alpha}\}$ is a defining family for $D$, then
$\{D_P \mid P \in X^1(D)\} \subseteq \{V_{\alpha}\}$.
\end{Thm}

Let $D$ be a Krull domain.
Theorem \ref{thm9} also says that if $\Lambda$ is a subset of $X^1(D)$,
then $D \subsetneq \bigcap\limits_{P \in \Lambda}D_P$ if and only if $\Lambda \subsetneq X^1(D)$.
In particular, for every $P\in X^1(D)$, $D_P$ is not superfluous for the defining family $\{D_P\mid P\in X^1(D)\}$, 
which therefore realizes the unique irredundant representation of $D$.


\begin{Cor} \label{cor10}
Let $D$ be a Krull domain, $\{V_{\alpha}\}$ be a defining family for $D$
and $V \in \{V_{\alpha}\}$ with maximal ideal $M$.

%{\rc [[ I shortened the following 2 statements, since "$V \in \{V_{\alpha}\}$ with maximal ideal $M$" has already been introduced in the previous line.]]}

\begin{enumerate}
\item[{\em (1)}] $V$ is superfluous in $D = \bigcap\limits_{\alpha}V_{\alpha}$
if and only if ht$(M \cap D) \geq 2$.
\item[{\em (2)}] The intersection $D = \bigcap\limits_{\alpha}V_{\alpha}$ is irredundant if and only if ht$(M_{\alpha} \cap D)=1$ for all $\alpha$.
\end{enumerate}
\end{Cor}

\begin{proof}
(1) Suppose that $V$ is not superfluous in $D = \bigcap\limits_{\alpha}V_{\alpha}$
and let $R = \bigcap\limits_{V_{\alpha} \neq V}V_{\alpha}$. Then $D = R \cap V$ and $R \nsubseteq V$.
Hence $V = D_{M \cap D}$ \cite[Lemma 1.3]{HeinzerOhm}, and since $V$ is a DVR,
ht$(M \cap D) = $ht$((M \cap D)D_{M \cap D}) = 1$. Conversely, assume that ht$(M \cap D) < 2$.
Then ht$(M \cap D)=1$, and hence $D_{M \cap D}$ is a DVR.
Note that $D_{M \cap D} \subseteq V$, so $D_{M \cap D} = V$.
Thus $V$ is not superfluous in $D = \bigcap\limits_{\alpha}V_{\alpha}$.

(2) This follows directly from (1).
\end{proof}




\begin{Cor}\label{center superfluous DVR}
Let $R$ be as in \eqref{representation R} and let $W=W_j$ for some $j\in\{1,\ldots,m\}$. If $W'$ is a DVR extending $\Z_{(p)}$ such that $w<w'$, then $R\subset W'$ and the center of $W'$ on $R$ has height two.

Moreover, if $W''$ is another DVR extending $\Z_{(p)}$ such that $w<w'\leq w''$, then $M\cap R\subset M'\cap R=M''\cap R$.
\end{Cor}

\begin{proof}
%Let $M'$ be the maximal ideal of $W'$. 
By assumption that $w<w'$, we have $R\subseteq W\cap\Q[X]\subseteq W'\cap\Q[X]$.
Hence $W'$ is superfluous in $R \cap W'$, and since $W \neq W'$, we have ht$(M' \cap R) \geq 2$ by Corollary \ref{cor10}. Note that dim$(R) \leq 2$.
Thus ht$(M' \cap R) = 2$.

For the last statement, we have $R\subset W'\cap\Q[X]\subseteq W''\cap\Q[X]$, so $M'\cap R\subseteq M''\cap R$. Since dim$(R) \leq 2$, equality holds in the previous containment.
\end{proof}

\begin{Cor}\label{nonunitary not superfluous}
Let $R$ be as in \eqref{representation R}. Then, for every $q\in\Q[X]$ irreducible, $\Q[X]_{(q)}$ is not superfluous in that representation of $R$. %{\rc Moreover, $Q$  is a maximal ideal of $R$.}
\end{Cor}
\begin{proof}
Let $Q$ be the center of $\Q[X]_{(q)}$ on $R$.
Then $Q \cap \Z = (0)$, and hence if $S = \Z \setminus \{0\}$, then $R_S = \bigcap\limits_{i=1}^n(W_i)_S \cap \Q[X] = \Q[X]$
\cite[Proposition 43.5]{Gilm} and $QR_S \subsetneq R_S$. Thus ht$Q=$ ht$(QR_S) = 1$,
so, by Corollary \ref{cor10}, 2., $\Q[X]_{(q)}$ is not superfluous.
\end{proof}
%{\rc For the last claim, note that $\Q\subseteq R/Q\subseteq \Q[X]_{(q)}/q\Q[X]_{(q)}=\Q[X]/(q)$. Since $\Q\subseteq\Q[X]/(q)$ is an algebraic extension, it follows that $R/Q$ is a field, i.e. $Q$ is a maximal ideal.}



\begin{Rem}\label{irredundant represent balls}
(1) By Proposition \ref{overring intersection DVRs},  geometrically the representation of $R$ in \eqref{representation R} is irredundant if and only if the balls $\bar B_p(\alpha_i,\delta_i)$, $(\alpha_i,\delta_i)\in\O_p^{\br}\times(\Q_{\geq0}\cup\{\infty\})$, for $i=1,\ldots,n$, are disjoint under the action of the Galois group $G_p$ in the following sense: for each pair of balls $\bar B_p(\alpha_{i_1},\delta_{i_1}),\bar B_p(\alpha_{i_2},\delta_{i_2})$, there is no $\sigma\in G_p$ such that $\bar B_p(\alpha_{i_1},\delta_{i_1})$ is contained in $\sigma(\bar B_p(\alpha_{i_2},\delta_{i_2}))$ (see Remarks \ref{ball containment}). Equivalently,  the balls $\bar B_p(\sigma(\alpha_i),\delta_i)$, for $i=1,\ldots,n$ and $\sigma\in  G_p$, are disjoint. We stress that the number of such balls is finite if and only if there is no $j\in\{1,\ldots,n\}$ such that $\alpha_j\in\O_p^{\br}\setminus\oZp$ (see Remark \ref{conjugate transcendental element}).

(2) By Corollary \ref{cor10} and Corollary \ref{nonunitary not superfluous},   
the representation of $R$ in \eqref{representation R} is irredundant 
if and only if $M_j\cap R$ has height one for $j=1,\ldots,m$, where $M_j$ is the maximal ideal of $W_j$.
\end{Rem}

Now we study some basic ring-theoretic properties of $R = \bigcap\limits_{j=1}^m W_j\cap\Q[X]$, 
which are very useful for the study of polynomial Krull domains. In particular, we characterize when the prime unitary ideals of $R$ are maximal. If $W$ is a DVR of $\Q(X)$ extending $\Z_{(p)}$ for some $p\in\PP$, we denote by $e(W\mid \Z_{(p)})$ the ramification index of the extension $W\supset\Z_{(p)}$. If $\alpha\in\C_p^{\br}$, we denote by $e(\Q_p(\alpha)\mid\Q_p)$ the ramification index of the extension $\Q_p(\alpha)\supseteq\Q_p$ (we stress that this extension is algebraic precisely when $\alpha\in \oQp$).
% between $\Z[X]$ and $\Q[X]$

\begin{Thm} \label{characterization local Krull domain}
Let $p\in\PP$. Suppose that $W_j=\Z_{(p),\alpha_j,\delta_j}$, for some $(\alpha_j,\delta_j)\in\C_p^{\br}\times(\Q\cup\{\infty\})$, 
is a DVR of $\Q(X)$ extending $\Z_{(p)}$ and $X \in W_j$ for $j = 1, \dots , m$. Now, let
 $M_j=M_{(p),\alpha_j,\delta_j}$ for each $j=1,\ldots,m$, and
\begin{equation}\label{representation R2}
R = \bigcap_{j=1}^m W_j\cap\Q[X].
\end{equation}
%{\rc how do we take into account the case $R=\Q[X]$ here? that is, $m$ can also be $0$}
\noindent
If the representation of $R$ in \eqref{representation R2} is  irredundant (namely,  ht$(M_j \cap R)=1$ for $j=1, \dots , m$), then the following statements hold:
\begin{enumerate}
\item[{\em (1)}] $R$ is a Krull domain such that $\Z_{(p)}[X] \subseteq R \subseteq \Q[X]$.
\item[{\em (2)}] $X^1(R) = \{M_j \cap R \mid j =1, \dots , m\} \cup \{f\Q[X] \cap R \mid f \in \Q[X]$ is irreducible$\}$.
\item[{\em (3)}] $M_{j} \cap R$ is a maximal ideal of $R$ if and only if $W_j$
is residually algebraic over $\Z_{(p)}$.
\item[{\em (4)}] Suppose that $(\alpha_j,\delta_j)$ is a minimal pair if $\delta_j\in\Q$. Then we have
 $$e(W_j|\Z_{(p)}) = \left\{\begin{array}{cc}
e(\Q_p(\alpha_j)|\Q_p),&\text{ if } \delta_j = \infty\\
e(\Q_p(\alpha_j)|\Q_p)e(\gamma_j|\Q_p(\alpha_j)) ,&\text{ if } \delta_j < \infty
\end{array}\right.$$
where $\gamma_j=w_j(f_{\alpha_j})$, $f_{\alpha_j}$ is the minimal polynomial of $\alpha_j$ over $\Q_p$ and $e(\gamma_j|\Q_p(\alpha_j))$ 
is the smallest natural number $n$ such that $n\gamma_j$ is in the value group of $\Q_p(\alpha_j)$.

\item[{\em (5)}] If $e_j = e(W_j|\Z_{(p)})$ for $j=1, \dots , m$ and $d= gcd(e_1, \dots , e_m)$, then
Cl$(R) = \Z/d\Z \oplus \Z^{m-1}$.

%{\bc
%\item[{\em (6)}] $R$ is a UFD if and only if either $m=0$ (in which case $R=\Q[X]$) or $m=1$, $\alpha_1$ is unramified over $\Q_p$ and either $\delta_1=\infty$  or $\gamma_1\in\Z$. }
%{\rc FIX here!!! it is not necessarily $\gamma_1\in\Z$,?? it is sufficient, though...}.

\item[{\em (6)}] $R$ is Dedekind if and only if $\delta_j=\infty$  and $\alpha_i$ is transcendental over $\mathbb{Q}$ for each $j=1,\ldots,m$.
\end{enumerate}
\end{Thm}
\begin{proof}
(1) Clear, since $R$ is a finite intersection of Krull domains with the same quotient field (\cite[Chapter VII, \S 1, 3. Exemples, 3)]{Bourb}).

(2) The containment $(\supseteq)$ follows by \cite[Corollary 2.3(3)]{Chang}
and the assumption that $M_j\cap R\in X^1(R)$ for all $j\in\{1,\ldots,m\}$.
The other containment follows easily by  Theorem  \ref{thm9}.

(3) If $W_j$ is residually transcendental over $\Z_{(p)}$, that is, $\delta_j\in\Q$ (Theorem \ref{unitary DVRs over Z(p)}), then the claim follows by Corollary \ref{center superfluous DVR}. Suppose instead that $W_j$ is residually algebraic over $\Z_{(p)}$, that is, $\delta=\infty$ (Theorem \ref{unitary DVRs over Z(p)}).  We have the following containments: $\Z/p\Z\subseteq R/(M_j\cap R)\subseteq W_j/M_j$; since $W_j/M_j$ is algebraic over $\Z/p\Z$, it follows that $M_j\cap R$ is a maximal ideal of $R$.

(4) If $\delta_j<\infty$, see \cite[Theorem 2.1, c)]{APZTheorem}, while if $\delta_j=\infty$, see \cite[Proposition 2.2]{PerTransc}.

(5) Let $G$ be a free abelian group of rank $m$ on the generators $g_1, \dots, g_m$
and $H$ be the subgroup of $G$ generated by $r = \sum\limits_{i=1}^me_ig_i$. Then Cl$(R) = G/H$ \cite[Theorem 2.4]{Chang},
and since $d= gcd(e_1, \dots , e_m)$, Cl$(R) = \Z/d\Z \oplus \Z^{m-1}$
\cite[Remark 1.6]{Chang}.

%(6) It is known that a Krull domain is a UFD if and only if it has a trivial divisor class group\cite[Proposition 6.1]{Fossum}. The proof then follows from items (4) and (5) above.

(6) A Krull domain is Dedekind if and only if it has (Krull) dimension one (\cite[ Theorem 28, Chapt. VI, \S 13]{ZS2}). The proof then follows from items (2), (3) and Theorem \ref{unitary DVRs over Z(p)}.
\end{proof}

%{\bc \begin{Rem}
%Let $R$ be  as in Theorem \ref{characterization local Krull domain}(4). 
%If $e(W_j|\Z_{(p)}) = 1$, then either $\delta_j = \infty$ or $\delta_j < \infty$ and $e(\gamma_j|\Q_p(\alpha_j)) =1$. Now assume $\delta_j < \infty$ and let $\Omega$ is the set of conjugates of $\alpha_j$ in $\oQp$, so that $f_{\alpha_j} = \prod\limits_{\alpha \in \Omega}(X- \alpha)$. Then $\gamma_j = w_j(f_{\alpha_j}) = \sum\limits_{\alpha \in \Omega}w_j(X-\alpha) = v_p(c) + k\delta_j$ for some $c \in \oQp$ and $k=\#\{\alpha\in\bar B_p(\alpha_j,\delta_j)\mid \alpha \in \Omega\}$. Hence, as $v_p(c) \in \Z$, then $\gamma_j \in \Z$ if and only if $k\delta_j \in \Z$.
%Therefore, Theorem \ref{characterization local Krull domain}(6) can be restated as follows: $R$ is a UFD if and only if either $m=0$ or $m=1$, $\alpha_1$ is unramified over $\Q_p$ and either $\delta_1=\infty$ or  $k\delta_1\in\Z$ for $k=\#\{\alpha\in\bar B_p(\alpha_1,\delta_1)\mid f_{\alpha_1}(\alpha)=0\}$. Unfortunately, we don't know whether $\delta_1 \in \Z$ or there is a UFD $R$ with $\delta_1 \not\in \Z$.
%where {\rc FIX here!!! it is not necessarily $\gamma_1\in\Z$,?? it is sufficient, though...}.\end{Rem} }

\begin{Def}
An infinite family $\{W_i\}_{i\in I}$ of DVRs of $\Q(X)$ will be said to be
\emph{polynomially finite} if there exists a finite set $\{V_1, \dots , V_n\}$ of DVRs of $\Q(X)$ such that
$\bigcap\limits_{i\in I}W_i\cap \Q[X]=\bigcap\limits_{j=1}^n V_j \cap\Q[X]$.
\end{Def}

For example, putting together \eqref{restriction valuation domain on polynomial ring} and \eqref{IntQpB2}, we have
$$\bigcap_{\alpha'\in\bar B_p(\alpha,\delta)}\O_{p,\alpha',\infty}\cap\C_p[X]=\O_{p,\alpha,\delta}\cap\C_p[X].$$
Note that an element of the family $\{\O_{p,\alpha',\infty}\mid\alpha'\in \bar B_p(\alpha,\delta)\}$ is a DVR
if and only if $\alpha'$ is transcendental over $\Q$ (\cite[Proposition 2.2]{PerTransc}), so in particular, not every of the valuation overring $\O_{p,\alpha_p,\infty}$ above is a DVR.

The next theorem establishes when an infinite family $\{W_i\}_{i\in I}$ of DVRs of $\Q(X)$ extending $\Z_{(p)}$ is polynomially finite.


%, namely:
%$$\bigcap_{i\in I}W_i\cap \Q[X]=\bigcap_{j=1}^m W_j'\cap\Q[X]$$
%where $\{W_j'\}_{j=1,\ldots,m}$ is a finite family of DVRs extending $\Z_{(p)}$? Obviously, without loss of generality we may assume that $\{W_j'\}_{j=1,\ldots,m}\subset\{W_i\}_{i\in I}$. We let then $W_j'=W_j$.

\begin{Thm}\label{polynomially finite}
Let $p\in\PP$. Suppose that $\Z_{(p),\alpha_i,\delta_i}$, for some $(\alpha_i,\delta_i)\in\C_p^{\br}\times(\Q\cup\{\infty\})$, 
is an infinite family of DVRs of $\Q(X)$ extending $\Z_{(p)}$, for $i\in I$, 
and let $R = \bigcap\limits_{i\in I}\Z_{(p),\alpha_i,\delta_i}\cap \Q[X]$.
Then the following statements are equivalent:
\begin{enumerate}
\item[{\em (1)}] $R$ is a Krull domain.
\item[{\em (2)}] The set $\{\Z_{(p),\alpha_i,\delta_i} \mid i \in I\}$ is polynomially finite.
\item[{\em (3)}] $R = \bigcap\limits_{j=1}^n \Z_{(p),\alpha_j,\delta_j}\cap\Q[X]$ for some $\{(\alpha_j,\delta_j) \mid j =1, \dots ,n\} \subseteq \C_p^{\br}\times(\Q\cup\{\infty\})$.
\item[{\em (4)}] There exists a finite subset $\{(\alpha_j,\delta_j) \mid j =1, \dots ,n\}$ of $\C_p^{\br}\times(\Q\cup\{\infty\})$, 
so that $R \subseteq \bigcap\limits_{j=1}^n\Z_{(p), \alpha_j, \delta_j}$
and for each $i\in I$, there exists $j\in\{1,\ldots,n\}$ such that $\delta_i\geq\delta_j$ and $v(\alpha_i-\sigma(\alpha_j))\geq \delta_j$,
for some $\sigma\in G_p$.
\end{enumerate}
%$\C_p^{\br}\times(\Q\cup\{\infty\})$
\end{Thm}

\begin{proof}
(1) $\Rightarrow$ (2) Let $\mathcal{A} = X^1_p(R)$. Then $\mathcal{A}$ is finite,
$R = \bigcap\limits_{Q \in \mathcal{A}}R_Q \cap \Q[X]$, $R_Q$ is a DVR of $\Q(X)$ extending $\Z_{(p)}$.
Thus $\{\Z_{(p),\alpha_i,\delta_i} \mid i \in I\}$ is polynomially finite.


(2) $\Rightarrow$ (3) This follows from the fact that each DVR
of $\Q(X)$ extending $\Z_{(p)}$ is of the form $\Z_{(p),\alpha,\delta}$ for some $(\alpha, \delta) \in \C_p^{\br}\times(\Q\cup\{\infty\})$
by Theorem \ref{unitary DVRs over Z(p)}.

(3) $\Rightarrow$ (1) Clear, since $R$ is a finite intersection of Krull domains with the same quotient field (\cite[Chapter VII, \S 1, 3. Exemples, 3)]{Bourb}).

(3) $\Rightarrow$ (4) Each $\Z_{(p),\alpha_i,\delta_i}$ is an overring of $\bigcap\limits_{j=1}^n \Z_{(p),\alpha_j,\delta_j}\cap\Q[X]$, 
so by Proposition \ref{overring intersection DVRs}, $\delta_i\geq\delta_j$ and $v(\alpha_i-\sigma(\alpha_j))\geq \delta_j$, 
for some $j\in\{1,\ldots,m\}$ and $\sigma\in G_p$.

(4) $\Rightarrow$ (3) Suppose that for each $i\in I$, there exists $j\in\{1,\ldots,n\}$ such that $\delta_i\geq\delta_j$ 
and $v(\alpha_i-\sigma(\alpha_j))\geq \delta_j$, for some $\sigma\in G_p$. 
Then by Proposition \ref{overring intersection DVRs}, $\Z_{(p),\alpha_i,\delta_i}$ is an overring of $\bigcap\limits_{j=1}^n \Z_{(p),\alpha_j,\delta_j}\cap\Q[X]$, 
so the containment $(\supseteq)$ holds and the equality follows because $R \subseteq \bigcap\limits_{j=1}^n\Z_{(p), \alpha_j, \delta_j}$.
\end{proof}



\begin{Rem}
Geometrically, condition 4. of Theorem \ref{polynomially finite}
means that the union of the balls $\bar B_p(\alpha_i,\delta_i)$ in $\oQp$, for $i\in I$ (each of them corresponds to a valuation domain $\Z_{(p),\alpha_i,\delta_i}$),  is contained in finitely many of them up to the action of the Galois group $G_p$, namely:
$$\bigcup_{i\in I}\bar B_p(\alpha_i,\delta_i)\subseteq \bigcup_{j=1,\ldots,m}\bigcup_{\sigma\in G_p}\bar B_p(\sigma(\alpha_j),\delta_j).$$
\end{Rem}
%\subsection{Polynomial Krull Domains of $\Z_{(p)}$}
%\vskip0.5cm

% and Corollary \ref{coro38}

We can now characterize Krull Domains $R$ between $\Z_{(p)}[X]$ and $\Q[X]$. By Theorem \ref{unitary DVRs over Z(p)}, we may restrict to consider in the defining family of DVRs of $R$  only those DVRs extending $\Z_{(p)}$ defined by  pairs $(\alpha,\delta)\in\O_p^{\br}\times(\Q_{\geq0}\cup\{\infty\})$.
We recall that  $\Z_{(p)}[X]$ is a Krull domain (for example see \cite[Proposition 1.6]{Fossum}).
If $\Z_{(p),0,0}$ is the valuation domain of the trivial Gaussian extension, then $\Z_{(p)}[X]=\Z_{(p),0,0}\cap\Q[X]=\IntQ(\O_p)$.

%More generally, given $(\alpha,\delta)\in\oQp\times\Q$, $\Z_{(p),\alpha,\delta}\cap\Q[X]$ is a Krull domain, being an intersection of finitely many Krull domains (see \cite[Chapt. VII, \S 1, 3. Examples, p. 201]{Bourb}).



\begin{Thm} \label{thm4.23}
Let $p\in\PP$ and  $R$ be a Krull domain  between $\Z_{(p)}[X]$ and $\Q[X]$. Then there exist finitely many pairs $(\alpha_i,\delta_i)\in\O_p^{\br}\times(\Q_{\geq0}\cup\{\infty\})$,  $i=1,\ldots,n$, with  $\alpha_i$ transcendental over $\Q$ if $\delta_i=\infty$, such that
\begin{equation}\label{representation Krull local}
    R=\bigcap_{i=1}^n\Z_{(p),\alpha_i,\delta_i}\cap\Q[X]=\IntQ\left(E_p,\O_p\right),
\end{equation}
where $E_p = \bigcup\limits_{i=1}^n \bar B_p(\alpha_i,\delta_i)$. Moreover, 
if the representation in \eqref{representation Krull local} is irredundant (see Remark \ref{irredundant represent balls}),
then $R$ is Dedekind if and only if either $n=0$ (and then $R=\Q[X]$) or $n\geq1$ and  $\delta_i=\infty$ for each $i=1,\ldots,n$;  thus, in either case, $R$ is Dedekind if and only if $E_p$ contains finitely many elements.
%we have the following:
%\begin{itemize}
%\item[-] $R$ is a UFD if and only if $n=0$ (and then $R=\Q[X]$) or $n=1$, $\alpha_1$ is unramified over $\Q_p$ and {\rc $\gamma_1\in\Z\cup\{\infty\}$}. {\bc either $\delta_1 = \infty$ or $\gamma_1=v_{p, \alpha_1, \delta_1}(f_{\alpha_1}) \in \Z$ for the minimal polynomial  $f_{\alpha_1}$ of $\alpha_1$ over $\Q_p$. ??}    \item[-] $R$ is Dedekind if and only if either $n=0$ (and then $R=\Q[X]$) or $n\geq1$ and  $\delta_i=\infty$ for each $i=1,\ldots,n$;  thus, in either case, $R$ is Dedekind if and only if $E_p$ contains finitely many elements.
%\end{itemize}
\end{Thm}
\begin{proof}
We remarked at the beginning of this section in \eqref{representation R} that $R=\bigcap\limits_{i=1}^n W_i\cap\Q[X]$, where $W_i$ is a DVR of $\Q(X)$ containing $\Z_{(p)}[X]$, for each $i=1,\ldots,n$.  By Theorem \ref{unitary DVRs over Z(p)},  there exist $(\alpha_i,\delta_i)\in\O_p^{\br}\times(\Q_{\geq0}\cup\{\infty\})$ such that $W_i=\Z_{(p),\alpha_i,\delta_i}$, for each $i=1,\ldots,n$. Note that, if $\delta_i=\infty$, then necessarily $\alpha$ must be transcendental over $\Q$ (\S \ref{res tr}). The representation of $R$ as a ring of integer-valued polynomials follows by \eqref{IntQpB}.

Clearly, $R\cap\Q$ is equal either to $\Q$ or to $\Z_{(p)}$ according to whether the family of DVRs $W_i$ is empty or not. In the former case, $R=\Q[X]$ which is a PID (thus, Dedekind and a UFD). In the latter case, by \cite[Theorem 3.14]{PerPrufer}, the Krull domain $R$ is Dedekind if and only if there is no $W_i$ which is residually transcendental over $\Z_{(p)}$. This corresponds to saying that $\delta_i=\infty$ for each $i=1,\ldots,n$, by Theorem \ref{unitary DVRs over Z(p)}.

%It is well-known that a Krull domain is a UFD if and only if it has a trivial divisor class group.  The stated characterization then follows by Theorem \ref{characterization local Krull domain} (6).  
\end{proof}


\subsubsection{Non-unitary prime ideals}\label{section: non-unitary prime ideals}

Let $R$ be a polynomial Krull domain between $\Z_{(p)}[X]$ and $\Q[X]$ as in \eqref{representation R}.
%, where $W$ is a residually transcendental extension of $\Z_{(p)}$ to $\Q(X)$. By Theorem \ref{rt over completion}, we know  that $W=\Z_{(p),\alpha,\delta}$ for some $(\alpha,\delta)\in\oQp\times\Q$.
Let $q\in\Q[X]$ be irreducible and let $Q=q(X)\Q[X]_{(q)}\cap R$. By Theorem \ref{characterization local Krull domain} (2), $Q$ is an height 1 prime ideal of $R$ and $R_Q=\Q[X]_{(q)}$. In this section we characterize when $Q$ is a maximal ideal of $R$.

In general, we remark that if $W$ is a valuation domain of the field of rational functions $K(X)$ over a field $K$ and $W$ has dimension at least $2$, then $W\cap K=V$ is a valuation domain of $K$ which is not $K$ itself. In fact, $W\cap K=K\Leftrightarrow K\subset W$ and in this case $W$ would be a DVR (\cite[p. 3]{Chevalley}).

\begin{Prop}\label{subvaluation of DVR}
Let $W$ be a valuation domain of $K(X)$ such that $W\subset K[X]_{(q)}$ for some irreducible $q\in K[X]$. Let $V=W\cap K$. Then, there exist a root $\alpha\in\olK$ of $q(X)$ and an extension $U$ of $V$ to $K(\alpha)$ such that $W=U_{\alpha}=\{\phi\in K(X)\mid \phi(\alpha)\in U\}$.
\end{Prop}
\begin{proof}
 Let $\pi:K[X]_{(q)}\to K[X]/(q)\cong K(\alpha)$ be the canonical residue map, where $\alpha$ is the residue class of $X$ and can be identified to a root of $q$. Note that for $\phi\in K[X]_{(q)}$, $\pi(\phi)=\phi(\alpha)$ because $\pi$ is an homomorphism.

It is well-known that every valuation domain $W$ strictly contained in $K[X]_{(q)}$ is of the form $\pi^{-1}(U)$, where $U$ is a valuation domain of $K(\alpha)$; moreover, $U=\pi(W)=W/Q$, where $Q$ is the maximal ideal of $K[X]_{(q)}$.  Hence,
\begin{equation}\label{residue W}
U=W/Q=\{\phi(\alpha)\mid \phi\in W\}=\{\phi(\alpha)\in K(\alpha)\mid \phi(\alpha)\in U\}.    
\end{equation}
(Actually, every element of $K(\alpha)$ is of the form $g(\alpha)$ for some $g\in K[X]$). Moreover, the maximal ideal of $U$ is equal to $M_W/Q$, where $M_W$ is the maximal ideal of $W$. In particular, $U\cap K=V$. In fact, if $\beta\in V=W\cap K$, then $\pi(\beta)=\beta$ because $\pi_{\mid K}$ is the identity. Hence, $\beta=\pi(\beta)\in U=W/Q$. Similary, one can show that $M_V\subseteq M_U=M_W/Q$, the maximal ideal of $U$. Therefore, $V=U\cap K$.

Finally, by \eqref{residue W} and the above considerations, 
$$W=\pi^{-1}(U)=\{\phi\in K(X)\mid \phi(\alpha)\in U\}=U_{\alpha}.$$
\end{proof}

We remark that if we consider the completion of $K(\alpha)$ with respect to the valuation domain $U$, we may consider $\alpha$ as an element of $\overline{\widehat{K}}$, algebraic over $K$, so that we have the equality:
$$W=\{\phi\in K(X)\mid \phi(\alpha)\in \widehat{U}\}$$
where $\widehat{U}$ is the completion of $U$.

\begin{Thm}\label{nonunitary prime ideals}
  Let $p\in\PP$ and let $R=\bigcap\limits_{j=1}^m \Z_{(p),\alpha_j,\delta_j}\cap \Q[X]$ be a Krull domain, where $(\alpha_j,\delta_j)\in \C_p\times (\Q\cup\{\infty\})$ for $j=1,\ldots,m$.  Let $Q=q(X)\Q[X]\cap R$ be a non-unitary prime ideal of $R$, for some irreducible $q\in \Q[X]$. Then $Q$ is a maximal ideal of $R$ if and only if $q$ has no root in $\bigcup\limits_{j=1}^m \bar B_p(\alpha_j,\delta_j)$.  
\end{Thm}
\begin{proof}
Note that when $R$ is Dedekind (that is, $\delta_j=\infty$ for $j=1,\ldots,m$ by Theorem \ref{characterization local Krull domain}, (7)), then $Q$ is automatically maximal. In this case each ball $\bar B_p(\alpha_j,\delta_j)$ is the singleton $\{\alpha_j\}$ which is necessarily transcendental over $\Q$ because otherwise $\Z_{(p),\alpha_j,\infty}$ would not be a DVR.

In general, suppose that $q(\alpha')=0$ for some $\alpha'\in\bigcup_{j=1}^m \bar B_p(\alpha_j,\delta_j)$. Then, we have these containments:
$$R\subset\Z_{(p),\alpha'}=\Z_{(p),\alpha',\infty}\subset\Q[X]_{(q)}$$
where the first containment follows by Proposition \ref{overring intersection DVRs} and the second one from the fact that $\Q[X]_{(q)}=\{\phi\in K(X)\mid \phi(\alpha')\in\oQp\}$. In particular, $Q\subset M_{(p),\alpha'}\cap R$, so $Q$ is not a maximal ideal.

Conversely, suppose that $Q\subset M\subset R$ for some height $2$ maximal ideal $M$ of $R$, which necessarily contains $p$.  By \cite[Corollary 19.7]{Gilm}, there exists a $2$-dimensional valuation overring $W$ of $R$ with non-zero prime ideals $Q'\subset M'$ such that $Q'\cap R=Q$ and $M'\cap R=M$. Note that $W\cap \Q=\Z_{(p)}$ and $W_{Q'}=R_Q=\Q[X]_{(q)}$. By Proposition \ref{subvaluation of DVR}, there exist an extension $U$ of $\Z_{(p)}$ to $\Q[X]/(q)$  and a root $\alpha'$ of $q(X)$ such that $W=\{\phi\in\Q(X)\mid \phi(\alpha')\in U\}$. By the remark after that Proposition, we may assume that $\alpha'\in\oQp$, necessarily algebraic over $\Q$, so that $W=\Z_{(p),\alpha'}=\Z_{(p),\alpha',\infty}$. Since $R\subset \Z_{(p),\alpha'}$, it follows by Proposition \ref{overring intersection DVRs} that $\alpha'$ (or possibly some of its conjugate) is in $\bigcup_{j=1}^m \bar B_p(\alpha_j,\delta_j)$, as claimed.
\end{proof}

\begin{Cor}
Let $R$ and $Q$ be as in Theorem \ref{nonunitary prime ideals} and suppose that $Q$ is not a maximal ideal. Then the maximal ideals of $R$ containing $Q$ are unitary and they are finitely many.
\end{Cor}
\begin{proof}
    If $Q\subset M\subset R$ where $M$ is a maximal ideal, the above proof of Theorem \ref{nonunitary prime ideals} shows that $M=M_{(p),\alpha'}\cap R$, where $\alpha'$ is a root of $q(X)$. Hence, such $M$'s are finitely many.
 \end{proof}
\vskip0.2cm

\subsubsection{Unique factorization domains}\label{section UFD local}

In this section we completely describe the UFDs between $\Z_{(p)}[X]$ and $\Q[X]$, where $p\in\Z$ is a fixed prime.
 It is known that an integral domain is a UFD if and only if it is a Krull domain with trivial divisor class group \cite[Proposition 6.1]{Fossum}.

The following lemma follows from \cite[Proposition 3.2 \& Proposition 3.7]{PS2}, but for the sake of the reader we give a self-contained proof.

\begin{Lem}\label{conjugates in the ball}
Suppose $\alpha\in\oQp$ is unramified over $\Q_p$ and let $q_{\alpha}\in\Q_p[X]$ be its minimal polynomial. Then, $v_{p,\alpha,\delta}(q_{\alpha})=k\delta+\gamma$, for some $\gamma\in\Z$ and $k=\#(\Omega_{\alpha}\cap\bar B_p(\alpha,\delta))$, where $\Omega_{\alpha}$ is the set of conjugates of $\alpha$ over $\Q_p$.
\end{Lem}
\begin{proof}
We remark that $\Omega_{\alpha}\subset\Q_p(\alpha)$ because $\Q_p(\alpha)/\Q_p$ is a Galois extension. Note also that $q_{\alpha}(X)=\prod\limits_{\alpha'\in\Omega_{\alpha}}(X-\alpha')$ over $\Q_p(\alpha)$.  
 For each $\alpha'\in\Omega_{\alpha}$, we have
$$v_{p,\alpha,\delta}(X-\alpha')=\min\{\delta,v_p(\alpha-\alpha')\}$$
so that 
$$v_{p,\alpha,\delta}(q_{\alpha})=\sum_{\alpha'\in\Omega_{\alpha}\cap\bar B_p(\alpha,\delta)}\delta+\sum_{\alpha'\in\Omega_{\alpha}\setminus \bar B_p(\alpha,\delta)}v_p(\alpha-\alpha')=k\delta+\gamma$$
where $\gamma=\sum\limits_{\alpha'\in\Omega_{\alpha}\setminus \bar B_p(\alpha,\delta)}v_p(\alpha-\alpha')$. Let $O_{\Q_p(\alpha)}$ be
the ring of integers of $\Q_p(\alpha)$, so $O_{\Q_p(\alpha)}$ is a valuation domain. 
Then, as $\alpha$ is unramified,
$O_{\Q_p(\alpha)}$ and $\Z_p$ have the same value group, and thus we have $\gamma\in\Z$. More precisely, $\gamma=v_p(h(\alpha))$, where $h(X)=\prod\limits_{\alpha'\in\Omega_{\alpha}\setminus \bar B_p(\alpha,\delta)}(X-\alpha')$.
\end{proof}


Next, we show that the cardinality of the set of conjugates of $\alpha$ over $\Q_p$ in $\bar B_p(\alpha,\delta)$,  namely, the quantity $\#(\Omega_{\alpha}\cap\bar B_p(\alpha,\delta))$ is 1  when $\alpha$ is unramified over $\Q_p$ and $(\alpha,\delta)$ is a minimal pair.

\begin{Prop}\label{unramified implies k=1}
Let $(\alpha,\delta)\in\oQp\times\Q$ be a minimal pair such that $\alpha$ is unramified over $\Q_p$. Then $\#(\Omega_{\alpha}\cap\bar B_p(\alpha,\delta))=1$. In particular, $\gamma=v_{p,\alpha,\delta}(q_{\alpha})=\delta+\lambda$, for some $\lambda\in\Z$ and we therefore have:
$$e(\gamma\mid\Q_p(\alpha))=e(\delta\mid \Q_p)$$
\end{Prop}

\begin{proof}
Suppose first that $\alpha\in\Q_p$, so that $q_{\alpha}(X)=X-\alpha$ and $\Omega_{\alpha}=\{\alpha\}$. Then $v_{p,\alpha,\delta}(q_{\alpha})=\delta$ and the statement is true in this case.

Suppose now that $\alpha\not\in\Q_p$. We consider the following constants:
\begin{eqnarray*}
\delta_{\Q_p}(\alpha)&=&\sup\{v_p(\alpha-c)\mid c\in\oQp, \ [\Q_p(c):\Q_p]<[\Q_p(\alpha):\Q_p]\}\\
\omega_{\Q_p}(\alpha)&=&\sup\{v_p(\alpha-\alpha')\mid\alpha'\in\Omega_{\alpha}\setminus\{\alpha\}\}
\end{eqnarray*}
which are both well-defined because $\alpha\not\in\Q_p$. It is known that $\delta_{\Q_p}(\alpha)$ is a maximum (see for example \cite[Lemma 1.2]{PerStacked}) and so in general we have the inequality $\delta_{\Q_p}(\alpha)\leq\omega_{\Q_p}(\alpha)$ (which follows by the classical Krasner's lemma). Since $\Q_p(\alpha)/\Q_p$ is unramified, by \cite[Theorem 1.3]{Khanduja} we have the equality $\delta_{\Q_p}(\alpha)=\omega_{\Q_p}(\alpha)$. Moreover, from the definitions it follows that $(\alpha,\delta)$ is a minimal pair if and only if $\delta>\delta_{\Q_p(\alpha)}$.  Now, if $\beta\in\bar B_p(\alpha,\delta)$, we have
$$v_p(\alpha-\beta)\geq\delta>\delta_{\Q_p(\alpha)}=\omega_{\Q_p}(\alpha)$$
so, the only conjugate of $\alpha$ in $\bar B_p(\alpha,\delta)$ is $\alpha$ itself, that is, $\#(\Omega_{\alpha}\cap \bar B_p(\alpha,\delta))=1$.

The last claim now follows by Lemma \ref{conjugates in the ball} and the displayed equation is clear from the fact that $\lambda\in\Z$.
\end{proof}

The next theorem characterizes when a DVR of $\Q(X)$ extending $\Z_{(p)}$ is unramified over $\Z_{(p)}$.

\begin{Thm}\label{unramified DVR over Zp}
Let $W$ be a DVR of $\Q(X)$ extending $\Z_{(p)}$ containing $X$. Then $W$ is unramified over $\Z_{(p)}$ if and only if there exists a minimal pair $(\alpha,\delta)\in\O_p^{\br}\times(\Z_{\geq0}\cup\{\infty\})$ such that $W= \Z_{(p), \alpha, \delta}$ and  $\Q_p(\alpha)/\Q_p$ is unramified.
\end{Thm}
\begin{proof}
By Theorem \ref{unitary DVRs over Z(p)}, there exists $(\alpha,\delta)\in\C_p^{\br}\times(\Q\cup\{\infty\})$ such that $W=\Z_{(p),\alpha,\delta}$. 

Suppose first that $W$ is residually algebraic over $\Z_{(p)}$, so, by the same Theorem $\delta=\infty$ and $\alpha$ is transcendental over $\Q$. In this case, if $\alpha\in\oQp$, then by \cite[Proposition 2.2]{PerTransc} we have $e(W\mid\Z_{(p)})=e(\Q_p(\alpha)\mid\Q_p)$. If instead $\alpha\in\C_p\setminus\oQp$, then  $e(W\mid\Z_{(p)})=e(\Z_{p,\alpha,\delta}\mid\Z_p)=e(\Q_p(\alpha)\mid\Q_p)$, where the first equality follows from  \cite[Proposition 2.24]{PerStacked} and the second equality from \cite[Proposition 2.14]{PerStacked}. In either case, the claim follows. 

Suppose now that $W$ is residually transcendental over $\Z_{(p)}$, so in particular, $\delta\in\Q$ by the aforementioned Theorem,  and we may assume that $\alpha\in\oQp$. By \cite[Theorem 2.1]{APZTheorem}, $e(W\mid\Z_{(p)})=1$ if and only if the following condition holds:
\begin{itemize}
\item[] there exists a minimal pair $(\alpha,\delta)\in\oQp\times\Q$ of definition of $W$ such that $\Q_p(\alpha)/\Q_p$ is unramified and $w(q_{\alpha})\in\Gamma_{\Q_p(\alpha)}$ where $q_{\alpha}$ is the minimal polynomial of $\alpha$ over $\Q_p$.
\end{itemize} 
By Lemma \ref{conjugates in the ball}, we have $w(q_{\alpha})=v_{p,\alpha,\delta}(q_{\alpha})=k\delta+\lambda$, where $k=\#(\Omega_{\alpha}\cap \bar B_p(\alpha,\delta))$ and $\lambda\in \Gamma_{\Q_p(\alpha)}=\Gamma_{\Q_p}=\Z$; in particular, we have $w(q_{\alpha})\in\Gamma_{\Q_p(\alpha)}\Leftrightarrow k\delta\in\Z$.

If $\alpha\in\Q_p$, then $q_{\alpha}(X)=X-\alpha$, so $v_{p,\alpha,\delta}(q_{\alpha})=\delta$ which is therefore an integer.
Suppose that $\alpha\not\in\Q_p$. Since $(\alpha,\delta)$ is  a minimal pair, $\delta>0$. By Proposition \ref{unramified implies k=1} we have $k=1$, so that above displayed condition is equivalent to the condition of the main statement of the theorem.
\end{proof}

\begin{Thm}\label{intersection DVR is UFD}
Let $R$ be a ring between $\Z_{(p)}[X]$ and $\Q[X]$. Then $R$ is a UFD if and only if either $R=\Q[X]$ or $R=W\cap\Q[X]$ where $W$ is a DVR of $\Q(X)$ extending $\Z_{(p)}$ unramified over $\Z_{(p)}$. More precisely, the last condition holds if and only if  there exists $(\alpha,\delta)\in\O_p^{\br}\times(\Z_{\geq0}\cup\{\infty\})$ such that $\alpha$ is unramified over $\Q_p$ and in this case $R=\IntQ(\bar B_p(\alpha,\delta),\O_p)$.    
    %as in \eqref{representation R2} and assume that the representation is irredundant. Then $R$ is a UFD if and only if either $m=0$ (in which case $R=\Q[X]$) or $m=1$, $\alpha_1$ is unramified over $\Q_p$ and $\delta\in\Z\cup\{\infty\}$.
\end{Thm}

\begin{proof}
A UFD is just a Krull domain with trivial divisor class group \cite[Proposition 6.1]{Fossum}. So, if $R$ is a UFD, then,  by Theorem \ref{thm4.23} we have $R=\bigcap\limits_{i=1}^m\Z_{(p),\alpha_i,\delta_i}\cap\Q[X]$, for some $(\alpha_i,\delta_i)\in\O_p^{\br}\times(\Q_{\geq0}\cup\{\infty\})$, for $i=1,\ldots,n$. By items (4) and (5) of Theorem \ref{characterization local Krull domain}, either $R=\Q[X]$ or $R=\Z_{(p),\alpha,\delta}\cap\Q[X]$ and $\Z_{(p),\alpha,\delta}$ is unramified over $\Z_{(p),\alpha,\delta}$. Conversely, if $R$ has the stated form,  then by Theorem \ref{characterization local Krull domain}, $R$ is a Krull domain with trivial class group, hence, a Krull domain.

The proof of the last claim follows from Theorem \ref{unramified DVR over Zp}.
\end{proof}


%{\rc
%\begin{Thm}\label{UFD local case}
%    Let $R$ be a  UFD between $\Z_{(p)}[X]$ and $\Q[X]$. Then $R$ is a UFD if and only if either $R=\Q[X]$ or $R=W\cap\Q[X]$, where $W$ is a DVR of $\Q(X)$ extending $\Z_{(p)}$ unramified over $\Z_{(p)}$.
%\end{Thm}
%\begin{proof}
%The proof follows from Theorems \ref{thm4.23} and \ref{intersection DVR is UFD}.
%\end{proof}


\begin{Ex}
 Let $R$ be a UFD between $\Z_{(p)}[X]$ and $\Q[X]$. We remark that $R$ is equal to $\Q[X]$ precisely when $p$ is a unit of $R$, otherwise $R=W\cap \Q[X]$ for some DVR $W$ of $\Q(X)$ extending $\Z_{(p)}$, such that $W$ is unramified over $\Z_{(p)}$. Note that $R=\Z_{(p)}[X]$ if and only if $W=\Z_{(p),0,0}$. 
\end{Ex}




\subsection{Global case}\label{section global case}

In this section, we globalize the results of Section \ref{section local case}. Namely, we consider a subset $\Lambda \subseteq \PP$, $W_{p, j}$ a DVR
of $\mathbb{Q}(X)$ with $X \in W_{p,j}$ and $W_{p,j} \cap \mathbb{Q} = \mathbb{Z}_{(p)}$,
and $R= \bigcap\limits_{p \in \Lambda}(\bigcap\limits_{j=1}^{m_p}(W_{p,j} \cap \mathbb{Q}[X]))$. 

\subsubsection{When a polynomial Krull domain is Dedekind}


Let  $R$ be a polynomial Krull domain. The next proposition shows that if all the unitary DVRs of the defining family of $R$ are residually algebraic extensions of $\Z_{(p)}$, for $p\in\PP$, then $R$ is a Dedekind domain. We first recall the following definition given in \cite{PerDedekind,PerStacked}.

\begin{Def}
Let $\uE=\prod\limits_{p\in\PP}E_p\subseteq\O=\prod_{p\in\PP}\O_p$ be a subset.
We say that $\uE$ is \emph{polynomially factorizable} if, for each $g\in\Z[X]$ and $\alpha=(\alpha_p)\in\uE$, there exist $n,d\in\Z$, $n,d\geq 1$ such that $\frac{g(\alpha)^n}{d}$ is a unit of $\ohZ$, that is, $v_p(\frac{g(\alpha_p)^n}{d})=0$ for all $p\in\PP$.
\end{Def}

If every $\alpha_p\in E_p$ is transcendental over $\Q$ for every $p\in\PP$, \cite[Lemma 2.12]{PerDedekind} shows that a set $\uE$ as in the above definition is polynomially factorizable if and only if, for each $g\in\Z[X]$, the following set is finite:
$$\PP_{g,\uE}=\{p\in\PP\mid\exists \alpha_p\in E_p \text{ such that }v_p(g(\alpha_p))>0\}.$$
% FALSE!!! $v_p(g())$ dovrebbe essere costante su tutto $E_p$... Actually, if we suppose that each $E_p$ is a \textbf{finite} set, then, for each $g\in\Z[X]$ we have the following equivalence (stronger than \cite[Lemma 2.12]{PerDedekind}):
%$$\PP_{g,\uE}\text{ is finite}\Leftrightarrow \exists n,d\geq 1 \text{ such that }\frac{g(\alpha)^n}{d}\in\O^*,\forall \alpha\in\uE$$
%Hence, if $E_p$ is finite for each $p\in\PP$, then we have
%$$\uE \text{ is polynomially factorizable }\Leftrightarrow\forall g\in\Z[X],\exists n,d\geq1 \text{ such that } \frac{g(\alpha)^n}{d}\in\O^*,\forall\alpha\in\uE$$
%that is, $g(X)^n/d$ is a unit-valued polynomial on $\uE$.
Clearly, we have
$$\PP_{g,\uE}=\{p\in\PP\mid\exists \alpha_p\in E_p \text{ such that }g\in\Z[X]\cap M_{(p),\alpha_p}\},$$
where $M_{(p),\alpha_p}=M_{(p),\alpha_p,\infty}$ is the maximal ideal of $\Z_{(p),\alpha_p}=\Z_{(p),\alpha_p,\infty}$.




\begin{Prop}\label{Krull Dedekind}
Let $R$ be a Krull domain between $\Z[X]$ and $\Q[X]$. Then the following statements 
are equivalent.
\begin{enumerate}
\item[{\em (1)}] $R$ is a Dedekind domain.
\item[{\em (2)}] Each unitary DVR $W$ of the defining family of $R$ is a residually algebraic extension of $W\cap\Q$. 
\item[{\em (3)}] $R = \IntQ(\uE,\O)$, for some subset $\uE=\prod\limits_{p\in\PP}E_p\subset\O=\prod\limits_{p\in\PP}\O_p$ such that $E_p\subset\O_p^{\br}$ is a finite set of transcendental elements over $\Q$ for each $p\in\PP$. 
%\item[{\em (3)}] $R$ is a Dedekind domain equal to $\IntQ(\uE,\O)$, for
%some subset %$\uE=\prod\limits_{p\in\PP}E_p\subset\O=\prod\limits_{p\in\PP}\O_p$ 
%such that $E_p\subset\O_p^{\br}$ is a finite set of transcendental 
%elements over $\Q$ for each $p\in\PP$ and $\uE$ is polynomially 
%factorizable.
\end{enumerate}
 In this case, $\uE$ is polynomially factorizable.
\end{Prop}


%\color{red} note that the necessary condition  of the statement ("$R$ is Dedekind") is also sufficient 
%(in the sense that it implies that "each unitary DVR %$W$ of the defining family of $R$ is a residually algebraic  
%extension of $W\cap\Q$."). This follows by Theorem \ref{characterization local Krull domain}, (3).
%\normalcolor

\begin{proof}
%\color{red} improve notation? for each $p\in\PP$ there might be multiple $\alpha_p\in\O_p^{\br}$ \normalcolor\\
(1) $\Rightarrow$ (2) This is an immediate consequence of Theorem \ref{characterization local Krull domain}, (3).

(2) $\Rightarrow$ (3) Let $\{W_i\}_{i\in I}$ be the proper subset of the defining family of DVRs of $R$ which are extensions of $\Z_{(p)}$ for some $p\in\PP$. By Theorem \ref{DVR ra}, if $W_i$ is a residually algebraic extension of $\Z_{(p)}$, then there exists $\alpha_p\in\O_p^{\br}$, transcendental over $\Q$, such that $W_i=\Z_{(p),\alpha_p}=\Z_{(p),\alpha_p,\infty}$. For each $p\in\PP$, we denote by $E_p$ the subset formed by these $\alpha_p\in\O_p^{\br}$. Then   the defining family of DVRs of $R$ is equal to:
$$\{\Z_{(p),\alpha_p}\mid p\in\PP,\alpha_p\in E_p\}\cup\{\Q[X]_{(g)}\mid g\in\Q[X], \text{ irreducible}\}$$

In particular, we have
$$R=\bigcap_{i\in I}W_i\cap\Q[X]=\bigcap_{p\in\PP}\bigcap_{\alpha_p\in E_p}\Z_{(p),\alpha_p}\cap\Q[X]=\bigcap_{p\in\PP}\IntQ(E_p,\O_p)$$
where the last equality follows from Remark \ref{balls and intval rings}. Note that we also have
$$\bigcap_{p\in\PP}\IntQ(E_p,\O_p)=\IntQ(\uE,\O)=\{f\in\Q[X]\mid f(\alpha)\in\O,\forall\alpha\in \uE\}$$
where $\uE=\prod\limits_p E_p$, $\O=\prod\limits_p \O_p$ and $f(\alpha)=(f(\alpha_p))_p$ if $\alpha=(\alpha_p)_p\in\uE$.
 Finally, let $p\in\PP$. Then, since $p\in R$, $p$ is contained in finitely many centers of the defining family of DVRs of $R$. This implies that $E_p$ is finite.

%in order to show that $R$ is a Dedekind domain, by \cite[Theorem 3.3]
%{PerStacked} (see also \cite[Theorem 2.15]{PerDedekind}) it is enough 
%to show that $E_p$ is finite for each $p\in\PP$ and $\uE$ is 
%polynomially factorizable.

%\color{red} Note that this last part of the proof follows also by Theorem \ref{Krull domains finite character} 
%below, since we know by assumption that $R$ is Krull. \normalcolor

(3) $\Rightarrow$ (1) Let $M$ be a maximal ideal of $R$. 
If $M \cap \Z = (0)$, then ht$M=1$, so assume that $M \cap \Z = p\Z$
for some $p \in \PP$. Note that
$R = \IntQ(\uE,\O) = \bigcap\limits_{p\in\PP}\bigcap\limits_{\alpha_p\in E_p}\Z_{(p),\alpha_p}\cap\Q[X]$, so if $S = \Z \setminus p\Z$, then
$R_S = \bigcap\limits_{\alpha_p\in E_p}\Z_{(p),\alpha_p}\cap\Q[X]$
by Lemma \ref{lemma43} and $\bigcap\limits_{\alpha_p\in E_p}\Z_{(p),\alpha_p}\cap\Q[X]$ is a Dedekind domain by Theorem \ref{thm4.23}. Hence, ht$M =$ ht$(MR_S) =1$.
Thus, $R$ is a one-dimensional Krull domain, therefore $R$ is a Dedekind domain.


% (see also the arguments of the proof of \cite[Proposition  2.10]{PerDedekind}).
Now for the proof of ``In this case", let $f\in\Z[X]\subset R$.  Since $R$ is a Krull domain, $f$ is contained in only finitely many centers of the valuation domains of the defining family for $R$ (clearly,  $f$ is contained in finitely many centers of the non-unitary valuation domains $\Q[X]_{(g)}$, which are the non-unitary prime ideals of $R$ ). This precisely means that the set $\PP_{g,\uE}=\{p\in\PP\mid \exists\alpha_p\in E_p, v_p(f(\alpha_p))>0\}$ is finite.  Hence, $\uE$ is polynomially factorizable by  \cite[Lemma 2.12]{PerDedekind}.
\end{proof}

In particular, a polynomial Krull domain which is not Dedekind must have an essential valuation
overring which is a residually transcendental extension of $\Z_{(p)}$, for some $p\in\PP$.

%{\rc [[I had to restate the statement since if $\mathcal W_{\rt}=\emptyset$ then $R_2=\Q[X]$ which is Dedekind!]]}

% between $\Z[X]$ and $\Q[X]$
\begin{Prop}\label{Krull intersection}
Let $R$ be a polynomial Krull domain. Then
$R=R_1\cap R_2$, where $R_1  $ is a Dedekind domain such that $\Z[X]\subset R_1\subseteq\Q[X]$, and $R_2$ is either
a non-Pr\"ufer Krull domain such that $\Z[X]\subseteq R_2\subset\Q[X]$ 
and $Q$ is not maximal for all $Q \in X^1(R_2)$ with $Q \cap \Z \neq (0)$ 
or $R_2=\Q[X]$.
\end{Prop}

\begin{proof}
Let $\mathcal W(R)=\{W_i\}_{i\in I}$ be the subset of the defining family of DVRs for $R$ which are unitary valuation domains.
Note that $\mathcal W(R)=\emptyset$ is equivalent to saying that $R=\Q[X]$; in this case we have $R_1=R_2=\Q[X]$.
Suppose now that $\mathcal W(R)\not=\emptyset$.
We partition $\mathcal W(R)$ into the subset $\mathcal{W}_{\text{ra}}(R)$ of residually algebraic extensions of $\Z_{(p)}$ for some $p\in\PP$
and the subset $\mathcal{W}_{\text{rt}}(R)$ of residually transcendental extensions of $\Z_{(p)}$ for some $p\in\PP$.

We set:
\begin{align*}
R_1=\bigcap_{W_i\in \mathcal{W}_{\text{ra}}(R)}W_i\cap\Q[X] \ \ \text{ and } \ \
R_2=\bigcap_{W_i\in \mathcal{W}_{\text{rt}}(R)}W_i\cap\Q[X].
\end{align*}
Then $R=R_1\cap R_2$. Clearly, $R_1=\Q[X]$ if and only if $\mathcal{W}_{\text{ra}}(R)=\emptyset$ and similarly $R_2=\Q[X]$ if and only if $\mathcal{W}_{\text{rt}}(R)=\emptyset$.
Note that since $R_1$ and $R_2$ are both subintersections of the Krull domain $R$, they both are Krull domains.
Moreover, by Proposition \ref{Krull Dedekind},
$R_1$ is Dedekind, while if $\mathcal{W}_{\text{rt}}(R)\not=\emptyset$, then $R_2$ is not Pr\"ufer,
since every $W_i\in \mathcal{W}_{\text{rt}}(R)$ is residually transcendental over $W_i\cap\Q$ (see \cite[Theorem 3.14]{PerPrufer}).
%In particular, $R_2$ is not Dedekind.

Now, suppose that $R_2\subset\Q[X]$. Let $Q \in X^1(R_2)$ be such that $Q \cap \Z\neq (0)$,
so $Q \cap \Q = p\Z_{(p)}$ for some $p \in \PP$.
Then ${(R_2)}_{Q} = W_k$ for some $W_k \in \mathcal{W}_{\text{rt}}(R)$ by \cite[Corollary 2.3(i)]{Chang}.
It is clear that if $S$ is the multiplicative subset of $\Z$ generated by $\PP - \{p\}$, then
${(R_2)}_S$ is a Krull domain between $\Z_{(p)}[X]$ and $\Q[X]$.
Note that $(R_2)_S = \bigcap\limits_{W_i\in \mathcal{W}_{\text{rt}}(R)}(W_i)_S\cap\Q[X]$
and ${(W_k)}_S = W_k$, so $Q_S$ is not a maximal ideal of ${(R_2)}_S$ by Theorem \ref{characterization local Krull domain} (3).
Thus $Q$ is not a maximal ideal of $R_2$.
\end{proof}

In Proposition \ref{Krull intersection},
$R_1=\Q[X]$ if $R$ has no valuation overrings which are residually algebraic extensions of $\Z_{(p)}$ for all $p\in\PP$,
and $R_2= \Q[X]$ if $R$ is a Dedekind domain (or, equivalently, $R$ has no valuation overrings
which are residually transcendental  extensions of $\Z_{(p)}$ for all $p\in\PP$).
Therefore, the objective is now to  characterize when, for a given family $\{W_i\}_{i\in I}$ of residually transcendental extensions of $\Z_{(p)}$, $p\in\PP$, the ring $\bigcap\limits_{i\in I}W_i\cap\Q[X]$ has finite character, and thus is a Krull domain.

%\color{red}
%I have a question regarding the statement of the above Proposition. Why don't we use the following one, instead?
%Let $R$ be a Krull domain between $\Z[X]$ and $\Q[X]$. Then
%$R=R_1\cap R_2$, where $R_1  $ is a Dedekind domain such that $\Z[X]\subset R_1\subseteq\Q[X]$, and $R_2$ is either
%a non-Pr\"ufer Krull domain such that $\Z[X]\subseteq R_2\subset\Q[X]$  or $R_2=\Q[X]$.\\

%I don't see the need to specify that "and $Q$ is not maximal for all $Q \in X^1(R_2)$ with $Q \cap \Z \neq (0)$" for $R_2$ (but of course, I agree with the truth of this statement). 
%{\color{blue} {\rc I have thought about this problem and I changed my mind because we need this part.}
%You are right, so I changed the statement. In fact, in Proposition \ref{Krull intersection}, we have that
%(i) a Krull domain $R$ is not Dedekind $\Leftrightarrow$ (ii) $R$ is non-Pr\"ufer, $\Leftrightarrow$ 
%(iii) $Q$ is not maximal for some $Q \in X^1(R_2)$ with $Q \cap \Z \neq (0)$.
%{\rc I have thought about this problem and I changed my mind. Note that if we don't need this fact, then we don't need this proposition,
%because if $R$ is not Dedekind, then $R_2$ is just $R$ and $R_1 = \mathbb{Q}[X]$.}} 
%\normalcolor





%\vskip1cm
\subsubsection{Centers of monomial valuations on $\Z[X]$}\label{centers}
%Corollary \ref{coro38}
Let $p\in\PP$ be fixed. Given $(\alpha,\delta)\in\C_p \times(\Q\cup\{\infty\})$, we denote the maximal ideal of $\Z_{(p),\alpha,\delta}$ by $M_{(p),\alpha,\delta}$. By Theorem \ref{unitary DVRs over Z(p)}, $\Z_{(p),\alpha,\delta}$ is an overring of $\Z_{(p)}[X]$ if and only if $(\alpha,\delta)\in\O_p\times(\Q_{\geq0}\cup\{\infty\})$. In this section, we restrict to this case and describe the center of $\Z_{(p),\alpha,\delta}$ on $\Z_{(p)}[X]$, namely, $M_{(p),\alpha,\delta}\cap (\Z_{(p)}[X])$.

Let $\mathbb M_p$ be the maximal ideal of $\O_p$. We set $\Int(\bar B_p(\alpha,\delta),\mathbb M_p)=\{f\in\C_p[X]\mid f(\bar B(\alpha,\delta))\subseteq\mathbb M_p\}$. It is easy to see that $\Int(\bar B_p(\alpha,\delta),\mathbb M_p)$ is an ideal of $\Int(\bar B_p(\alpha,\delta),\O_p)$. The following lemma shows that this ideal is precisely the center of $\O_{p,\alpha,\delta}$ on $\O_{p,\alpha,\delta}\cap\C_p[X]=\Int(\bar B_p(\alpha,\delta),\O_p)$ (see \eqref{restriction valuation domain on polynomial ring}).

\begin{Lem}\label{IntBallmaximal}
Let $p\in\PP$, $(\alpha,\delta)\in\C_p\times(\Q\cup\{\infty\})$ and $\mathbb M_{p,\alpha,\delta}$ be 
the maximal ideal of $\O_{p,\alpha,\delta}$. Then,  we have
$$\mathbb M_{p,\alpha,\delta}\cap(\O_{p,\alpha,\delta}\cap\C_p[X])=\Int(\bar B_p(\alpha,\delta),\mathbb M_p).$$
\end{Lem}
\begin{proof}
The statement clearly holds if $\delta=\infty$. Suppose then that $\delta<\infty$.

$(\subseteq)$ Let $f\in \mathbb{M}_{p,\alpha,\delta}\cap(\O_{p,\alpha,\delta}\cap\C_p[X])$. Then, if $f(X)=\sum\limits_k a_k(X-\alpha)^k$, we have $v_p(a_k)+k\delta>0$ for each $k$. In particular, if $\alpha'\in\C_p$ is such that $v_p(\alpha-\alpha')\geq\delta$, we have $f(\alpha')=\sum\limits_k a_k(\alpha'-\alpha)^k$, so that $v_p(f(\alpha'))\geq\min_k\{v_p(a_k)+k\delta\}>0$. Hence, $f(\bar B(\alpha,\delta))\subseteq\mathbb M_p$.

$(\supseteq)$ Let $f\in \Int(\bar B_p(\alpha,\delta),\mathbb M_p)\subset\Int(\bar B_p(\alpha,\delta),\O_p)$. Let $E=\{\alpha_n\}_{n\in\N}\subset\bar B_p(\alpha,\delta)$ be a pseudo-stationary sequence such that $\delta_E=\delta$ and let $v_{p,E}$ be the associated valuation (see \cite{PS2}). By \cite[Proposition 3.7]{PS2}, $v_{p,\alpha,\delta}=v_{p,E}$ and by \cite[Theorem 3.4]{PS2}, $f(X)$ is in the maximal ideal of $\O_{p,E}=\O_{p,\alpha,\delta}$, because by assumption $v_p(f(\alpha_n))>0$ for every $n\in\N$; thus, $f\in \mathbb{M}_{p,\alpha,\delta}$.
\end{proof}

\vskip0.2cm
%Given $\alpha\in\O_p$, we denote by $\overline{\alpha}$ the residue of $\alpha$ modulo $\mathbb M_p$ (which is the same as the residue of $\alpha$ modulo the maximal ideal $M_{O_{\Q_p(\alpha)}}$ of $O_{\Q_p(\alpha)}$); note that $\overline{\alpha}\in\overline{\F_p}$. We denote by $g_{\overline{\alpha}}(X)$ the minimal polynomial of $\overline{\alpha}$ over $\F_p$.  Similarly, given $g\in\Z[X]$, we denote by $\overline{g}\in(\Z/p\Z)[X]$ the reduction modulo $p$ of the polynomial $g(X)$. Finally, for each $\alpha\in\O_p$, we fix an arbitrary monic $g_{\alpha}\in\Z_{(p)}[X]$ such that $\overline{g_{\alpha}}(X)=g_{\overline{\alpha}}(X)$. Note that $v(g_{\alpha}(\alpha))>0$, that is, $g_{\alpha}\in M_{(p),\alpha,\infty}$.

For each $p \in \PP$, let $\mathbb{F}_p$ be the finite field of $p$ elements, i.e., 
$\mathbb{F}_p = \mathbb{Z}/p\mathbb{Z}$, and $\overline{f}$ be the image of $f \in \mathbb{Z}[X]$ in $\mathbb{F}_p[X]$.
Let Spec$(\Z[X])$ (resp., Max$(\Z[X])$) be the set of prime ideals (resp., maximal ideals) of $\Z[X]$.
Note that dim$(\Z[X]) =2$ and each maximal ideal of $\Z[X]$ has height two,
so
\begin{center}
Spec$(\Z[X]) = \{0\} \cup X^1(\Z[X]) \cup$ Max$(\Z[X])$ and $X^1(\Z[X]) \cap$ Max$(\Z[X]) = \emptyset$.
\end{center}
It is well known and easy to see that
\begin{enumerate}
\item[(a)] $X^1(\Z[X]) = \{p\Z[X] \mid p \in \PP\} \cup \{f\Z[X] \mid f\in\Z[X]$ and $f$ is irreducible in $\Z[X]\}$
\item[(b)] Max$(\Z[X]) = \{(p, f) \mid p \in \PP, f\in\Z[X]$ and $\overline{f}\F_p[X]$ is a prime ideal$\}$.
\end{enumerate}
In particular, for each $p \in \PP$,
%\begin{multline}
\begin{eqnarray*}
\text{Spec}(\Z_{(p)}[X]) &=& \{0\} \cup  \{p\Z_{(p)}[X]\}\\
                         &\cup& \{f\Z_{(p)}[X] \mid f\in\Z_{(p)}[X] \text{ and }f \text{ is irrducible in  }\Z_{(p)}[X]\} \\
&\cup&\{(p, f) \mid \overline{f}\F_p[X] \text{ is a prime ideal }\}.
\end{eqnarray*}
%\end{multline}


%$$\text{Spec}(\Z_{(p)}[X]) = \{0\} \cup  \{p\Z_{(p)}[X]\} \cup \{f\Z_{(p)}[X] \mid f\in\Z_{(p)}[X], c(f)\Z_{(p)} = \Z_{(p)} \text{ and }f \text{ irr. over }\Q\}\cup \{(p, f) \mid \overline{f}\F_p[X] \text{ is a prime ideal }\}.$$
Recall that if $\alpha \in \mathbb{O}_p$, then $\overline{\alpha}$ is the image of $\alpha$
modulo $\mathbb{M}_p$ and $g_{\alpha}(X) \in \mathbb{Z}_{(p)}[X]$ denotes a monic polynomial
such that $\overline{g_{\alpha}(X)}$ is the minimal polynomial over $\overline{\alpha}$ over $\mathbb{F}_p$.


\begin{Prop}\label{contraction maximal ideal}
Let $p\in\PP$ and $(\alpha,\delta)\in\O_p\times(\Q_{\geq0}\cup\{\infty\})$. Then
$$\Z_{(p)}[X]\cap M_{(p),\alpha,\delta}=\left\{\begin{array}{lc}
(p),&\text{ if } \delta=0\\
(p,g_{\alpha}(X)),&\text{ if } \delta>0
\end{array}\right.$$
\end{Prop}
\begin{proof} Suppose first that $\delta=0$. Then $\Z_{(p),\alpha,0}=\Z_{(p),0,0}= \Z_{(p)}[X]_{p\Z_{(p)}[X]}$ (i.e., the trivial Gaussian extension).
In particular, $\Z_{(p)}[X]=\Z_{(p),0,0}\cap\Q[X]$, so, the center of $M_{(p),\alpha,0}=M_{(p),0,0}$ on $\Z_{(p)}[X]$ is equal to $p\Z_{(p)}[X]$.

Suppose now that $\delta=\infty$. Since $(p,g_{\alpha}(X))$  is a maximal ideal of $\Z_{(p)}[X]$ and $g_{\alpha}\in M_{(p),\alpha,\infty}$,  $(p,g_{\alpha}(X))$ is the center of $\Z_{(p),\alpha,\infty}$ on $\Z_{(p)}[X]$.

Finally, suppose that $0<\delta<\infty$. By Lemma \ref{IntBallmaximal}, 
$$M_{(p),\alpha,\delta}\cap(\Z_{(p),\alpha,\delta}\cap\Q[X])=\IntQ(\bar B_p(\alpha,\delta),\mathbb M_p),$$ 
so $\Z_{(p)}[X]\cap M_{(p),\alpha,\delta}=\{f\in\Z_{(p)}[X]\mid f(\bar B_p(\alpha,\delta))\subseteq\mathbb M_p\}$. 
Hence, in order to show that this last intersection is equal to $(p,g_{\alpha}(X))$, it is sufficient to show that $g_{\alpha}$ maps $\bar B_p(\alpha,\delta)$ into $\mathbb M_p$. If $\alpha'\in\C_p$ is such that $v(\alpha'-\alpha)\geq\delta>0$, then $g_{\alpha}(\alpha')\equiv g_{\alpha}(\alpha)\equiv 0\pmod{\mathbb M_p}$, which is what we wanted to prove.
\end{proof}

\begin{Rem}\label{collapse centers DVRs}
Clearly, if $\bar B_p(\alpha,\delta)\subset\O_p$ with $\delta>0$, then for all $\alpha'\in \bar B_p(\alpha,\delta)$, 
we have $\overline{\alpha}=\overline{\alpha'}$, so $\overline{g_{\alpha}}=\overline{g_{\alpha'}}$. 
Hence, for $\delta>0$, the center of $\Z_{(p),\alpha,\delta}$ on $\Z_{(p)}[X]$ does not depend on the particular center of the ball $\bar B(\alpha,\delta)$.
Most importantly, by Lemma \ref{contraction maximal ideal}, given $\alpha\in\O_p$, 
the center of $\Z_{(p),\alpha,\delta}$ on $\Z_{(p)}[X]$ is the same maximal ideal for all $\delta\in\Q_{>0}\cup\{\infty\}$.  

Putting these two remarks together, for each $\delta\in\Q,\delta>0$, we have that
$$M_{(p),\alpha,\delta}\cap\Z[X]=M_{(p),\alpha',\infty}\cap\Z[X],\;\;\forall\alpha'\in\bar B_p(\alpha,\delta).$$
In particular, for $\delta>0$ and $g\in\Z[X]$, we have that $g\in M_{(p),\alpha,\delta}$ 
if and only if $g\in M_{(p),\alpha',\infty}$ for some  (hence, all) $\alpha'\in \bar B_p(\alpha,\delta)$.
\end{Rem}


\begin{Cor}\label{equality centers on Z[X]}
Let $p\in\PP$ and
%\begin{itemize}
%\item[i)] Let $\alpha_{i}\in\O_p$ for $i=1,2$. Then
%$$\Z_{(p)}[X]\cap M_{(p),\alpha_{1}}=\Z_{(p)}[X]\cap M_{(p),\alpha_{2}}\Leftrightarrow \overline{\alpha_{1}},\overline{\alpha_{2}}\text{ are conjugate over }\F_p.$$
%\item[ii)]
$(\alpha_i,\delta_i)\in\O_p\times(\Q_ {\geq0}\cup\{\infty\})$, for $i=1,2$.\\ Then $\Z_{(p)}[X]\cap M_{(p),\alpha_1,\delta_1}=\Z_{(p)}[X]\cap M_{(p),\alpha_2,\delta_2}$ if and only if one of the following conditions holds:
\begin{enumerate}
\item[{\em (a)}]  $\delta_1=\delta_2=0$;
\item[{\em (b)}] $\delta_i\not=0$, for $i=1,2$,  and $\overline{\alpha_1},\overline{\alpha_2}$  are  conjugate over $\F_p$. 
Equivalently, $v_p(\sigma(\alpha_1)-\alpha_2)>0$ for some $\sigma\in G_p=\Gal(\oQp/\Q_p)$.
\end{enumerate}
%\item[ii)OLD!] Let $(\alpha_i,\delta_i)\in\overline{\Z_{(p)}}\times\Q$, $\delta_i\geq0$ for $i=1,2$; let $F$ be a finite extension of $\Q$ containing $\alpha_i$, for $i=1,2$, and such that $\delta_i\in\Gamma_{U_F}$, where  $U_F=U\cap F=O_{F,P}$ and $P\subset O_F$ is a prime ideal above $p$. Then $\Z_{(p)}[X]\cap M_{(p),\alpha_1,\delta_1}=\Z_{(p)}[X]\cap M_{(p),\alpha_2,\delta_2}$ if and only if either one of the following conditions holds:
%\begin{enumerate}
%\item  $\delta_1=\delta_2=0$;
%\item  $\delta_i\not=0$, for $i=1,2$,  and $\alpha_1\pmod{PO_{F,P}},\alpha_2\pmod{PO_{F,P}}$  are  conjugate over $\F_p$.
%\end{enumerate}

%\end{itemize}
\end{Cor}
\begin{proof}
$(\Rightarrow)$ Suppose first that $\Z_{(p),\alpha_1,\delta_1},\Z_{(p),\alpha_1,\delta_1}$ have the same center on $\Z_{(p)}[X]$.
 By Proposition \ref{contraction maximal ideal}, we have either $\delta_1=\delta_2=0$ or $\delta_i\not=0$ for $i=1,2$. In the last case, we have $\Z_{(p)}[X]\cap M_{(p),\alpha_i,\delta_i}=(p,g_{\alpha_i}(X))$, where $g_{\alpha_i}\in\Z_{(p)}[X]$ is a monic polynomial such that $\overline{g_{\alpha_i}}$ is the minimal polynomial of $\overline{\alpha_i}$ over $\F_p$, for $i=1,2$.  Then, these prime ideals are the same if and only if $\overline{g_{\alpha_1}}=\overline{g_{\alpha_2}}$, that is, $\overline{\alpha_1},\overline{\alpha_2}$ are conjugate over $\F_p$. Note that this last condition amounts 
 to saying that there exists some $\overline{\sigma}\in \Gal(\overline{\F_p}/\F_p)$ such that $\overline{\sigma}(\overline{\alpha_1})=\overline{\alpha_2}$. Since $\overline{\sigma}(\overline{\alpha_1})=\overline{\sigma(\alpha_1)}$ for some $\sigma\in G_p$  \cite[Chapt. V, §2.2, Proposition 6(ii)]{Bourb}), it follows that $v_p(\sigma(\alpha_1)-\alpha_2)>0$, as claimed.

$(\Leftarrow)$ If $\delta_1=\delta_2=0$, then $\Z_{(p)}[X]\cap M_{(p),\alpha_i,\delta_i}=(p)$ for $i=1,2$. 
If instead $\delta_i\not=0$ for $i=1,2$ and $\overline{\alpha_1},\overline{\alpha_2}$ are conjugate over $\F_p$, 
then by the above arguments we have $\Z_{(p)}[X]\cap M_{(p),\alpha_1,\delta_1}=\Z_{(p)}[X]\cap M_{(p),\alpha_2,\delta_2}$.
\end{proof}

In Corollary \ref{equality centers on Z[X]}, if $\delta_i\not=0$ for $i=1,2$, we have 
the following three possibilities: (i) $\alpha_{1},\alpha_{2}$ are conjugate by the action of $G_p$ and $\delta_1=\delta_2$ 
and so $\Z_{(p),\alpha_{1},\delta_1}=\Z_{(p),\alpha_{2},\delta_2}$ (and so in particular they have the same center on $\Z_{(p)}[X]$), 
(ii) $\alpha_{1},\alpha_{2}$ are not conjugate over $\Q_p$ or (iii)
$\delta_1\not=\delta_2$ so that $\Z_{(p),\alpha_{1},\delta_1}\not=\Z_{(p),\alpha_{2},\delta_2}$ but  
$\overline{\alpha_{1}},\overline{\alpha_{2}}$ are conjugate over $\F_p$ so the corresponding valuation domains have the same center on $\Z_{(p)}[X]$. 
For example, if $\alpha_1,\alpha_2\in \O_p$ are such that $v_p(\alpha_1-\alpha_2)>0$ (that is, they belong to the same ball contained in $\O_p$ of strictly positive radius), then $\overline{\alpha_1}=\overline{\alpha_2}$.



\subsubsection{Finite character}

A Polynomial Krull domain $R$ must be of the form $R = \bigcap\limits_{p \in \Lambda}(\bigcap\limits_{j=1}^{m_p}(W_{p,j} \cap \Q[X])),$ for some $\Lambda\subseteq\PP$ and DVRs $W_{p,j}$, for each $p\in\Lambda$ and $j\in\{1,\ldots,m_p\}$. But, in general, a ring of this form need not be a polynomial Krull domain (see, for example, \cite[Theorem 3.5]{Chang}
or Example \ref{ex430}). 

\begin{Ex} \label{ex430}
For each prime $p \in \PP$, let $\delta_p\in\Q$, $\delta_p>0$, and let
$$R = \bigcap_{p\in\PP}\Z_{(p),0,\delta_p}\cap\Q[X].$$
Then $R$ is not Krull, because $X$ is contained in each center of
the residually transcendental extensions $\Z_{(p),0,\delta_p}$ of $\Z_{(p)}$, for $p\in\PP$.
%Note that ht$(M_{(p),0,\delta_p}\cap R)=1$ by Corollary \ref{cor10} and $pR = (M_{(p),0,\delta_p})^{n_p} \cap R$
%for all $p \in \PP$, where $n_p$ is the smallest positive integer with $n_p\delta_p$ an integer, so $R$ is a $t$-almost Dedekind domain with Cl$(R) = \bigoplus\limits_{p \in \PP}\Z/n_p\Z$ by \cite[Theorem 2.4]{Chang}.
%In particular, $R$ is an almost GCD domain,  i.e., given a pair $a, b \in R$, there is a positive integer $n = n(a,b)$ such that $a^nR \cap b^nR$ is principal \cite{z85},and $R$ is a GCD domain if and only if $n_p=1$, i.e., $p\Z_{(p),0,\delta_p} = M_{(p),0,\delta_p}$ for all $p \in \PP$. It is clear that a GCD domain is an almost GCD domain.
\end{Ex}

\noindent
In this section, we determine when such an intersection  is a Krull domain.
%$R = \bigcap_{p \in \Lambda}(\bigcap_{j=1}^{m_p}(W_{p,j} \cap \Q[X]))$ 

%Assume that for each $p\in\PP$, every $\alpha_p\in E_p$ is transcendental over $\Q$;  

We consider the case of a subset $E=\prod\limits_{p\in\PP}E_p$ of $\O$ such that, for each $p\in\PP$, $E_p$ is a finite union of balls in $\O_p^{\br}$:
$$E_p=\bigcup_{j=1}^{m_p}\bar B_p(\alpha_{p,j},\delta_{p,j}),\;\;(\alpha_{p,j},\delta_{p,j})\in\O_p^{\br}\times(\Q_{\geq0}\cup\{\infty\}) \;\text{ for } j=1,\ldots,m_p,$$
with the assumption that $\alpha_{p,j}$ is transcendental over $\Q$ for each $p\in\PP$ and $j\in\{1,\ldots,m_p\}$. If $\delta_{p,j}\not=\infty$ this can always be assumed by a density argument and if $\delta_{p,j}=\infty$ this condition has to be imposed. For each $p\in\PP$, we consider the set of centers of the balls in $E_p$ of strictly positive radius   given by $C(E_p)=\{\alpha_{p,1},\ldots,\alpha_{p,m_p}\mid \delta_{p,j}>0,\forall j=1,\ldots,m_p\}$ and set
$$C(\uE)=\prod_{p\in\PP}C(E_p)$$
Note that, for each $p\in\PP$, the set of the centers of $\{\Z_{(p),\alpha_{p,j},\delta_{p,j}}\mid j=1,\ldots,m_p\}$ on $\Z[X]$, namely
$$\{\Z[X]\cap M_{(p),\alpha_{p,j},\delta_{p,j}}\mid j=1,\ldots,m_p\}$$
is finite (possibly, by Corollary \ref{equality centers on Z[X]} its cardinality  can be strictly smaller than $m_p$).

We remark that the choice of a center of an ultrametric ball $\bar B_p(\alpha,\delta)\subset \O_p$ is completely arbitrary (every element of the ball is a center of the ball itself). However, for our scopes, by Remark \ref{collapse centers DVRs}
 a possible different choice does not affect the next Theorem \ref{Krull domains finite character}, \textbf{provided} the center of each  ball is transcendental over $\Q$ 
 %({\rc the reason is due to the fact that otherwise \cite[Lemma 2.12]%{PerDedekind} does not hold anymore})
  (because, in this case, \cite[Lemma 2.12]{PerDedekind} can be applied for the proof of Theorem \ref{Krull domains finite character}). This property can always be assumed if $\delta_{p,j}\in\Q_{\geq0}$ and it has to be imposed when $\delta_{p,j}=\infty$.
%{\rc \textbf{Question}: is there an analogous definition for polynomially factorizable for sets like $E=\prod_p E_p$ where each $E_p$ is a finite union of ultrametric balls?}
%\begin{Def}
%Let $\mathcal{E}$ be as above. We say that $\mathcal{E}$ is \emph{polynomially factorizable} if, for each $g\in\Z[X]$ and ...?? there exist $n,d\in\Z,n,d\geq1$ such that $v_{p,\alpha_{p,j},\delta_{p,j}}(\frac{g^n}{d})=0$. ??? INCOMPLETE
%\end{Def}
%\vskip0.5cm
%\begin{Lem}
%Let $\mathcal{E}$ be as above and let $g\in\Z[X]$.  Then there exists $n,d\geq1$ such that $n\inf_{x\in\mathcal{E}_p}v_p(g(x))=v_p(d)$ for each $p\in\PP$ ???? if and only if the following set is finite:
%$$\PP_{g,\mathcal{E}}=\{p\in\PP\mid\exists j\in\{1,\ldots,m_p\} \text{ such that }g\in\Z[X]\cap\ M_{p,\alpha_{p,j},\delta_{p,j}}\}$$
%\end{Lem}
\vskip0.5cm
\begin{Thm}\label{Krull domains finite character}
 Let $\Lambda \subseteq \PP$ be a subset,  $E=\prod\limits_{p\in\Lambda}E_p\subset\O$ where for each $p\in\Lambda$, $E_p=\bigcup\limits_{j=1}^{m_p}\bar B_p(\alpha_{p,j},\delta_{p,j})$ with $(\alpha_{p,j},\delta_{p,j})\in\O_p^{\br}\times(\Q_{\geq0}\cup\{\infty\})$, $\alpha_{p,j}$ transcendental over $\Q$, for each 
$j\in\{1,\ldots,m_p\}$ and let $R=\IntQ(\uE,\O)$, where $\uE=\prod\limits_{p\in\Lambda}E_p$, namely:
$$R=\bigcap_{p\in\Lambda}\bigcap_{j=1}^{m_p}\Z_{(p),\alpha_{p,j},\delta_{p,j}}\cap\Q[X].$$
Assuming that the above representation of $R$ is irredundant, 
%[I guess if we assume that the above family of DVRs $\Z_{(p),\alpha_{p,j},\delta_{p,j}}$ is a defining family as we do in the proof,  than we must assume that the representation is irredundant, right?],
then the following conditions are equivalent:
\begin{itemize}
    \item[{\em (1)}] $R$ is a Krull domain.
  %  \item[ii)] $\mathcal E$ is polynomially factorizable.
    \item[{\em (2)}] For each (nonconstant and irreducible) $g\in\Z[X]$, $g$ is contained in finitely many heigth one unitary prime ideals of $R$.
   \item[{\em (3)}]  For each $g\in\Z[X]$, the following set is finite:
$$\mathbb D_{g,\uE}=\{p\in\Lambda \mid\exists j\in\{1,\ldots,m_p\} \text{ such that }g\in\Z[X]\cap\ M_{(p),\alpha_{p,j},\delta_{p,j}}\}.$$
\item[{\em (4)}] $C(\uE)=\prod\limits_{p\in\Lambda}\{\alpha_{p,1},\ldots,\alpha_{p,m_p}\mid \forall j=1,\ldots,m_p,\delta_{p,j}>0\}$ is polynomially factorizable.
\end{itemize}
 \end{Thm}
We remark that if for some $p\in\Lambda$ we have $\delta_{p,\alpha_j}=0$ for $j=1,\ldots,m_p$, then $m_p=1$ by the assumption that the above  representation of $R$ is irredundant. For such a prime $p$, we have 
  $$R_{(p)}=\Z_{(p),\alpha_p,0}\cap\Q[X]=\IntQ(\O_p)=\Z_{(p)}[X]$$ (see for example \eqref{IntQpB}). 
\begin{proof}
(1) $\Leftrightarrow$ (2) Note that $R$ is a Krull domain if and only if, for every $f\in R$, $f\not=0$, 
$f$ is contained in finitely many centers of the defining family of DVR overrings of $R$, 
namely, $\Z_{(p),\alpha_{p,j},\delta_{p,j}}$, for $p\in\Lambda$ and $j=1,\ldots,m_p$ 
and $\Q[X]_{(q)}$, for $q\in\Q[X]$ irreducible.

 Since every $f\in R$ can be written as $\frac{g}{d}$, for some $g\in\Z[X]$    
 and $d\in\Z$, $d\not=0$, it is clear that $f$ is contained in finitely many centers of the defining family of DVR overrings of $R$ if and only if the same holds for $g$, 
 since the integer $d$ is divisible by finitely many $p\in\Lambda$ and for each $p\in\Lambda$ there are only finitely many DVRs overring of $R$ extending $\Z_{(p)}$. 
 Moreover, $g$ is clearly contained in finitely many nonunitary prime ideals of $R$ (whose height is one by an argument similar 
 to the one of  Corollary \ref{nonunitary not superfluous}), so $R$ is Krull if and only if condition (2) holds (by assumption and  by Corollary \ref{cor10}, the centers on $R$ of the DVRs $\Z_{(p),\alpha_{p,j},\delta_{p,j}}$ are of height one).

(2) $\Leftrightarrow$ (3)
 Since for each $p\in\Lambda$, we only have finitely many unitary DVR overrings of $R$ extending $\Z_{(p)}$ 
 (namely, $\Z_{(p),\alpha_{p,j},\delta_{p,j}}$, for $j=1,\ldots,m_p$) and their centers on $R$ are of height one by Corollary \ref{cor10}, it follows that (2) is equivalent to (3).

(3) $\Leftrightarrow$ (4)
 Let $\PP_0$ be the set of primes in $\Lambda$ for which there exists $j\in\{1,\ldots,m_p\}$ such that $\delta_{p,j}=0$. If $p \in \PP_0$, then $m_p =1$ 
 by Remark \ref{ball containment} and the assumption that the above  representation of $R$ is irredundant, and $M_{(p),\alpha_{p,1},\delta_{p,1}}\cap\Z_{(p)}[X]=(p)$ by Proposition \ref{contraction maximal ideal}. Hence, for each nonzero $g\in\Z[X]$, $\mathbb D_{g,\uE}\cap\PP_0$ is a finite set. By the same aforementioned Proposition, if $p\in\Lambda \setminus\PP_0$, we have $M_{(p),\alpha_{p,j},\delta_{p,j}}\cap\Z_{(p)}[X]=M_{(p),\alpha_{p,j},\infty}\cap\Z_{(p)}[X]$, so $\mathbb D_{g,\uE}\setminus\PP_0=\PP_{g,C(\uE)}$. Therefore, $\mathbb D_{g,\uE}$ is finite if and only if $\PP_{g,C(\uE)}$ is finite. Since each $\alpha_{p,j}$ is transcendental over $\Q$ by assumption,  \cite[Lemma 2.12]{PerDedekind} shows that  $C(\uE)$ is polynomially factorizable if and only if  for every $g\in\Z[X]$, $\PP_{g,C(\uE)}$ is finite. Hence, (3) and (4) are equivalent. %$C(\uE)$ being polynomially factorizable.
 %Let $\PP_{g,C(\uE)}=\{p\in\PP\mid \exists j\in\{1,\ldots,m_p\}\text{ such that }g\in\Z[X]\cap M_{(p),\alpha_{p,j},\infty}\}.$
%Then, by Remark \ref{collapse centers DVRs}, $\mathbb D_{g,\uE} = \PP_{g,C(\uE)}$. 
%Hence, by \cite[Lemma 2.12]{PerDedekind}, the fact that $\PP_{g,C(\uE)}$ is
%finite for each $p\in\PP$ is equivalent to $C(\uE)$ being polynomially factorizable.
\end{proof}

Let $\Lambda$ and $R$ be as in Theorem \ref{Krull domains finite character}. It is helpful to note that $\Lambda = \emptyset$ if and only if $R = \Q[X]$, and, if $\Lambda$ is finite, then $R$ is always a Krull domain. 
The next theorem is the global version of Theorem \ref{thm4.23} and describes the class of polynomial Krull Domains. 
By a result of Heinzer, this is precisely the class of integrally closed Noetherian domains between $\Z[X]$ and $\Q[X]$ (see \cite{Heinzer}).

\begin{Thm}\label{main thm polynomial krull domains}
Let $R$ be a polynomial Krull domain. Then  
\begin{equation}\label{Krull domain as int ring}
R=\bigcap_{p\in\Lambda}\bigcap_{j=1}^{m_p}\Z_{(p),\alpha_{p,j},\delta_{p,j}}\cap\Q[X]=\IntQ(\uE,\O)    
\end{equation}
for some subset $\Lambda\subseteq\PP$, where $(\alpha_{p,j},\delta_{p,j})\in\O_p^{\br}\times(\Q_{\geq0}\cup\{\infty\})$ 
with $\alpha_{p,j}$ transcendental over $\Q$ if $\delta_{p,j}=\infty$ for each $p\in\Lambda$ and $j=1,\ldots,m_p$; 
$\uE=\prod\limits_{p\in\Lambda} E_p$ with  $E_p=\bigcup\limits_{j=1}^{m_p}\bar B_p(\alpha_{p,j},\delta_{p,j})$; and $C(\uE)=\prod\limits_{p\in\Lambda}\{\alpha_{p,1},\ldots,\alpha_{p,m_p} \mid \forall j=1,\ldots,m_p,\delta_{p,j}>0\}$ is polynomially factorizable.


Moreover, if ht$(M_{(p),\alpha_{p,j},\delta_{p,j}} \cap R) =1$ for all $p \in \Lambda$ and
$j=1, \dots, m_p$, then $$\text{Cl}(R) = \bigoplus_{p \in \Lambda}(\Z/d_p\Z \oplus \Z^{m_p-1}),$$ where
$e_{p,j} = e(\Z_{(p),\alpha_{p,j},\delta_{p,j}}|\Z_{(p)})$ for $j=1, \dots , m_p$ and $d_p= gcd(e_{p,1}, \dots , e_{p,m_p})$
for all $p \in \Lambda$.
\end{Thm}
%    In particular, $R$ is a UFD if and only if, for each $p\in\PP$, either $m_p=0$ or $m_p=1$ and $d_p=1$. If this last case occurs for some $p\in\PP$, then this amounts to saying that $\alpha_p$ is unramified over $\Q_p$ and $\delta_p\in\Z\cup\{\infty\}$.
\begin{proof}
As we remarked in \eqref{representation Krull domain}, we have
$$R = \bigcap_{p \in \Lambda}(\bigcap_{j=1}^{m_p}W_{p,j}) \cap \Q[X],$$
for some subset $\Lambda\subseteq\PP$ where $W_{p,j}$ is a DVR of $\Q(X)$ lying over $\Z_{(p)}$ for each $j=1,\ldots,m_p$ and $p\in\PP$.

By  Theorem \ref{unitary DVRs over Z(p)}, for each $p\in\Lambda$, there exist  $(\alpha_{p,j},\delta_{p,j})\in\O_p^{\br}\times(\Q_{\geq0}\cup\{\infty\})$ such that $W_{p,j}=\Z_{(p),\alpha_{p,j},\delta_{p,j}}$, for each $j=1,\ldots,m_p$. In particular, $R$ has the stated form.

Finally, $C(\uE)=\prod\limits_{p\in\Lambda}\{\alpha_{p,1},\ldots,\alpha_{p,m_p} \mid \forall j=1,\ldots,m_p,\delta_{p,j}>0\}$ is polynomially factorizable by Theorem \ref{Krull domains finite character}
and Cl$(R) = \bigoplus\limits_{p \in \Lambda}(\Z/d_p\Z \oplus \Z^{m_p-1})$ by the proof of \cite[Theorem 2.4(4)]{Chang}.
\end{proof}


In the next corollary, we completely describe UFDs between $\Z[X]$ and $\Q[X]$.
\begin{Cor}\label{global UFD}
Let $R$ be a UFD between $\Z[X]$ and $\Q[X]$. Then there exist a subset 
$\Lambda\subseteq\PP$ and pairs $(\alpha_p,\delta_p)\in\O_p^{\br}\times(\Z_{\geq0}\cup\{\infty\})$ such that 
\begin{equation}\label{UFD int val}
R=\bigcap_{p\in\Lambda}\Z_{(p),\alpha_p,\delta_p}\cap\Q[X]=\bigcap_{p\in\Lambda}\IntQ(\bar B_p(\alpha_p,\delta_p),\O_p)
\end{equation}
where, for each $p\in\Lambda$, $\alpha_p$ is unramified over $\Q_p$ and is  transcendental over 
$\Q$ if $\delta_p=\infty$ and the set  
$C(\uE)=\prod\limits_{p\in\Lambda}\{\alpha_{p} \mid \delta_{p}>0\}$ is
polynomially factorizable.
\end{Cor}
\begin{proof}
It is well-known that a UFD is a Krull domain with trivial class group (see for example \cite[Proposition 6.1]{Fossum}). By Theorem \ref{main thm polynomial krull domains}, $R$ is as in \eqref{Krull domain as int ring} (and so, in particular, it is a Noetherian domain by \cite{Heinzer}). By the characterization of the class group of $R$ in the aforementioned Theorem and also Theorem \ref{intersection DVR is UFD}, $R$ is as in the statement of the corollary. Note that the second equality of \eqref{UFD int val} follows from Remark \ref{balls and intval rings}.
\end{proof}


\begin{Cor}\label{coro433}
Let $\Lambda\subseteq\PP$ be a subset and  $(\alpha_{p,j},\delta_{p,j}), (\alpha_{p,j},\delta_{p,j}') \in\O_p^{\br}\times(\Q_{\geq0}\cup\{\infty\})$ 
with $\delta_{p,j} \leq \delta_{p,j}'$, and $\alpha_{p,j}$ transcendental over $\Q$ if $\delta_{p,j}=\infty$, 
for each $p\in\Lambda$ and $j=1,\ldots,m_p$. Let 
%\begin{eqnarray*}
$$R = \bigcap_{p\in\Lambda}\bigcap_{j=1}^{m_p}\Z_{(p),\alpha_{p,j},\delta_{p,j}}\cap\Q[X] \text{ and }
R' = \bigcap_{p\in\Lambda}\bigcap_{j=1}^{m_p}\Z_{(p),\alpha_{p,j},\delta_{p,j}'}\cap\Q[X].$$
%\end{eqnarray*}
%Then $\Z[X] \subseteq R \subseteq R' \subseteq \Q[X]$. 
Let $\Delta = \{p \in \Lambda \mid 0 = \delta_{p, j} < \delta_{p,j}'$ for some $j=1, \dots, m_p\}$ and assume that $\Delta$ is a finite set. Then the following statements hold.
\begin{enumerate}
\item[\em (1)]  $R$ is a Krull domain if and only if $R'$ is a Krull domain. %In this case, if $\delta_{p,j}' = \infty$ for each $p\in\Lambda$ and %$j=1,\ldots,m_p$, then $R'$ is a Dedekind domain.   
\item[\em (2)] Let $R$ be a Krull domain. Then Cl$(R')$ is a factor group of Cl$(R)$. Moreover, if $\delta_{p,j}' < \infty$ for all $\delta_{p,j}$
with $\delta_{p,j} < \infty$, then Cl$(R') =$ Cl$(R)$.
\item[\em (3)] If $R$ is a UFD, then $R'$ is also a UFD.
\end{enumerate}
\end{Cor}

\begin{proof}
(1) For each $g\in\Z[X]$, let
\begin{eqnarray*}
\mathbb D_{g,\uE}(R) &=& \{p\in\Lambda \mid\exists j\in\{1,\ldots,m_p\} \text{ such that }g\in\Z[X]\cap\ M_{(p),\alpha_{p,j},\delta_{p,j}}\}\\
\mathbb D_{g,\uE}(R') &=& \{p\in\Lambda \mid\exists j\in\{1,\ldots,m_p\} \text{ such that }g\in\Z[X]\cap\ M_{(p),\alpha_{p,j},\delta_{p,j}'}\}.
\end{eqnarray*}
Then, by Theorem \ref{Krull domains finite character}, it suffices to show that $\mathbb D_{g,\uE}(R)$ is finite if and only if $\mathbb D_{g,\uE}(R')$ is finite. 

It is clear that $\Z[X]\cap\ M_{(p),\alpha_{p,j},\delta_{p,j}} \subseteq \Z[X]\cap\ M_{(p),\alpha_{p,j},\delta_{p,j}'}$ for each $p\in\Lambda$ and $j=1,\ldots,m_p$. Hence, $\mathbb D_{g,\uE}(R) \subseteq \mathbb D_{g,\uE}(R')$. Moreover, if $\delta_{p,j} \neq 0$ for $j =1, \dots, m_p$, then 
$\Z[X]\cap\ M_{(p),\alpha_{p,j},\delta_{p,j}} = \Z[X]\cap\ M_{(p),\alpha_{p,j},\delta_{p,j}'}$ by Proposition \ref{contraction maximal ideal}. Hence, $\mathbb D_{g,\uE}(R') \setminus \mathbb D_{g,\uE}(R) \subseteq \Delta$, and since $\Delta$ is finite, $\mathbb D_{g,\uE}(R') \setminus \mathbb D_{g,\uE}(R)$ is finite. Thus $\mathbb D_{g,\uE}(R)$ is finite if and only if $\mathbb D_{g,\uE}(R')$ is also finite. 

(2) Let the notation be as in Theorems \ref{characterization local Krull domain} and \ref{main thm polynomial krull domains}. 
 We may assume that the representation of $R$ above is irredundant. Hence, by Proposition \ref{overring intersection DVRs}, the representation of $R'$ above is also irredundant. 
Note that $(\alpha_{p,j}, \delta_{p,j})$ is a minimal pair if $\delta_{p,j} \in \Q$, so if $\delta_{p,j}' < \infty$, then 
$(\alpha_{p,j}, \delta_{p,j}')$ is a minimal pair because 
$\bar B_p(\alpha_{p,j},\delta_{p,j}')\subseteq \bar B_p(\alpha_{p,j},\delta_{p,j})$. Hence 
Cl$(R') = \bigoplus\limits_{p \in \Lambda}(\Z/d'_p\Z \oplus \Z^{m_p-1})$ for some positive integer $d'_p$ with $d'_p|d_p$ for each $p \in \Lambda$ by Theorem \ref{characterization local Krull domain}. Hence, $\Z/d_p'\Z \cong (\Z/d_p\Z)/(d_p'\Z/d_p\Z)$, so if we let $H = \bigoplus\limits_{p \in \Lambda}(d_p'\Z/d_p\Z \oplus \{0\}^{m_p-1})$, then $H$ is a subgroup of Cl$(D)$ and Cl$(R)/H =$ Cl$(R')$. In particular, if $\delta_{p,j}' < \infty$ for all $\delta_{p,j}$
with $\delta_{p,j} < \infty$, then $d_p = d_p'$ by Theorem \ref{characterization local Krull domain} again. Thus Cl$(R)=$ Cl$(R')$.  

(3) This follows by (2) above because a UFD is a Krull domain with trivial class group.
\end{proof}

\begin{Ex}
Let $R = \bigcap\limits_{p \in \PP}\Z_{(p), 0, 0} \cap \Q[X]$. Then $R = \Z[X]$, and hence $R$ is a Krull domain (in fact, a UFD). However, if $R'= \bigcap\limits_{p \in \PP}\Z_{(p), 0, p} \cap \Q[X]$, then $R$ is not a Krull domain because $X \in X\Z_{(p), 0, p} = p\Z_{(p), 0, p}$ for all $p \in \PP$. Therefore the condition of $\Delta$ being finite is necessary for Corollary \ref{coro433}.  
\end{Ex}


\subsubsection{Construction of Krull domains with prescribed class group}\label{section: construction Krull}


For convenience, we say that a polynomial Krull domain $R$ is a \emph{pure} Krull domain if $Q$ is not maximal for all $Q \in X^1(R)$ with $Q \cap \Z \neq (0)$. By Proposition \ref{Krull intersection}, if $R$ is a polynomial Krull domain that is not Dedekind, then $R$ can be represented as an intersection of a Dedekind domain $R_1$ and a pure Krull domain $R_2$, i.e., $R = R_1 \cap R_2$. 
The divisor class group of a polynomial Krull domain is a direct sum of a countable family of finitely generated abelian groups by Theorem \ref{main thm polynomial krull domains}. Conversely, if $G$ is a direct sum of a countable family of finitely generated abelian groups, then there is a Dedekind domain between $\Z[X]$ and $\Q[X]$ with ideal class group $G$ \cite[Theorem 3.1]{PerDedekind}. 

We next construct a pure Krull domain with prescribed divisor class group.
 
\begin{Prop}\label{existence Krull local case}
Let $G$ be a finitely generated abelian group. Then there exists a pure Krull domain $R$ between $\Z[X]$ and $\Q[X]$ such that Cl$(R)=G$. 
\end{Prop}

\begin{proof}
Suppose $G=\Z/e_1\Z \oplus \cdots \oplus \Z/e_n\Z \oplus \Z^m$
for some integers $e_1, \dots , e_n \in \N$ and $m \geq 0$. 
Let $p_1, \dots , p_n \in \PP$ be distinct. 

We pick $m+1$ distinct elements $\alpha_{p_n, 1},\ldots,\alpha_{p_n, m+1}$ in $\Z_{p_n}$ which are transcendental over $\Q$.  Let $N=\max\{v_{p_n}(\alpha_{p_n,i}-\alpha_{p_n, j})\mid i\not=j\in\{1,\ldots,m+1\}\}$. 
Let $W_{p_n, j} = \Z_{(p_n),\alpha_{p_n, j}, N+1/e_n}$, for $j\in\{1,\ldots,m +1\}$.  Note that the $W_{p_n,j}$'s are distinct (because $v_{p_n}(\alpha_{p_n, i}-\alpha_{p_n,j})<N+1/e_n$ for each 
$i\not=j$) and $(\alpha_{p_n,j}, N+ 1/e_n)$ is a minimal pair for each 
$j\in\{1,\ldots,m+1\}$. Then by Theorem \ref{characterization local Krull domain}, (4) and Proposition \ref{unramified implies k=1} we have 
 $$e_{p_n,j}=e(W_{p_n,j}\mid\Z_{(p_n)})=e(N+1/e_n\mid\Q_p)=e_n.$$
By the same argument, we can also choose $\alpha_{p_i}$ in $\Z_{p_i}$
which is transcendental over $\Q$, so that $(\alpha_{p_i}, 1/e_i)$ is minimal pair for $i=1, \dots, n-1$. Let $W_{p_i} = \Z_{(p_i), \alpha_{p_i}, 1/e_i}$ for $i =1, \dots, n-1$  so that, if we set $R= (\bigcap\limits_{i=1}^{n-1}W_{p_i}) \cap (\bigcap\limits_{j=1}^{m+1}W_{p_n,j}\cap\Q[X])$, 
we have that $R$ is a Krull domain such that Cl$(R)=G$
by Theorem \ref{main thm polynomial krull domains}
and $Q$ is not maximal for all $Q \in X^1(R)$ with $Q \cap \Z \neq (0)$
by Theorem \ref{characterization local Krull domain}(3).
\end{proof}

 
\begin{Thm}
Let $G$ be a direct sum of a countable family $\{G_i \mid i \in I\}$ of finitely generated abelian groups (which are not necessarily distinct).
Then there exists a pure Krull domain $R$ between $\Z[X]$ and $\Q[X]$ with Cl$(R)=G$. Moreover, for each $i \in I$, there is a multiplicative subset $S_i$ of $\Z$ such that $R_{S_i}$ is a pure Krull domain with Cl$(R_{S_i}) = G_i$.  
\end{Thm}

\begin{proof}
For each $i\in I$, let $G_i=\Z/e_{i1}\Z \oplus \cdots \oplus \Z/e_{in_i}\Z \oplus \Z^{m_i}$ for some $e_{i1}, \dots, e_{in_i} \geq 1$, $n_i \geq 1$ and $m_i\geq 0$. Without loss of generality, we may suppose that we have a partition of $\PP$ into a family of finite subsets which is indexed by $I$, namely,  $\{\PP_i\}_{i \in I}$ is a partition of $\PP$ with $|\PP_i| = n_i$ for each $i\in I$. By Proposition \ref{existence Krull local case}, there exists a pure Krull domain $R_i$, $\Z[X]\subseteq R_i\subset \Q[X]$, such that Cl$(R_i)=G_i$. 
%{\rc Note that we have
%$$R_i=\IntQ(E_i,\overline{\Z_{p_i}}),\; \text{where } E_i=\bigcup_{j=1}^{m_i+1}\bar B_{p_i}(\alpha_{p_i,j},\delta_{i})$$
%for some $\alpha_{p_i,j}\in\Q_{p_i}$ and $\delta_i=N_i+1/e_i$ if $G_i$ is not trivial (either $m_i\geq1$ or $e_i\geq2$) or $\delta_{i}=0$ otherwise.} In the last case, $E_i=\overline{\Z_{p_i}}$ and $R_i=\Z_{(p_i)}[X]$. 
%Note that $\delta_i>0\Leftrightarrow G_i$ is not trivial. 

Now set $R=\bigcap\limits_{i\in I}R_i$. Note that if $S_i = \Z \setminus \bigcup\limits_{p \in \PP_i}p\Z$ for each $i \in I$, then $R_{S_i} = R_i$ by Lemma \ref{lemma43}, thus
$R_{S_i}$ is a pure Krull domain with Cl$(R_{S_i}) = G_i$.
However, it may be the case that $R$ is not a Krull domain, that is, it has not the finite character property.

Let $i\in I$ be such that $G_i$ is not trivial, and let $p$ be a fixed prime in the finite set $\PP_i$. Suppose that for the construction of $R_i$ we chose in $\Z_p$ $k_p\geq 1$ distinct elements, say for example $E_p=\{\alpha_{p,1},\ldots,\alpha_{p,k_p}\}$; let also $\delta_{p,j}$ be the radius of the ultrametric ball of center $\alpha_{p,j}$, for $j=1,\ldots,k_p$. We then   replace the elements in $E_p$ by possibly other  $k_p$ distinct $\alpha_p$'s in $\Z_p$ which lie in the same residue class $\lfloor\log(p)\rfloor+p\Z_{p}$, such that these $\alpha_p$'s are transcendental over $\Q$. We also modify the corresponding $\delta_{p,j}$'s accordingly so that the previous facts about the corresponding ring $R_i$ still hold. If $G_i$ is trivial, then we simply choose an element $\alpha_p\in\Z_{p}$. We claim that the set 
$$C(\underline E)=\prod_{p\in \PP}\{\alpha_{p,1},\ldots,\alpha_{p,k_p}\in \Z_p\mid \delta_{p,j}>0\}$$
is polynomially factorizable, so by Theorems \ref{Krull domains finite character} and \ref{characterization local Krull domain}(3), $R$ is a pure Krull domain. Let $g\in\Z[X]$ and let $\underline{\alpha}=(\alpha_p)\in C(\underline E)$. Then we have
$$g(\alpha_p)\equiv g(\lfloor\log(p)\rfloor)\pmod{p}$$
Since for $p$ sufficiently large, $g(\lfloor\log(p)\rfloor)$ is not divisible by $p$ we get that $\PP_{g,C(\underline{E})}$ is finite.
\end{proof}

 \vskip1cm

\subsubsection{Almost Dedekind domains  and its generalization} 

Let $R=\bigcap\limits_{p\in\Lambda}\bigcap\limits_{j=1}^{m_p}\Z_{(p),\alpha_{p,j},\delta_{p,j}}\cap\Q[X]$ be as in Theorem \ref{Krull domains finite character}. We know that $R$ need not be a Krull domain (see Example \ref{ex430}). We now investigate some ring-theoretic properties of $R$ even when it is not a Krull domain.


\begin{Cor} \label{thm3}
Let $\Lambda$ be a nonempty subset of $\PP$.
For each $p \in \Lambda$, suppose that $\Z_{(p),\alpha_{p,j},\delta_{p,j}}$,
for some $(\alpha_{p,j},\delta_{p,j})\in\O_p^{\br}\times(\Q_{\geq0}\cup\{\infty\})$ is a family of DVRs of $\Q(X)$ extending $\Z_{(p)}$
and $M_{p,j} = M_{(p),\alpha_{p,j},\delta_{p,j}}$ for $j = 1, \dots , m_p$. Now, let
\begin{equation}\label{representation R3}
R = \bigcap_{p \in \Lambda} (\bigcap_{j=1}^{m_p}\Z_{(p),\alpha_{p,j},\delta_{p,j}} \cap\Q[X]).
\end{equation}
Then $R_P$ is a DVR for all $P \in X^1(R)$, $R = \bigcap\limits_{P \in X^1(R)}R_P$ and $I_t = R$ for all ideals $I$ of $R$ with $I \nsubseteq P$ for all $P \in X^1(R)$. 
Moreover, if the representation in \eqref{representation R3} is irredundant, then we have the following.
\begin{enumerate}
%\item[{\em (1)}] $R$ is a $t$-almost Dedekind domain. 
\item[{\em (1)}] $X^1(R) = (\bigcup\limits_{p \in \Lambda}\{M_{p,j} \cap R \mid j=1, \dots , m_p\}) \cup \{f\Q[X] \cap R \mid f \in \Q[X]$ is irreducible over $\Q\}$. 

%\item[{\em (3)}] If $e_{p,j} = %e(\Z_{(p),\alpha_{p,j},\delta_{p,j}}|\Z_{(p)})$ for $j=1, \dots , m_p$ %and $d_p= gcd(e_{p,1}, \dots , e_{p,m_p})$
%for all $p \in \Lambda$, then Cl$(R) = \bigoplus\limits_{p \in \Lambda}%(\Z/{d_p}\Z \oplus \Z^{m_p-1})$.

\item[{\em (2)}] $R = \IntQ(\uE,\O)$, for some subset $\uE=\prod_{p\in\Lambda}E_p\subset\O=\prod\limits_{p\in\Lambda}\O_p$ 
such that $E_p = \bigcup\limits_{j=1}^{m_p} \bar B_p(\alpha_{p,j},\delta_{p,j})$ $\subset\O_p^{\br}$, where $\alpha_{p,j}$ is transcendental over $\Q$ if $\delta_{p,j}=\infty$.
\end{enumerate}
\end{Cor}

\begin{proof}
The first part follows directly from \cite[Proposition 2.2(3)]{Chang}
and (1) is from \cite[Corollary 2.3]{Chang}.
For (2), note that, by Theorem \ref{thm4.23}, $R = \bigcap\limits_{p \in \Lambda}\IntQ(E_p,\O_p)$,
 where $E_p$ is a finite union of balls $\bar B_p(\alpha_{p,j},\delta_{p,j})\subset\O_p^{\br}$, for $j=1,\ldots,m_p$, with  $\alpha_{p,j}$ transcendental over $\Q$ if $\delta_{p,j} =\infty$. Thus, $R = \IntQ(\uE,\O)$, where $\uE=\prod\limits_{p\in\Lambda} E_p$.
\end{proof}



 We recall that an integral domain $D$ is an almost Dedekind domain if $D_M$ is a DVR for each maximal ideal $M$ of $D$. The next corollary characterizes when an intersection of DVRs with $\Q[X]$ is an almost Dedekind domain.
%{\rc I am wondering now if we can now describe ALL the almost Dedekind domains between $\Z[X]$ and $\Q[X]$; I guess that in that case we must allow the quantity $m_p$ to be infinite, isn't it?}
%{\bc I don't know wheter $R$ is almost Dedekind or not when $m_p = \infty$}

\begin{Cor} \label{coro1}
Let $\Lambda$ be a nonempty subset of $\PP$.
For each $p \in \Lambda$ and $j=1, \dots, m_p$, suppose that $\Z_{(p),\alpha_{p,j},\delta_{p,j}}$,
for some $(\alpha_{p,j},\delta_{p,j})\in \O_p^{\br}\times(\Q_{\geq0}\cup\{\infty\})$, is a family of DVRs of $\Q(X)$ extending $\Z_{(p)}$.  
%and $M_{j,p} = M_{(p),\alpha_{j,p},\delta_{j,p}}$ for $j = 1, \dots , m_p$.
Let
\begin{equation}\label{representation R4}
R = \bigcap_{p \in \Lambda} (\bigcap_{j=1}^{m_p}\Z_{(p),\alpha_{p,j},\delta_{p,j}} \cap\Q[X]).
\end{equation}
%If ht$(M_{p,j} \cap R)=1$ for $p \in  \Lambda$ and $j=1, \dots , m_p$,
Then the following statements are equivalent:
\begin{enumerate}
\item[{\em (1)}] $R$ is an almost Dedekind domain.
\item[{\em (2)}] $\Z_{(p),\alpha_{p,j},\delta_{p,j}}$
is residually algebraic over $\Z_{(p)}$ for all $p \in \Lambda$ and $j =1, \dots, m_p$.
\item[{\em (3)}] Each non-unitary prime ideal of $R$ is a maximal ideal. 
\end{enumerate}
\end{Cor}
%\color{red}
%\textbf{Question}: if we want to characterize Almost Dedekind domains between $\Z[X]$ and $\Q[X]$ then why for each $p\in\PP$ we restrict to only finitely %many of them, that is, $m_p<\infty$? 
%I mean, in general we may have an almost Dedekind domain for which $m_p$ is also infinite (provided, I think, the topological closure of $\{\alpha_{p,j}\}$ %does not contain any algebraic element).
%{\color{blue} You are right. This is because I could prove only the finite case,
%but I don't know how to solve the infinite case.}
%\normalcolor
\begin{proof}
(1) $\Rightarrow$ (2) For a fixed $p \in \Lambda$, let $S = \Z \setminus p\Z$.
Then $R_S = \bigcap\limits_{j=1}^{m_p}\Z_{(p),\alpha_{p,j},\delta_{p,j}} \cap\Q[X]$ by Lemma \ref{lemma43}
and $R_S$ is an almost Dedekind domain which is also a Krull domain, so $R_S$ is a Dedekind domain, and hence the intersection
$R_S = \bigcap\limits_{j=1}^{m_p}\Z_{(p),\alpha_{p,j},\delta_{p,j}} \cap\Q[X]$ is irredundant.
Hence, by Theorem \ref{characterization local Krull domain}(3),  $\Z_{(p),\alpha_{p,j},\delta_{p,j}}$
is residually algebraic over $\Z_{(p)}$ for all $p \in \Lambda$ and $j =1, \dots, m_p$.

(2) $\Rightarrow$ (3) By \cite[Corollary 2.6(1)]{Chang},
$R$ is an almost Dedekind domain. Hence, each nonzero prime ideal of $R$ is maximal. Thus, each non-unitary prime ideal of $R$ is also maximal.

(3) $\Rightarrow$ (1) Assume to the contrary that $R$ is not an almost Dedekind domain. Then $R_{\mathbb{Z} \setminus p\mathbb{Z}}$ is not an almost Dedekind domain for some $p \in \Lambda$. Let $S = \Z \setminus p\Z$. Then $R_{S} = \bigcap\limits_{j=1}^{m_p}\Z_{(p),\alpha_{p,j},\delta_{p,j}} \cap\Q[X]$ by Lemma \ref{lemma43}, so $R_{S}$ is a Krull domain but not a Dedekind domain. Hence, dim$R_S =2$,
and thus there is a height two prime ideal $M$ of $R$ such that $MR_S \cap \Q = p\Z_{(p)}$. Note that ht$(MR_S) =$ ht$M =2$, so $MR_S$ is a union of
infinitely many height one prime ideals of $R_S$ by the prime avoidance theorem \cite[Proposition 4.9]{Gilm} or \cite[Proposition 2.1]{Heitmann}.
Note that the height one prime ideals of $R_S$ containing $p$ is finite, so $MR_S$ contains infinitely
many prime ideals $Q$ of $R$ with $p \not\in Q$; in particular, these prime ideals $Q$ are not maximal. 
Clearly such a prime ideal $Q$ must be $Q \cap \Z = (0)$. 
\end{proof}

Let $V$ be a rank-one nondiscrete valuation domain with value group $G \subsetneq \mathbb{R}$. Then $V$ is completely integrally closed,
Pic$(V) = \{0\}$ and Cl$(V) = \mathbb{R}/G$ \cite[Theorem 2.7]{afz08}.
Hence, in general, Pic$(D) \neq$ Cl$(D)$ for a completely integrally closed domain $D$. An almost Dedekind domain $D$ is completely integrally closed, so Cl$(D)$ is well defined and Pic$(D)$ is a subgroup of Cl$(D)$. In the next corollary we calculate the Picard of an almost Dedekind domain but we don't known how to calculate its divisor class group.   

\begin{Cor} \label{coro1-1}
Let the notation be as in Corollary \ref{coro1}.
If $R$ is an almost Dedekind domain, then
the following statements are satisfied. 
\begin{enumerate}
\item[{\em (1)}] $R = \IntQ(\uE,\O)$, for some subset $\uE=\prod\limits_{p\in\Lambda}E_p\subset\O=\prod\limits_{p\in\Lambda}\O_p$ 
such that $E_p = \{\alpha_{p,1}, \dots , \alpha_{p,m_p}\}$ $\subseteq \O_p^{\br}$, $\alpha_{p,j}$ is transcendental over $\Q$ and $\delta_{p,j}=\infty$ for $j=1,\ldots,m_p$.
\item[{\em (2)}] Pic$(R)= \bigoplus\limits_{p \in \Lambda}(\Z/d_p\Z \oplus \Z^{m_p-1})$, where
$e_{p,j} = e(\Z_{(p),\alpha_{p,j},\delta_{p,j}}|\Z_{(p)})$ for $j=1, \dots , m_p$ and $d_p= gcd(e_{p,1}, \dots , e_{p,m_p})$ for all $p \in \Lambda$.
\end{enumerate} 
\end{Cor}

\begin{proof}
(1)  By Corollary \ref{thm3}, $R = \IntQ(\uE,\O)$, where $\uE=\prod\limits_{p\in\Lambda} E_p$ and $E_p = \bigcup\limits_{j=1}^{m_p} \bar B_p(\alpha_{p,j},\delta_{p,j})$ $\subset\O_p^{\br}$.
Moreover, by Corollary \ref{coro1}, $\Z_{(p),\alpha_{p,j},\delta_{p,j}}$
is residually algebraic over $\Z_{(p)}$ for all $p \in \Lambda$ and $j =1, \dots, m_p$. Hence, $\delta_{p,j}=\infty$, $\bar B_p(\alpha_{p, j},\delta_{p,j}) = \{\alpha_{p, j}\}$, $\alpha_{p,j}$ is transcendental over $\Q$ for $j=1, \dots , m_p$ and $E_p = \{\alpha_{p,1}, \dots , \alpha_{p,m_p}\}$.

(2) This follows directly from the proof of \cite[Theorem 2.4(4)]{Chang}.
\end{proof} 


In general, the almost Dedekind domain of (\ref{representation R4}) need not be a Dedekind domain 
(see \cite[Theorems 3.1 and 3.5]{Chang} for a concrete example).



\subsection*{Acknowledgments}
This research was completed while the first author visited Universita of Padova during January 2026. He thanks the Department of Mathematics for its hospitality. The first author was supported by Basic Science Research Program through the National Research Foundation of Korea (NRF) funded by the Ministry of Education (2017R1D1A1B06029867). The second author  is a member of the National Group for Algebraic and Geometric Structures and their Applications (GNSAGA) of the Italian Mathematics Research Institute (INdAM). 
 
%\newpage
\begin{thebibliography}{99}
\bibliographystyle{plain}

\bibitem{AlePop} V. Alexandru, N. Popescu, \emph{Sur une classe de prolongements \`a K(X) d'une valuation sur un corps $K$}, Rev. Roumaine Math. Pures Appl. 33 (1988), no. 5, 393-400.

\bibitem{APZTheorem} V. Alexandru, N. Popescu, A. Zaharescu, \emph{A theorem of characterization of residual transcendental extensions of a valuation}. J. Math. Kyoto Univ. 28 (1988), no. 4, 579-592.

\bibitem{APZMinimal} V. Alexandru, N. Popescu, A. Zaharescu, \emph{Minimal pairs of definition of a residual transcendental extension of a valuation}. J. Math. Kyoto Univ. 30 (1990), no. 2, 207-225.

\bibitem{APZAll} V. Alexandru, N. Popescu, A. Zaharescu, \emph{All valuations on $K(X)$}. J. Math. Kyoto Univ., 30 (1990), no. 2, 281–296.

\bibitem{APZclosed} V. Alexandru, N. Popescu, A. Zaharescu, \emph{On the closed subfields of $\mathbb{C}_p$}. J. Number Theory 68 (1998), no. 2, 131-150.

\bibitem{afz08} D.F. Anderson, M. Fontana, and M. Zafrullah,
{\em Some remarks on Pr\"ufer $*$-multiplication domains and class groups}, J. Algebra 319 (2008), 272-295.


\bibitem{Bourb} N. Bourbaki, {\em Alg\`ebre commutative}, Hermann, Paris, 1961.

\bibitem{cc97} P.-J. Cahen and J.-L. Chabert, {\em Integer-Valued Polynomials},
 Amer. Math. Soc. Surveys and Monographs, Vol. 48, Providence, 1997.
 
%\bibitem{CaChNew} P.-J. Cahen, J.-L. Chabert, \emph{What's new about integer-valued polynomials on a subset?}. 
%Non-Noetherian commutative ring theory, 75–96. Math. Appl., 520 Kluwer Academic Publishers, Dordrecht, 2000.

\bibitem{ChabPolCloVal} J.-L. Chabert, \emph{On the polynomial closure in a valued field}, J. Number Theory 130 (2010), 458-468.

%\bibitem{ChabPers} J.-L. Chabert, personal communication. %compact iff no pmt with non-zero breadth ideal

%\bibitem{ChabPer} J.-L. Chabert, G. Peruginelli, \emph{Polynomial overrings of $\Int(\Z)$}, J. Commut. Algebra 8 (2016), no. 1, 1-28.


\bibitem{Chang} G. W. Chang, \emph{The ideal class group of polynomial overrings of the ring of integers}, J. Korean Math. Soc. 59 (2022), No. 3, pp. 571-594.

\bibitem{Chevalley} C. Chevalley, \emph{Introduction to the Theory of Algebraic Functions of One Variable.} Mathematical Surveys, No. VI. American Mathematical Society, New York, N. Y., 1951.

%\bibitem{Claborn} L. Claborn, \emph{Every abelian group is a class group}, Pacific J. Math. 18 (1966), 219-222.

%\bibitem{EakHei} P. Eakin, W. Heinzer, \emph{More noneuclidian PID's and Dedekind domains with prescribed class group}. Proc. Amer. Math. Soc. 40 (1973), 66-68.

\bibitem{EngPre} A. J. Engler, A. Prestel, \emph{Valued fields}. Springer-Verlag, Berlin, 2005.


\bibitem{Fossum} R. M.  Fossum, \emph{The divisor class group of a Krull domain}. Ergebnisse der Mathematik und ihrer Grenzgebiete, Band 74. Springer-Verlag, New York-Heidelberg, 1973.


\bibitem{Gilm} R. Gilmer, \emph{Multiplicative ideal theory}, Corrected reprint of the 1972 edition, Queen's University, Kingston, ON, 1992.

\bibitem{gh68} R. Gilmer and W. Heinzer, \emph{Irredundant intersections of valuation rings}, Math. Z., 103 (1968), 306-317.

\bibitem{ghls90} R. Gilmer, W. Heinzer, D. Lantz, and W. Smith,
{\em The ring of integer-valued polynomials of a Dedekind domain}, Proc. Amer. Math. Soc.
108 (1990), 673-681.

\bibitem{g97} F.Q. Gouv\^ea, {\em P-adic Numbers: An introduction}, 2nd Edition, New York, Springer, 1997.

\bibitem{Heinzer} W. Heinzer, \emph{On Krull overrings of a Noetherian domain}. Proc. Amer. Math. Soc. 22 (1969), 217-222.

\bibitem{HeinzerOhm} W. Heinzer, J. Ohm, \emph{Noetherian intersections of integral domains}. Trans. Amer. Math. Soc. 167 (1972), 291-308.

\bibitem{Heitmann} R. C. Heitmann, \emph{PID's with specified residue fields}. Duke Math. J. 41 (1974), 565-582.

\bibitem{k89} B.G. Kang,  {\em On the converse of a well-known fact about Krull domains}, J. Algebra 124 (1989), 284-299.


\bibitem{Khanduja} S. K. Khanduja and J. Saha, \emph{A generalized fundamental principle}, Mathematika, 46:1 (1999), 83-92.

%\bibitem{k89-1} B.G. Kang, {\em Pr\"{u}fer $v$-multiplication domains and the ring $R[X]_{N_v}$},J. Algebra 123 (1989), 151-170.

%\bibitem{McQ Gilmer} D. L. McQuillan, \emph{On a theorem of R Gilmer}. J. Number Theory, 39 (1991) 3, 245-250.


\bibitem{n55} M. Nagata, {\em On the derived normal rings of Noetherian integral domains},
Mem. Coll. Sci. Univ. Kyoto Ser. A. Math. 29 (1955), 293-303.

%\bibitem{Nark} W. Narkiewicz, \emph{Elementary and analytic theory of algebraic numbers}. Third edition. Springer Monographs in Mathematics. Springer-Verlag, Berlin, 2004.

\bibitem{NN} E. Nart, J. Novacoski, \emph{Geometric parametrization of valuations on a polynomial ring in one variable}. Math. Z. 304 (2023), no.4, 59.

%\bibitem{OhmRuled} J. Ohm, \emph{The ruled residue theorem for simple transcendental extensions of valued fields}. Proc. Amer. Math. Soc. 89 (1983), no. 1, 16-18.


\bibitem{PerTransc} G. Peruginelli, \emph{Transcendental extensions of a valuation domain of rank one}, Proc. Amer. Math. Soc. 145 (2017), no. 10, 4211-4226.

\bibitem{PerPrufer} G. Peruginelli, \emph{Pr\"ufer intersection of valuation domains of a field of rational functions}, J. Algebra 509 (2018), 240-262.


\bibitem{PerDedekind} G. Peruginelli, \emph{Polynomial Dedekind domain with finite residue fields of prime characteristic}, Pacific J. Math. 324 (2023), no. 2, 333–351.

\bibitem{PerStacked} G. Peruginelli, \emph{Stacked Pseudo-Convergent Sequences and Polynomial Dedekind Domains}, 
Algebra \& Number Theory 19 (2025), no. 10, 1947-1972.
%DOI: 10.2140/ant.2025.19.1947.



%\bibitem{PS1} G. Peruginelli, D. Spirito, \emph{The Zariski-Riemann space of valuation domains associated to pseudo-convergent sequences}, Trans. Amer. Math. Soc. 373 (2020), no. 11, 7959–7990.

\bibitem{PS2} G. Peruginelli, D. Spirito, \emph{Extending valuations to the field of rational functions using pseudo-monotone sequences}, J. Algebra 586 (2021), 756-786.

\bibitem{PSpolclo} G. Peruginelli, D. Spirito,  \emph{The polynomial closure is not topological}.
J. Pure Appl. Algebra 226 (2022), no. 12, No. 107133.

%\bibitem{PerWerNontrivial} G. Peruginelli, N. J. Werner. \emph{Non-triviality Conditions for Integer-valued Polynomials on Algebras},  Monatsh. Math. 183 (2017), no. 1, 177-189.

%\bibitem{PZstructure} N. Popescu, A. Zaharescu, \emph{On the structure of the irreducible polynomials over local fields}. J. Number Theory 52 (1995), no. 1, 98-118.

%\bibitem{RibVal} P. Ribenboim, \emph{Th\'eorie des valuations}. S\'eminaire de Math\'ematiques Sup\'erieures, No. 9 (\'Et\'e, 1964) Les Presses de l'Universit\'e de Montr\'eal, Montreal, Que. 1968.

\bibitem{z85}  M. Zafrullah, {\em A general theory of almost factoriality}, Manuscripta Math. 51 (1985), 29-62.

\bibitem{ZS2} O. Zariski, P. Samuel, \emph{Commutative Algebra, vol. II}, Springer-Verlag, New York-Heidelberg-Berlin, 1975.

\end{thebibliography}


\end{document} 
